\documentclass[12pt, twoside]{memoir}

% Version Description
\newcommand{\bookversion}{v260901
%	\begin{flushleft}
%		\begin{itemize}
%		\item This version is set up for 7x10, 14pt main text font. documentclass was changed to "memoir" to avoid each chapter starting on right-facing page (this is for e-book layout). 
%		\item Color formats are defined in the byrne.mp file. These are defined as RGB colors -- tried to change to cmykcolor but the diagrams failed to load. Instead, added a line in the .tex file to set to cmyk color output.
%		\item Def, Post, CN, and each Proposition is set as a Chapter with font and line spacing set in the preamble.
%		\item The usepackage lines were all moved to the preamble.
%		\item New commands were added to the preamble to simplify the formatting. the /justification[1] command will move a justification to the right margin of a line or put it on a new line and move it all the way to the right, both without interfering with centering. The /centeredjust[2] command has two inputs: centered main text(1) and justification text (2). It centers the main text on the page and adds the justification on the same line if there is room or on the subsequent line if no room.
%		\item All prop diagrams were modified such that there are no dashed lines in the prop diagrams. The reason for this is to avoid confusion of some lines looking like they are different kinds of things than other lines. Also, dashed lines look like a collection of points...
%		\item Changed preamble to use "memoir" class inherent commands to adjust spacing before/after chapter titles. Adjusting other ways wasn't working.
%		\end{itemize}
%	\end{flushleft}
%%% v260901 CORRECTED TYPO IN CN 3 ("wholes" changed to "remainders") %%%
} % change the book version here with each printing

% Input Preamble
\usepackage{byrne}
\usepackage{etoolbox}
\usepackage{ccicons,bxtexlogo,listings}
\usepackage{tikz}
\usepackage{comment}
\usepackage{fontspec}
\usepackage{lineno}
\usepackage{ebgaramond}
%\usepackage{microtype} % Improves line breaking and spacing
\usepackage{titling}
\usepackage{fancyhdr}
\usepackage{anyfontsize}
\usepackage{titlesec}
\usepackage[cmyk]{xcolor}
\selectcolormodel{cmyk}
%\usepackage{hyperref}


% 1. Format for Unnumbered Chapters (\chapter*)
\titleformat{name=\chapter,numberless}[block]
  {\normalfont\bfseries\Large}{}{0pt}{} % Large font size (17.28pt in 12pt document)

% 2. Format for Sections
\titleformat*{\section}{\normalfont\bfseries\large}

% 1. Remove the "Chapter Sink" (vertical drop)
\setlength{\beforechapskip}{-3\baselineskip} 
% Use a negative value (e.g., -20pt or -1cm) to pull it even higher if needed.

% 2. Adjust space between "Proposition X" and Enunciation
\setlength{\midchapskip}{10pt} 

% 3. Adjust space between "Proposition X" and the text body
\setlength{\afterchapskip}{20pt}

%%%% set format for 7x10 
\usepackage[
paperwidth=7in, 
paperheight=10in, 
inner=1.25in,
outer=0.75in, 
top=1in, 
bottom=1in
\setlength{\headsep}{15pt}
]{geometry} 

%%%% sets format for 8.5x11
%\usepackage[
%paperwidth=8.5in, 
%paperheight=11in, 
%inner=1in,
%outer=1in,
%top=1in,
%bottom=1in,
%headsep=15pt
%]{geometry}


\pagestyle{fancy}
\setlength{\parskip}{0pt}
\setlength{\headheight}{15pt}
\setlength{\parindent}{0pt}  % Set auto-indent to 0 point

% Define header/footer styles
\fancyhf{}
\fancyhead[LE,RO]{\leftmark}  % chapters (Def, CN, post, prop# in header
\fancyfoot[C]{\thepage} % page numbers centered at bottom

% Format footnotes (removes indent at start of footnote)
\makeatletter
\renewcommand\@makefntext[1]{%
  \noindent\makebox[1em][l]{\@thefnmark}#1}
\makeatother

% Number every nth line
\modulolinenumbers[1]

% Main Text Toggles
\label{Toggles}

% Set EB Garamond as the main font
\setmainfont{EB Garamond}

% Define line numbering toggle
\newtoggle{lineNumbering}

% Define text version toggle Byrnified HEATH original or RAUM-modified Heath
\newtoggle{heath} % Byrnified Heath text
\newtoggle{raum} % Raumified Heath
\newtoggle{heath1} % Byrnified Heath text without paragraph formatting
% Set values of toggles (uncomment one to enable it)
\toggletrue{heath}
\togglefalse{raum}
\togglefalse{heath1}

% Define and set toggle for pictures to be included in definitions or not
\newtoggle{defpics}
\toggletrue{defpics}

% Define text version toggle
\newtoggle{inline}

% Set line numbering on or off
\toggletrue{lineNumbering}

% Set inline drawings toggle (comment it out to turn off inline)
\toggletrue{inline}

\label{MetaPost Definitions}

\def\mpPre{
	u:=1cm;
    lineWidth := 2pt; % default is 2pt - increase for thicker lines in drawings
	textLabels := true;
	textLabelScaleFactor := 1; % set label size for diagrams
	textLabelAvSize := abs(ulcorner(textLabelRender("A")) - llcorner(textLabelRender("A")));
	angleSize := 1cm;
	compensateLineLabels := false; % this one removes vertical compensation for text labels for line segments
	lineLabelsOnTop := false; % this one moves line segment labels to the bottom
        angleModeFallback := true;
	}
 
\def\mpPost{
	textLabels := false;
    lineWidth := 2pt; % default is 2pt - increase for thicker lines in drawings
    textLabelScaleFactor := 1; % set label size for inline lines
	}

\label{Commands, Parameters}

\label{Justification Command}
% Robust command for formatting justifications (e.g. "C.N. 1")
 
\newcommand{\justification}[1]{%
  \unskip                       % 1. Remove space before command
  \hfill                        % 2. Push to the right
  \penalty100                   % 3. "If you must break the line, do it here"
  \hskip 1em plus 1fill         % 4. The "Smart Spring"
  \mbox{[\textit{#1}]}%         % 5. The justification (Italicized, normal size)
  \par                          % 6. Finish the line
}

\newcommand{\justificationNew}[1]{%
  \unskip                % Remove extra space
  \hspace*{2em}          % Minimum gap before the note
  \nobreak\hfill         % Try to stay on this line and push right
  \allowbreak            % If it's too tight, allow a break here
  \hspace*{\fill}        % If it breaks, push the note to the right on the new line
  \mbox{[\textit{#1}]}% The note itself
  \par                   % End the line/para
}

%%%%%
% Command for centering main text (#1), right-flushing the justification (#2)

\usepackage{calc}
\newlength{\textlength}
\newlength{\justlength}

\newcommand{\centeredjust}[2]{%
	%\unskip
    % 1. Measure the text and the justification
    \settowidth{\textlength}{#1}
    \settowidth{\justlength}{[\textit{#2}]}
    
    % 2. Logic: If (text + justification + buffer) > linewidth, drop to next line
    \ifdim\dimexpr\textlength+\justlength+1em\relax > \linewidth
        {\centering #1 \par}
        \hspace*{\fill}[\textit{#2}]\par
    \else
        % 3. If it fits, use the "Balanced Ghost" method to keep text centered
        \noindent\hspace*{\fill}\makebox[0pt][l]{\phantom{[\textit{#2}]}}#1%
        \hspace*{\fill}\llap{[\textit{#2}]}\par
    \fi
}
%%%%%

\def\pointlabel#1{\textsf{\textbf{#1}}} % This one determines the style of text labels both on the main diagram and in the text

\newdimen\textLabelShift % This thing allows you to move labeled line segments up and down to make the labels appear like normal characters

\textLabelShift = -1.6pt % This value works for 12pt EB Garamond, but might not work with a different font/size

\newdimen\circleSpace % Space between inline circles and the letters

\circleSpace = 2pt % This value works for 12pt EB Garamond, but might not work with a different font/size

\newdimen\pointShift % This allows you to move the inline points backward (negative) or forward (positive) to fine tune the spacing -- the \pointlabel function seems to usually insert a space before points. 

\pointShift = -3.5pt % This value works for 12pt EB Garamond, but might not work with a different font/size

\newdimen\lineShift % This allows you to move the inline lines backward (negative) or forward (positive) to fine tune the spacing -- the \pointlabel function seems to usually insert a space before points. 

\lineShift = -5pt % This value works for 12pt EB Garamond, but might not work with a different font/size

\newdimen\inlineLength % this is the length of the inline lines. 

\inlineLength = 0.55cm % Length of inline lines should be ~0.046 * fontsize (it's 0.55cm for 12pt)

\DeclareDocumentCommand{\drawUnitLineAlt}{o m O{0.55cm}}{% This is an alternative macro for inline line segments. Optional input is to rename the line
    \iftoggle{inline}{
        \hspace{\lineShift}
        \advance\pictOffsetBottom by \textLabelShift
	    \advance\pictOffsetTop by \textLabelShift
    	\IfNoValueTF{#1}%
    		{\drawFromCurrentPicture[top][uline#2]{%
    			textLabels := true;%
    			startAutoLabeling;%
    			draw byNamedCompoundLine(#3, 0)(#2);
    			stopAutoLabeling;%
    			textLabels := false;%
    			}}%
    		{\drawFromCurrentPicture[top][uline#2]{%
    			textLabels := true;%
    			startAutoLabeling;%
    			draw byNamedLineSeqByPoint(1)(#2);%
    			stopAutoLabeling;%
    			textLabels := false;%
    			}}%
	    \advance\pictOffsetBottom by -\textLabelShift%
	    \advance\pictOffsetTop by -\textLabelShift}%
        {%
            \IfNoValueTF{#1}%
                {\pointlabel{#2}%
                }%
                {\pointlabel{#1}}%
        }%
}%

\DeclareDocumentCommand{\drawAngleAlt}{o m}{% This alternative in-line angle command draws the named angle, then puts the name (or the optional argument) in bold. 
    \iftoggle{inline}{%
        \drawAngle{#2}%
	\IfNoValueTF{#1}%
	{\pointlabel{#2}}%
	{\pointlabel{#1}}%
	}%
    {%
        \pointlabel{#2}%
    }%
}%

 \DeclareDocumentCommand{\drawPolygonAlt}{o m}{% This alternative in-line polygon command draws the named polygon, then puts the name (or the optional argument) in bold. 
    \iftoggle{inline}{%
        \drawPolygon{#2}%
	\IfNoValueTF{#1}%
    	{%
            \pointlabel{#2}}%
    	{%
            \pointlabel{#1}}%
    	}%
    {%
        \pointlabel{#2}%
    }%
}%

%% Define bycolors for inline circles. Other uses of colors are defined in byrne.mp and byrnebook.cls
% RGB equivalents
\definecolor{byblue}{rgb}{.15, .35, .6}
\definecolor{byred}{rgb}{.85, .3, .1}
\definecolor{byyellow}{rgb}{.95, .7, 0.1}
\definecolor{bygreen}{rgb}{0.122, 0.529, 0.122}
\definecolor{byblack}{rgb}{0, 0, 0}
\definecolor{bygrey}{rgb}{.8, .8, .8}
\definecolor{bywhite}{rgb}{1, 1, 1}

% CMYK equivalents
%\definecolor{byred}{cmyk}{0.0, 0.88, 0.88, 0.15}
%\definecolor{byblue}{cmyk}{0.75, 0.417, 0.0, 0.4}
%\definecolor{byyellow}{cmyk}{0.0, 0.263, 0.895, 0.05}
%\definecolor{byblack}{cmyk}{0, 0, 0, 1}
%\definecolor{bygrey}{cmyk}{0.0, 0.0, 0.0, 0.2}

%%%%% drawCircleAlt
\DeclareDocumentCommand{\drawCircleAlt}{o m}{%
  \iftoggle{inline}{%
    \tikz[baseline=-1ex]{
      \draw[line width=1.6pt, draw=#1] (0,0) circle [radius=0.5em];
    }%
    \hspace{\circleSpace}\pointlabel{#2}
  }{%
    \pointlabel{#2}%
  }%
}

%%%%% ^drawCircleAlt

\def\drawPointAlt#1{%
    \iftoggle{inline}{%
        \drawPointM{#1}\hspace{\pointShift}\pointlabel{#1}%
        }%
    {%
    \pointlabel{#1}%
    }%
}%
%\def\drawAngleAlt#1{\drawAngle{#1}\pointlabel{#1}}

\DeclareDocumentCommand{\drawLineAlt}{o m}{% Alternative inline drawing command to draw connected sequence of lines (e.g. for triangles or polygons in general)
    \iftoggle{inline}{%
        \drawLineByPoints{#2}%
    \IfNoValueTF{#1}%
    {\pointlabel{\directlua{cleanPolygonName("#2")}}}%
    {\pointlabel{#1}}%
    }%
    {%
        \pointlabel{#2}%
    }%
}%

\def\RemoveCommas#1,{%
  \ifx\relax#1\relax% Check if we've reached the end
  \else%
    \textbf{#1}% Print in bold without a comma
    \expandafter\RemoveCommas% Continue with the next item
  \fi%
}%



%%%%%%%%%%%%%%%%%%%%%%%%%%%%%%%%%%%%%%%%

\lstset{
language=MetaPost,
alsolanguage=TeX,
numbers=none,
basicstyle=\ttfamily\scriptsize
}

\directlua{
function cleanPolygonName(s)
	local resultingName = s:gsub ("[, ]","")
	resultingName = resultingName:gsub ("-","minus")
	tex.print(resultingName)
end
}


%%%%%%%%%%%%%%%%%%%%%%%%%%%%%%%%%%%%%%%%%%%%%%%%%%%%%%%%%%%%%%%%%%%%
% The following is the original \drawLineAlt command, which had no optional argument to rename the figure
% \def\drawLineAlt#1{
%     \iftoggle{inline}{
%     \drawLineByPoints{#1}\pointlabel{\directlua{cleanPolygonName("#1")}}
%         }
%     {
%         \pointlabel{#1}
%     }
% }
%%%%%%%%%%%%%%%%%%%%%%%%%%%%%%%%%%%%%%%%%%%%%%%%%%%%%%%%%%%%%%%%%%%%

%%%%%%%%%%%%%%%%%%%%%%%%%%%%%%%%%%%%%%%%%%%%%%%%%%%%%%%%%%%%%%%%%%%%
%\def\drawPolygonAlt#1{\drawPolygon{#1}\pointlabel{#1}}
%drawPolygon[verticalAlignment][PolygonNewName] \{Poly\-gon\-Na\-me\}}
%%%%%%%%%%%%%%%%%%%%%%%%%%%%%%%%%%%%%%%%%%%%%%%%%%%%%%%%%%%%%%%%%%%%

%%%%%%%%%%%%%%%%%%%%%%%%%%%%%%%%%%%
% DBM attempt at new drawPolygonAlt command
\DeclareDocumentCommand{\drawMergedPolygonAlt}{o m}{% This is an alternative macro for polygons, allowing for merged polygons
    \iftoggle{inline}{%
        \drawFromCurrentPicture[middle][polygon#2]{%
            startAutoLabeling;%
            draw byNamedPolygon(#2);%
            stopAutoLabeling;%
            \pointlabel{#1}%
            }%
        {%
        \pointlabel{#1}%
        }%
    }%
}%
%%%%%%%%%%%%%%%%%%%%%%%%%%%%%%%%%%%


%%%%%%%%%%%%%%%%%%%%%%%%%%%%%%%%%%%%%%%%%%%%%%%%%%%%%%%%%%%%%%%%%%%%
% KDP 7x10 Left-Right Spread Formatting for Propositions (Fail-Safe)
%%%%%%%%%%%%%%%%%%%%%%%%%%%%%%%%%%%%%%%%%%%%%%%%%%%%%%%%%%%%%%%%%%%%
\usepackage{afterpage}

% Global state tracking
\newtoggle{inproposition}
\togglefalse{inproposition}
\newcommand{\currentpropdiagram}{}

% Automatically inserts the diagram on subsequent left (even) pages if text overflows
\newcommand{\insertEvenPageDiagram}{%
  \iftoggle{inproposition}{%
    % REMOVED the initial \clearpage that was breaking line-numbering and paragraph flow!
    \begin{center}
      \vspace*{\fill}
      \currentpropdiagram
      \vspace*{\fill}
    \end{center}
    \clearpage
    % Re-queue this check for the next page in case text overflows yet again
    \afterpage{\insertEvenPageDiagram}%
  }{}%
}

\newcommand{\startProposition}[2]{%
  \clearpage
  % Force an even (left-facing) page for the primary diagram
  \ifodd\value{page}\hbox{}\clearpage\fi 
  
  % Save the diagram code and set the active proposition state
  \renewcommand{\currentpropdiagram}{#1}
  \toggletrue{inproposition}
  
  % 1. Render the initial diagram on the left-facing page
  \begin{center}
    \vspace*{\fill}
    #1
    \vspace*{\fill}
  \end{center}
  \clearpage % Move cleanly to the right-facing (odd) page
  
  % 2. Defer queuing the even-page diagram until the actual text paragraph begins.
  \everypar{\everypar{}\afterpage{\insertEvenPageDiagram}}%
  
  % 3. Typeset the text block
  #2
  
  % 4. Clean up state before shipping out the final page to prevent trailing loops
  \togglefalse{inproposition}
  \clearpage
}

% Create Header Style
% 1. Create a global variable for the right-hand header
\newcommand{\currentpropheader}{Euclid in Color}

% 2. Define a native memoir page style (bypasses standard LaTeX marks)
\makepagestyle{euclidstyle}
\makeevenhead{euclidstyle}{\fontsize{12}{14}\selectfont Euclid in Color}{}{}
\makeoddhead{euclidstyle}{}{}{\fontsize{12}{14}\selectfont\rightmark}
\makeevenfoot{euclidstyle}{}{\thepage}{}
\makeoddfoot{euclidstyle}{}{\thepage}{}

% 3. Activate the style globally
\pagestyle{euclidstyle} % Input the preamble######.tex file


\begin{document}
% Title Page
\label{Title Page}
\begin{titlingpage}
    \centering
    \vspace*{2cm}
    {\fontsize{36pt}{40pt}\scshape\bfseries Euclid in Color\par}
    \vspace*{3cm}
    {\normalsize\scshape The Classic Translation of\par}
    {\Huge\scshape Euclid's \textit{Elements}, Book I\par}
    \vspace{0.2cm}
    {\normalsize\scshape by\par}
	{\Large\scshape Sir Thomas Heath\par}
    \vspace{0.3cm}
    {\normalsize\scshape with Diagrams Inspired by the Style of \par}
    {\Large\scshape Oliver Byrne\par}
    \vspace{2cm}
    {\large Edited by\par}
    \vspace{0.2cm}
    {\Large Daniel Murphy\par}
    \vspace{0.6cm}
    {\large with\par}
    \vspace{0.2cm}
    {\large Benjamin George \& Dillon Raum\par}
    \vfill
    {\Large 2026\par}
\end{titlingpage}

%\clearpage
%\thispagestyle{empty}

\vspace*{\fill}
\begin{center}
Copyright © 2026 \par
Version \bookversion
\end{center}
\vspace*{\fill}

\clearpage

\frontmatter % uses Roman numeral page numbering

\tableofcontents
\clearforchapter

\setlength{\parindent}{20pt} % include indents for new paragraphs in the Preface and Acknowledgements
\label{Preface}
\chapter*{Preface}
\addcontentsline{toc}{chapter}{Preface}
\markboth{Euclid in Color}{Preface}
\section*{Why this edition of Euclid's \textit{Elements}?}
\addcontentsline{toc}{section}{Why this edition of Euclid's \textit{Elements}?}

This book is an attempt to combine some of the positive features of two previous editions of Euclid's \textit{Elements} without thereby detracting from the text itself. Those two editions are Oliver Byrne's 1847 edition and Sir Thomas Heath's masterful 1908 translation. \\

In 1847, Oliver Byrne compiled an edition of the first six books of Euclid's \textit{Elements} with several unique features. His purpose was to help the learner acquire full knowledge of the propositions ``in less than one third the time usually employed.'' The features of Byrne's edition were colored diagrams, abbreviated text, and omission of letters from the diagrams and the text. All references to parts of a diagram were redrawn in the text of the proposition. \\

There are two features of Byrne's edition that we consider advantageous. Byrne's inclusion of color in the diagrams reflects a technique shared by many teachers of Euclid's \textit{Elements}. In the classroom, teachers will often draw parts of the diagrams in different colors to highlight importance or to facilitate making reference to ``the red line'' or ``the green angle.'' Byrne's inline redrawings of parts of the diagram reflect this same teaching technique. Whereas in the classroom, it is easier to speak of ``the red line'' or ``the green angle,'' in a text like Byrne's or this one, it is easier to simply draw a red line or a green angle. \\

The full adoption of Byrne's edition for classroom teaching, however, is not desirable if one wishes to teach Euclid's \textit{Elements} faithfully. One undesirable feature of Byrne's edition is that in many cases the colors of the lines and angles in his diagrams signify equality. In Proposition I.5 (Book I, Proposition 5), for example, some lines and angles that are given as equal have the same color \textit{and} some lines and angles that will be proven equal have the same color. While this technique is sometimes employed by teachers in the classroom, there are problems with doing it in a static text. Drawing lines or angles with the same color which at the outset are not yet proven equal is misleading at best. It can confuse students about what is or is not given, for example. It can also mislead students about the meaning of equality (that the color of lines has something to do with it). Generally speaking, Byrne's edition is inconsistent in this regard. This edition aims to use colors only for reference, not to show equality.\\

Another undesirable feature of Byrne's edition is the heavy modification of the texts of the propositions. We consider this a mistake. A text that has taught students for over two thousand years should be altered only with serious forethought. Byrne retained the logic of the propositions for the most part, but in many places he errs far too much on the side of brevity, and in some places he omits essential logical steps in pursuit of brevity and clarity.\\

Those of us lacking proficiency in ancient Greek must read great texts in translation. Fortunately, Sir Thomas Heath provided an excellent substitute for Euclid's original Greek in his masterful 1908 translation, which we use in this edition.
 
\section*{Advantages of this edition}
\addcontentsline{toc}{section}{Advantages of this edition}
Many of our decisions about how to format this edition were based on consideration of an excellent review written by Dr. Andrew Seeley\footnote{Seeley, Andrew. “Oliver Byrne’s Elements of Euclid: The First Six Books with Coloured Diagrams and Symbols.” Principia: A Journal of Classical Education, vol. 2, no. 1, 2023, pp. 107–112, https://doi.org/10.5840/principia2023219.} in \textit{Principia}. Seeley's review is largely negative. We agree with most of Seeley's criticisms, and we hope that we have produced a version that Seeley would prefer to that of Oliver Byrne.\\

Much of Seeley's criticism centers on the modification of the text of Euclid's \textit{Elements}. Since we include Heath's translation in its entirety, these particular criticisms do not apply to this edition. Our inclusion of the entirety of Heath's translation, for example, makes it possible to construct a diagram entirely from the text itself.\\

Byrne's edition of the \textit{Elements} omits letters from proposition diagrams and text. Seeley takes issue with this because the letters reinforce the importance of points as the principles of geometry. We chose to take Seeley's advice and include letters in both the diagrams and the text. \\

Another of Seeley's criticisms that we took into account was Byrne's ``egregious'' treatment of Postulate 5. Seeley points out that Byrne's diagram represents parallel lines, whereas Postulate 5 is about lines that meet. We chose to alter Byrne's diagram accordingly.\\

This version of the \textit{Elements} adds shading and color to previous editions of Heath's translation. Even the diagrams are drawn similar to those found in the Green Lion Press and Cambridge University Press editions of Heath's work. Andrew Seeley closes his review of Byrne's edition with the following remark, which we believe is precisely what we have done:
\begin{quotation}
\noindent\itshape
\hyphenpenalty=10000\exhyphenpenalty=10000
Euclid is difficult for the beginner. It is often helpful to make his propositions more intuitive with visual aids, such as shading or coloring.
\end{quotation}

\section*{Using this edition}
\addcontentsline{toc}{section}{Using this edition}
\textit{Euclid in Color} only contains the first book of the \textit{Elements}, an introduction to the geometry and arithmetic of Euclid's thirteen-book masterpiece. It is intended for one approaching the \textit{Elements} for the first time, perhaps with an understandable degree of trepidation. The colorization of the diagrams does not replace but merely facilitates the logical thinking the propositions require. Byrne's abbreviated text allows a student to follow a proposition without ever working through its full argument, which can produce the appearance of understanding rather than understanding itself. Heath's complete translation offers no such shortcut. Every step Euclid takes is on the page, and a student who follows it has genuinely followed the proof and hopefully understood it. This is why we include Heath's translation in its entirety. \\

Since many have already studied Euclid using Heath's translation and diagrams, experienced teachers should find little challenge in adapting their classes to this edition. The similarity of this text to Heath's will also soften the transition from a colorized Book I to a non-colorized Book II and beyond. \\

A companion guide for teachers and independent learners is planned for release by the end of 2026. It will treat the definitions, postulates, and common notions in greater depth, and it will explain the recurring structure that every proposition in the \textit{Elements} shares. Readers who would like this material sooner can find much of it at \textit{What Schools Forget}, the editor's newsletter on education and the teaching of mathematics, at \texttt{whatschoolsforget.substack.com}.


\label{Acknowledgements}
\chapter*{Acknowledgements}
\renewcommand{\currentpropheader}{Acknowledgements}
\addcontentsline{toc}{chapter}{Acknowledgements}
This project was jointly conceived by Ben George, Dillon Raum, and Daniel Murphy. Daniel Murphy served as the lead editor, responsible for the \LaTeX\ programming, typesetting, and final production. Ben George and Dillon Raum reviewed drafts and gave numerous conceptual contributions. \\

Special thanks to Sergey Slyusarev, whose publicly available work provided the foundational material for this adaptation: \par
\noindent https://github.com/jemmybutton/byrne-euclid. Sergey also provided important help with the early stages of \LaTeX\ programming for the modified inline drawings. \\

The text for this version of Euclid's \textit{Elements} was taken from the Perseus Digital Library, based on the translation by Sir Thomas L. Heath (1908). The stable URI for this work is: \par 
\noindent http://data.perseus.org/catalog/urn:cts:greekLit:tlg1799.tlg001.perseus-eng1. \\

For inquiries, contact Daniel Murphy at danbmurphy@gmail.com.

\setlength{\parindent}{0pt} % Remove indenting for text of Elements

\mainmatter % Uses Hindu-Arabic page-numbering

\label{Definitions}
\chapter*{Definitions}
\addcontentsline{toc}{chapter}{Definitions}
\renewcommand{\currentpropheader}{Definitions}
\begin{enumerate}
	\setlength{\itemsep}{0pt}
    \setlength{\parsep}{0pt}
    \setlength{\topsep}{4pt}
    \item A \emph{point} is that which has no part.\label{Def:I.I}
    \item A \emph{line} is a breadthless length.\label{Def:I.II}
    \item The extremities of a line are points.\label{Def:I.III}
    \item A \emph{straight line} is that which lies evenly with the points on itself.\label{Def:I.IV}
    \item A \emph{surface} is that which has length and breadth only.\label{Def:I.V}
    \item The extremities of a surface are lines.\label{Def:I.VI}
    \item A \emph{plane surface} is a surface which lies evenly with the straight lines on itself.\label{Def:I.VII}
    \item A \emph{plane angle} is the inclination to one another of two lines in a plane which meet one another and do not lie in a straight line.\label{Def:I.VIII}
    \item And when the lines containing the angle are straight, the angle is called \emph{rectilineal}.\label{Def:I.IX} 
        
        \defineNewPicture{
	pair A, B, C;
	A := (0, 0);
	B := (2/3u, 2/3u);
	C := (u, ypart(A));
		byAngleDefine(B, A, C, byyellow, SOLID_SECTOR);
		draw byNamedAngleResized();
		byLineDefine(A, B, byblue, SOLID_LINE, REGULAR_WIDTH);
		byLineDefine(A, C, byred, SOLID_LINE, REGULAR_WIDTH);
		draw byNamedLineSeq(0)(AC,AB);
        }
        \iftoggle{defpics}{~\hfill\drawCurrentPicture\hfill~}{}
    \item When a straight line set up on a straight line makes the adjacent angles equal to one another, each of the equal angles is \emph{right}, and the straight line standing on the other is called a \emph{perpendicular} to that on which it stands.\label{Def:I.X}
        
        \defineNewPicture{
	pair A, B, C, D;
	A := (0, 0);
	B := (u, 0);
	C := (0, 2/3u);
	D := (-u, 0);
		byAngleDefine(B, A, C, byblack, ARC_SECTOR);
		byAngleDefine(D, A, C, byblack, ARC_SECTOR);
		%draw byNamedAngleResized();
		draw byLine(D, B, byblack, SOLID_LINE, REGULAR_WIDTH);
		draw byLine(A, C, byblack, SOLID_LINE, REGULAR_WIDTH);
        }
        \iftoggle{defpics}{~\hfill\drawCurrentPicture\hfill~}{}
    \item An \emph{obtuse angle} is an angle greater than a right angle.\label{Def:I.XI}
        
        \defineNewPicture{
	pair A, B, C;
	A := (0, 0);
	B := (-2/3u, 2/3u);
	C := (u, ypart(A));
		byAngleDefine(B, A, C, byred, SOLID_SECTOR);
		draw byNamedAngleResized();
		byLineDefine(A, B, byyellow, SOLID_LINE, REGULAR_WIDTH);
		byLineDefine(A, C, byblue, SOLID_LINE, REGULAR_WIDTH);
		draw byNamedLineSeq(0)(AC,AB);
        }
        \iftoggle{defpics}{~\hfill\drawCurrentPicture\hfill~}{}
    \item An \emph{acute angle} is an angle less than a right angle.\label{Def:I.XII}
        
        \defineNewPicture{
	pair A, B, C;
	A := (0, 0);
	B := (2/3u, 2/3u);
	C := (u, ypart(A));
		byAngleDefine(B, A, C, byblue, SOLID_SECTOR);
		draw byNamedAngleResized();
		byLineDefine(A, B, byyellow, SOLID_LINE, REGULAR_WIDTH);
		byLineDefine(A, C, byred, SOLID_LINE, REGULAR_WIDTH);
		draw byNamedLineSeq(0)(AC,AB);
        }
        \iftoggle{defpics}{~\hfill\drawCurrentPicture\hfill~}{}
    \item A \emph{boundary} is that which is an extremity of anything.\label{Def:I.XIII}
    \item A \emph{figure} is that which is contained by any boundary or boundaries.\label{Def:I.XIV}
    \item A \emph{circle} is a plane figure contained by one line such that all the straight lines falling upon it from one point among those lying within the figure are equal to one another;\label{Def:I.XV}
        
        \defineNewPicture{
	pair O, A, B, C, D, E;
	numeric r;
	r := 1/2u;
	O := (0, 0);
	A := dir(0) scaled r;
	B := dir(60) scaled r;
	C := dir(130) scaled r;
	D := dir(180) scaled r;
	E := dir(-60) scaled r;
		draw byLine(O, B)(byblack, SOLID_LINE, REGULAR_WIDTH);
		draw byLine(O, C)(byred, SOLID_LINE, REGULAR_WIDTH);
		draw byLine(O, E)(byyellow, SOLID_LINE, REGULAR_WIDTH);
		draw byLine(D, A)(byblue, SOLID_LINE, REGULAR_WIDTH);
		draw byCircleR(O, r, byred, 0, 0, 0);
        }
        \iftoggle{defpics}{~\hfill\drawCurrentPicture\hfill~}{}
    \item And the point is called the \emph{center} of the circle.\label{Def:I.XVI}
    \item A \emph{diameter} of the circle is any straight line drawn through the center and terminated in both directions by the circumference of the circle, and such a straight line also bisects the circle.\label{Def:I.XVII}
        
        \defineNewPicture{
	pair O, A, B;
	numeric r;
	r := 1/2u;
	O := (0, 0);
	A := dir(0) scaled r;
	B := dir(180) scaled r;
		draw byLine(A, B)(byyellow, SOLID_LINE, REGULAR_WIDTH);
		draw byCircleR(O, r, byred, 0, 0, 0);
        }
        \iftoggle{defpics}{~\hfill\drawCurrentPicture\hfill~}{}
    \item A \emph{semicircle} is the figure contained by the diameter and the circumference cut off by it. And the center of the semicircle is the same as that of the circle.\label{Def:I.XVIII}
        
        \defineNewPicture{
	pair O, A, B;
	numeric r;
	r := 1/2u;
	O := (0, 0);
	A := dir(0) scaled r;
	B := dir(180) scaled r;
		draw byLine(A, B)(byblue, SOLID_LINE, REGULAR_WIDTH);
		draw byArc(O, A, B)(r, byyellow, 0, 0, 0, 0);
		draw byArc(O, B, A)(r, byyellow, 1, 0, 0, 0);
        }
        \iftoggle{defpics}{~\hfill\drawCurrentPicture\hfill~}{}
    \item \emph{Rectilineal figures} are those which are contained by straight lines, \emph{trilateral} figures being those contained by three, \emph{quadrilateral} those contained by four, and \emph{multilateral} those contained by more than four straight lines.\label{Def:I.XIX}
    \item Of trilateral figures, an \emph{equilateral triangle} is that which has its three sides equal, an \emph{isosceles triangle} that which has two of its sides alone equal, and a \emph{scalene triangle} that which has its three sides unequal.\label{Def:I.XX}
    \item Further, of trilateral figures, a \emph{right-angled triangle} is that which has a right angle, an \emph{obtuse-angled triangle} that which has an obtuse angle, and an \emph{acute-angled triangle} that which has its three angles acute.\label{Def:I.XXI}
    \item Of quadrilateral figures, a \emph{square} is that which is both equilateral and right-angled; an \emph{oblong} (or ``rectangle'') that which is right-angled but not equilateral; a \emph{rhombus} that which is equilateral but not right-angled; and a \emph{rhomboid} (or ``parallelogram'') that which has its opposite sides and angles equal to one another but is neither equilateral nor right-angled. And let quadrilaterals other than these be called \emph{trapezia}.\label{Def:I.XXII}
    \item \emph{Parallel} straight lines are straight lines which, being in the same plane and being produced indefinitely in both directions, do not meet one another in either direction.\label{Def:I.XXIII}
\end{enumerate}

\chapter*{Postulates}
\renewcommand{\currentpropheader}{Postulates}
\addcontentsline{toc}{chapter}{Postulates}
Let the following be postulated:
\begin{enumerate}
    \item To draw a straight line from any point to any point. \label{post:I.I}
    \item To produce a finite straight line continuously in a straight line. \label{post:I.II}
    \item To describe a circle with any center and distance. \label{post:I.III} 
    \item That all right angles are equal to one another. \label{post:I.IV}
    \item That, if a straight line falling on two straight lines make the interior angles on the same side less than two right angles, the two straight lines, if produced indefinitely, meet on that side on which are the angles less than the two right angles. \label{post:I.V}
    
        \defineNewPicture{
    	pair A, B, C, D, E, F, G, H;
    	numeric s;
    	s := 3/2u;
    	A := (0, 0);
    	B := (6s, s/2);
    	C := (0, s);
    	D := (6s, s/2);
    	E := (s, 8/6s);
    	F := (xpart(E), -2/6s);
    	G = whatever[A, B] = whatever[E, F];
    	H = whatever[C, D] = whatever[E, F];
    		byAngleDefine(B, G, E, byred, SOLID_SECTOR);
    		byAngleDefine(D, H, F, byyellow, SOLID_SECTOR);
    		draw byNamedAngleResized();
    		draw byLine(A, B, byblue, SOLID_LINE, REGULAR_WIDTH);
    		draw byLine(C, D, byred, SOLID_LINE, REGULAR_WIDTH);
    		draw byLine(E, F, byblack, SOLID_LINE, REGULAR_WIDTH);
        }
        \iftoggle{defpics}{~\hfill\drawCurrentPicture\hfill~}{}
\end{enumerate}

\chapter*{Common Notions}
 \renewcommand{\currentpropheader}{Common Notions}
\addcontentsline{toc}{chapter}{Common Notions}
\begin{enumerate}
    \item Things which are equal to the same thing are also equal to one another. \label{cn:I.I}
    \item If equals be added to equals, the wholes are equal. \label{cn:I.II}
    \item If equals be subtracted from equals, the remainders are equal. \label{cn:I.III}
    \item Things which coincide with one another are equal to one another. \label{cn:I.IV}
    \item The whole is greater than the part. \label{cn:I.V}
\end{enumerate}

\resetlinenumber

\startProposition{
	\defineNewPicture[1][5/2]{
		textLabels := true;
		pair A, B, C, D, E;
		path P[];
		numeric r;
		r := 3/2u;
		A := (0, 0);
		B := (r, 0);
		P1 := fullcircle scaled 2r;
		P2 := fullcircle scaled 2r shifted B;
		C := P1 intersectionpoint P2;
		D := A shifted (unitvector(A-B) scaled r); % defining points D and E relative to the other points
		E := B shifted (unitvector(B-A) scaled r);
		byLineDefine(A, B, byblack, SOLID_LINE, REGULAR_WIDTH);
		byLineDefine(B, C, byred, SOLID_LINE, REGULAR_WIDTH);
		byLineDefine(C, A, byyellow, SOLID_LINE, REGULAR_WIDTH);
		draw byNamedLineSeq(-1)(AB,CA,BC);
		draw byCircleABC(B, C, D, byblue, 0, 0, 1/2); % `byCircleABC` defines a circle by three points on its circumference
		draw byCircleABC(A, C, E, byred, 0, 0, 1/2);
		%draw byCircleABC(D, C, E, byyellow, 0, 0, 1/2);
		draw byLabelsOnPolygon(A, C, B)(OMIT_FIRST_LABEL+OMIT_LAST_LABEL, 1);
		draw byLabelsOnCircle(B, D)(BCD);
		draw byLabelsOnCircle(A, E)(ACE);
		textLabels := false;
	}
	~\hfill\drawCurrentPicture\hfill~
}{
	\chapter*{Proposition I.1}
    \renewcommand{\currentpropheader}{Proposition I.1}
	\addcontentsline{toc}{chapter}{Proposition I.1}
	\textit{On a given finite straight line to construct an equilateral triangle.}
	\vskip 0.5\baselineskip
	
	% Check line numbering toggle
	\iftoggle{lineNumbering}{\linenumbers}{}

	\iftoggle{heath}{
	    Let \drawUnitLineAlt{AB} be the given finite straight line.\par 
	    Thus it is required to construct an equilateral triangle on the straight line \drawUnitLineAlt{AB}. 
	    \vskip \baselineskip 
	    With center \drawPointAlt{A} and distance \drawUnitLineAlt{AB} let the circle \drawCircleAlt[byblue]{BCD} be described;\justification{Post. 3}
	    again, with center \drawPointAlt{B} and distance \drawUnitLineAlt{BA} let the circle \drawCircleAlt[byred]{ACE} be described;\par\justification{Post. 3}
	    and from the point \drawPointAlt{C}, in which the circles cut one another, to the points \drawPointAlt{A}, \drawPointAlt{B} let the straight lines \drawUnitLineAlt{CA}, \drawUnitLineAlt{CB} be joined.\justification{Post. 1}
	    Now, since the point \drawPointAlt{A} is the center of the circle \drawCircleAlt[byblue]{CDB}, \par
	    \centeredjust{\drawUnitLineAlt{AC} is equal to \drawUnitLineAlt{AB}.}{Def. 15}
	    Again, since the point \drawPointAlt{B} is the center of the circle \drawCircleAlt[byred]{CAE},\par
	    \centeredjust{\drawUnitLineAlt{BC} is equal to \drawUnitLineAlt{BA}.}{Def. 15}
	    But \drawUnitLineAlt{CA} was also proved equal to \drawUnitLineAlt{AB};\par
	    {\centering therefore each of the straight lines \drawUnitLineAlt{CA}, \drawUnitLineAlt{CB} is equal to \drawUnitLineAlt{AB}. \par}
	    And things which are equal to the same thing are also equal to one another; \par\justification{C.N. 1}
	    {\centering therefore \drawUnitLineAlt{CA} is also equal to \drawUnitLineAlt{CB}. \par}
	    Therefore the three straight lines \drawUnitLineAlt{CA}, \drawUnitLineAlt{AB}, \drawUnitLineAlt{BC} are equal to one another.\par
	    Therefore the triangle \drawLineAlt{A,B,C} is equilateral; and it has been constructed on the given finite straight line \drawUnitLineAlt{AB}.\par
	    \hfill(Being) what it was required to do.
	    }{}
	\iftoggle{raum}{
	    RaumText
	}{}
	\nolinenumbers
}
\resetlinenumber

\startProposition{
    \defineNewPicture[1][2]{
    textLabels := true;
	pair A, B, C, D, E, F, G, H, K, L;
	path P[];
	numeric r[];
	A := (-3/5u, -2/5u);
	B := (0, 0);
	C := (0, 2u);
	r1 := abs(A-B);
	D := (fullcircle scaled 2r1 shifted A) intersectionpoint (fullcircle scaled 2r1 shifted B); % Wow this is prop 1 in mp
	E := 14/3[D, A];
	F := 14/3[D, B];
	r2 := abs(B-C);
	r3 := r1 + r2;
	P1 := fullcircle scaled 2r2 shifted B; % Circle BC
	P2 := fullcircle scaled 2r3 shifted D; % Circle DG
	G := (D -- F) intersectionpoint P1;
	H := (B -- 10[B,D]) intersectionpoint P1;
	K := (B -- 10[B,D]) intersectionpoint P2;
	L := (D -- E) intersectionpoint P2;
	byLineDefine(A, B, byred, SOLID_LINE, REGULAR_WIDTH);
	byLineDefine(B, C, byblack, SOLID_LINE, REGULAR_WIDTH);
	byLineDefine(B, D, byblue, SOLID_LINE, REGULAR_WIDTH);
	byLineDefine(D, A, byyellow, SOLID_LINE, REGULAR_WIDTH); % improvement: change style of either BD or DA?
	byLineDefine(L, E, bygreen, SOLID_LINE, REGULAR_WIDTH);
	byLineDefine(G, F, byyellow, SOLID_LINE, REGULAR_WIDTH);
	byLineDefine(B, G, byyellow, SOLID_LINE, REGULAR_WIDTH);
	byLineDefine(A, L, bygreen, SOLID_LINE, REGULAR_WIDTH);
	draw byNamedLineSeq(0)(AB,BC);
	draw byNamedLineSeq(0)(BG,BD,DA,AL,LE,GF);
	draw byCircleABC(C, G, H, byblue, 0, 0, -1/2);
	draw byCircleABC(G, K, L, byred, 0, 0, 1/2);
	draw byLabelsOnPolygon(A, D, B)(0, 1);
	draw byLabelsOnCircle(H, C)(CGH);
	draw byLabelsOnCircle(K)(GKL);
	draw byLabelPoint(L, 60, -4);
	draw byLabelPoint(E, 90, -3);
	draw byLabelPoint(G, 15, 2);
	draw byLabelPoint(F, -30, 2);
	%draw byLabelPoint(H, 0, -2);
    }
    ~\hfill\drawCurrentPicture\hfill~
}{
    \chapter*{Proposition I.2}
    \renewcommand{\currentpropheader}{Proposition I.2}
    \addcontentsline{toc}{chapter}{Proposition I.2}
    \textit{To place at a given point (as an extremity) a straight line equal to a given straight line.}
	\vskip 0.5\baselineskip

    % Check line numbering toggle
\iftoggle{lineNumbering}{\linenumbers}{}

% Prop text
\iftoggle{heath}{
    Let \drawPointAlt{A} be the given point, and \drawUnitLineAlt{BC} the given straight line. 
    Thus it is required to place at the point \drawPointAlt{A} (as an extremity) a straight line equal to the given straight line\drawUnitLineAlt{BC}. 
    \vskip \baselineskip 
    From the point \drawPointAlt{A} to the point \drawPointAlt{B} let the straight line\drawUnitLineAlt{AB} be joined;\justification{Post. 1}
    and on it let the equilateral triangle \drawLineAlt{D,A,B} be constructed.\justification{I. 1}
    Let the straight lines \drawUnitLineAlt{AE}, \drawUnitLineAlt{BF} be produced in a straight line with \drawUnitLineAlt{DA}, \drawUnitLineAlt{DB}; \hfill\llap{} \par
    with center \drawPointAlt{B} and distance \drawUnitLineAlt{BC} let the circle \drawCircleAlt[byblue]{CGH} be described;\justification{Post. 3}
    and again, with center \drawPointAlt{D} and distance \drawUnitLineAlt{DG} let the circle \drawCircleAlt[byred]{GKL} be described.\par\justification{Post. 3}
    Then, since the point \drawPointAlt{B} is the center of the circle \drawCircleAlt[byblue]{CGH}, \par
    {\centering \drawUnitLineAlt{BC} is equal to \drawUnitLineAlt{BG}. \par}
    Again, since the point \drawPointAlt{D} is the center of the circle \drawCircleAlt[byred]{GKL},\par
    {\centering \drawUnitLineAlt{DL} is equal to \drawUnitLineAlt{DG}. \par}
    And in these \drawUnitLineAlt{DA} is equal to \drawUnitLineAlt{DB};\par
    \centeredjust{therefore the remainder \drawUnitLineAlt{AL} is equal to the remainder \drawUnitLineAlt{BG}.}{C.N. 3}
    But \drawUnitLineAlt{BC} was also proved equal to \drawUnitLineAlt{BG}; \par
    {\centering therefore each of the straight lines \drawUnitLineAlt{AL}, \drawUnitLineAlt{BC} is equal to \drawUnitLineAlt{BG}. \par}   
    And things which are equal to the same thing are also equal to one another;\par\justification{C.N. 1}
    {\centering therefore \drawUnitLineAlt{AL} is also equal to \drawUnitLineAlt{BC}. \par}
    Therefore at the given point \drawPointAlt{A} the straight line \drawUnitLineAlt{AL} is placed equal to the given straight line \drawUnitLineAlt{BC}. \par
    \hfill(Being) what it was required to do.
}{}

\iftoggle{raum}{
    Raum_Text
}{}

\nolinenumbers
\vskip \baselineskip


\resetlinenumber
}

\startProposition{
    \markboth{Euclid in Color}{Proposition I.3}
    \defineNewPicture[1][3]{
    textLabels := true;
	pair A, B, C, D, E, F, G;
	path P;
	numeric r;
	A := (0, 0);
	r := 7/4u;
	E := A shifted (r, 0);
	B := A shifted (4/3r, 0);
	D := A shifted dir(135)*r;
	C := A shifted (-1/2r, 6/5r);
	F := 2[D,A];
	G := A shifted (1/2r, 6/5r);
	byLineDefine(A, E, byblack, SOLID_LINE, REGULAR_WIDTH);
	byLineDefine(E, B, byblack, SOLID_LINE, REGULAR_WIDTH);
	byLineDefine(A, D, byyellow, SOLID_LINE, REGULAR_WIDTH);
	draw byNamedLineSeq(0)(EB,AE,AD);
	draw byLine.C(C, G, byred, SOLID_LINE, REGULAR_WIDTH);
	draw byCircleABC(D, E, F, byblue, 0, 0, 0);
	%draw byCircle.A(A, D, byblue, 0, 0, 0);
	draw byLabelsOnPolygon(E, A, D)(OMIT_FIRST_LABEL+OMIT_LAST_LABEL, 1);
	draw byLabelLineEnd(D, A, 0);
	draw byLabelLineEnd(B, A, 0);
	draw byLabelPoint(E, angle(E-A) + 45, 2);
	draw byLabelPoint(F, - 45, 2);
	draw byLabelsOnPolygon(G, C)(OMIT_FIRST_LABEL, 0);
	textLabels := false;
    }
    ~\hfill\drawCurrentPicture\hfill~
}{
    \chapter*{Proposition I.3}
    \renewcommand{\currentpropheader}{Proposition I.3}
    \addcontentsline{toc}{chapter}{Proposition I.3}
    \textit{Given two unequal straight lines, to cut off from the greater a straight line equal to the less.}
	\vskip 0.5\baselineskip

    % Check line numbering toggle
\iftoggle{lineNumbering}{\linenumbers}{}

% Prop text
\iftoggle{heath}{
    Let \drawUnitLineAlt{AB}, \drawUnitLineAlt{C} be the two given unequal straight lines, and let \drawUnitLineAlt{AB} be the greater of them.
    Thus it is required to cut off from \drawUnitLineAlt{AB} the greater a straight line equal to \drawUnitLineAlt{C} the less.
    \vskip \baselineskip
    At the point  \drawPointAlt{A} let \drawUnitLineAlt{AD} be placed equal to the straight line \drawUnitLineAlt{C}; \justification{I. 2}
    and with center  \drawPointAlt{A} and distance \drawUnitLineAlt{AD} let the circle \drawCircleAlt[byblue]{DEF} be described. \par\justification{Post. 3}
    Now, since the point  \drawPointAlt{A} is the center of the circle \drawCircleAlt[byblue]{DEF},\par
    \centeredjust{\drawUnitLineAlt{AE} is equal to \drawUnitLineAlt{AD}.}{Def. 15}
    But \drawUnitLineAlt{C} is also equal to \drawUnitLineAlt{AD}. \par
    Therefore each of the straight lines \drawUnitLineAlt{AE}, \drawUnitLineAlt{C} is equal to \drawUnitLineAlt{AD}; \par
    \centeredjust{so that \drawUnitLineAlt{AE} is also equal to \drawUnitLineAlt{C}.}{C.N. 1}
    Therefore, given the two straight lines \drawUnitLineAlt{AB}, \drawUnitLineAlt{C}, from \drawUnitLineAlt{AB} the greater \drawUnitLineAlt{AE} has been cut off equal to \drawUnitLineAlt{C} the less. \par
    \hfill(Being) what it was required to do.
}{}

\iftoggle{raum}{
    Raum_Text
}{}

\nolinenumbers
}

\resetlinenumber

\startProposition{
	\markboth{Euclid in Color}{Proposition I.4}    
    \defineNewPicture[1/4][3/2]{
    textLabels := true;
	pair A, B, C, D, E, F, G, H, K, d, g;
	path P[];
	numeric r[];
	A := (0, 0);
	B := A shifted (-5/2u, -7/2u);
	C := A shifted (1/2u, -3u);
	d := (5u, 0);
	D := A shifted d;
	E := B shifted d;
	F := C shifted d;
	byAngleDefine(B, A, C, byyellow, SOLID_SECTOR);
	byAngleDefine(A, B, C, byblue, SOLID_SECTOR);
	byAngleDefine(B, C, A, byblack, SOLID_SECTOR);
	draw byNamedAngleResized(BAC, ABC, BCA);
	byLineDefine(A, B, byred, SOLID_LINE, REGULAR_WIDTH);
	byLineDefine(B, C, bygreen, SOLID_LINE, REGULAR_WIDTH);
	byLineDefine(C, A, byblue, SOLID_LINE, REGULAR_WIDTH);
	draw byNamedLineSeq(0)(CA,BC,AB);
	byAngleDefine(E, D, F, byred, SOLID_SECTOR);
	byAngleDefine(D, E, F, byblack, SOLID_SECTOR);
	byAngleDefine(E, F, D, byblue, SOLID_SECTOR);
	draw byNamedAngleResized(EDF, DEF, EFD);
	byLineDefine(D, E, bygreen, SOLID_LINE, REGULAR_WIDTH);
	byLineDefine(E, F, byyellow, SOLID_LINE, REGULAR_WIDTH);
	byLineDefine(F, D, byblack, SOLID_LINE, REGULAR_WIDTH);
	draw byNamedLineSeq(0)(FD,EF,DE);
	draw byLabelsOnPolygon(F, E, D)(ALL_LABELS, 0);
	draw byLabelsOnPolygon(B, A, C)(ALL_LABELS, 1);
	%%% Total hack job getting the arc on EF to line up properly...
	r1 := 2abs(A-F);
	P1 := fullcircle scaled 2r1 shifted F;
	P2 := fullcircle scaled 2r1 shifted E;
	G := P1 intersectionpoint P2;
	H := 1/2[E,F];
	K := 1/2[H,G];
	byLineDefine(K, F, byblue, SOLID_LINE, REGULAR_WIDTH);
	byArcDefine.EF(K, E, F, r1/2, byyellow, SOLID_LINE, REGULAR_WIDTH, 2, 0);
	draw byNamedArcExact(EF);
    }
    \drawCurrentPicture
}{
    \chapter*{Proposition I.4}
    \renewcommand{\currentpropheader}{Proposition I.4}
    \addcontentsline{toc}{chapter}{Proposition I.4}
    \textit{If two triangles have the two sides equal to two sides respectively, and have the angles contained by the equal straight lines equal, they will also have the base equal to the base, the triangle will be equal to the triangle, and the remaining angles will be equal to the remaining angles respectively, namely those which the equal sides subtend.}
	\vskip 0.5\baselineskip

    % Check line numbering toggle
\iftoggle{lineNumbering}{\linenumbers}{}

% Prop text
\iftoggle{heath}{
    Let \drawLineAlt{A,B,C}, \drawLineAlt{D,E,F} be two triangles having the two sides \drawUnitLineAlt{AB},\drawUnitLineAlt{AC} equal to the two sides \drawUnitLineAlt{DE},\drawUnitLineAlt{DF} respectively, namely \drawUnitLineAlt{AB} to \drawUnitLineAlt{DE} and \drawUnitLineAlt{AC} to \drawUnitLineAlt{DF}, and the angle \drawAngleAlt{BAC} equal to the angle \drawAngleAlt{EDF}. \par
    I say that the base \drawUnitLineAlt{BC} is also equal to the base \drawUnitLineAlt{EF}, the triangle \drawLineAlt{A,B,C} will be equal to the triangle \drawLineAlt{D,E,F}, and the remaining angles will be equal to the remaining angles respectively, namely those which the equal sides subtend, that is, the angle \drawAngleAlt{ABC} to the angle \drawAngleAlt{DEF}, and the angle \drawAngleAlt{ACB} to the angle \drawAngleAlt{DFE}. 
    \vskip \baselineskip
    For, if the triangle \drawLineAlt{A,B,C} be applied to the triangle \drawLineAlt{D,E,F}, and if the point  \drawPointAlt{A} be placed on the point  \drawPointAlt{D} and the straight line \drawUnitLineAlt{AB} on \drawUnitLineAlt{DE}, then the point \drawPointAlt{B} will also coincide with  \drawPointAlt{E}, because \drawUnitLineAlt{AB} is equal to \drawUnitLineAlt{DE}. \par
    Again,\drawUnitLineAlt{AB} coinciding with \drawUnitLineAlt{DE}, the straight line \drawUnitLineAlt{AC} will also coincide with \drawUnitLineAlt{DF}, because the angle \drawAngleAlt{BAC} is equal to the angle \drawAngleAlt{EDF}; \par
    hence the point  \drawPointAlt{C} will also coincide with the point  \drawPointAlt{F}, because \drawUnitLineAlt{AC} is again equal to \drawUnitLineAlt{DF}. \par
    But \drawPointAlt{B} also coincided with  \drawPointAlt{E}; \par
    {\centering hence the base \drawUnitLineAlt{BC} will coincide with the base \drawUnitLineAlt{EF}, \par}
    [\textit{For if, when \drawPointAlt{B} coincides with  \drawPointAlt{E} and  \drawPointAlt{C} with  \drawPointAlt{F}, the base \drawUnitLineAlt{BC} does not coincide with the base \drawUnitLineAlt{EF}, two straight lines will enclose a space: which is impossible. Therefore the base \drawUnitLineAlt{BC} will coincide with \drawUnitLineAlt{EF}}] \par
    \centeredjust{and will be equal to it.}{C.N. 4}
    \nolinenumbers
    \clearpage
    \linenumbers
    Thus the whole triangle \drawLineAlt{A,B,C} will coincide with the whole triangle \drawLineAlt{D,E,F}, \par
    \centeredjust{and will be equal to it.}{C.N. 4}
    And the remaining angles will also coincide with the remaining angles and will be equal to them, \par
    {\centering the angle \drawAngleAlt{ABC} to the angle \drawAngleAlt{DEF}, \par}
    \centeredjust{and the angle \drawAngleAlt{ACB} to the angle \drawAngleAlt{DFE}.}{C.N. 4}
    Therefore etc. \par
    \hfill (Being) what it was required to prove.
}{}

\iftoggle{raum}{
    Raum_Text
}{}

\nolinenumbers

}

\resetlinenumber

\startProposition{
    \markboth{Euclid in Color}{Proposition I.5}
    \defineNewPicture[1/3][2]{
    textLabels := true;
	angleSize := 1cm;
	pair A, B, C, D, E, F, G;
	A := (0, 0);
	B := A shifted (-1.5u, -3u);
	C := B xscaled -1;
	D := 15/8[A,B];
	E := 15/8[A,C];
	F := 3/2[A,B];
	G := 3/2[A,C];
	byAngleDefine(B, A, C, byyellow, SOLID_SECTOR);
	byAngleDefine(A, B, C, byred, SOLID_SECTOR);
	byAngleDefine(B, C, A, byblue, SOLID_SECTOR);
	byAngleDefine(C, B, G, bygreen, SOLID_SECTOR);
	byAngleDefine(B, C, F, byyellow, SOLID_SECTOR);
	byAngleDefine(B, G, C, bygreen, SOLID_SECTOR);
	byAngleDefine(C, F, B, byyellow, SOLID_SECTOR);
	byAngleDefine(G, B, D, byblack, SOLID_SECTOR);
	byAngleDefine(F, C, E, byred, SOLID_SECTOR);
	draw byNamedAngleResized();
	byLineDefine(B, G, byred, SOLID_LINE, REGULAR_WIDTH);
	byLineDefine(C, F, bygreen, SOLID_LINE, REGULAR_WIDTH);
	byLineDefine(A, B, bygreen, SOLID_LINE, REGULAR_WIDTH);
	byLineDefine(A, C, byred, SOLID_LINE, REGULAR_WIDTH);
	byLineDefine(B, C, byblack, SOLID_LINE, REGULAR_WIDTH);
	byLineDefine(B, F, byblue, SOLID_LINE, REGULAR_WIDTH);
	byLineDefine(F, D, byblue, SOLID_LINE, REGULAR_WIDTH);
	byLineDefine(C, G, byblack, SOLID_LINE, REGULAR_WIDTH);	
	byLineDefine(G, E, byblack, SOLID_LINE, REGULAR_WIDTH);
	draw byNamedLineSeq(0)(AC,AB,BC);
	draw byNamedLineSeq(0)(FD,GE,BG,CF,BF,CG);
	draw byLabelsOnPolygon(D, F, B, A, C, G, E)(ALL_LABELS, 1);
    }
    ~\hfill\drawCurrentPicture\hfill~
}{
    \chapter*{Proposition I.5}
    \renewcommand{\currentpropheader}{Proposition I.5}
    \addcontentsline{toc}{chapter}{Proposition I.5}
    \textit{In isosceles triangles the angles at the base are equal to one another, and, if the equal straight lines be produced further, the angles under the base will be equal to one another.}
	\vskip 0.5\baselineskip

    % Check line numbering toggle
\iftoggle{lineNumbering}{\linenumbers}{}

% Prop text
\iftoggle{heath}{
    Let \drawLineAlt{A,B,C} be an isosceles triangle having the side \drawUnitLineAlt{AB} equal to the side \drawUnitLineAlt{AC};\par
    and let the straight lines \drawUnitLineAlt{BD},\drawUnitLineAlt{CE} be produced further in a straight line with \drawUnitLineAlt{AB},\drawUnitLineAlt{AC}.  \justification{Post. 2}
    I say that the angle \drawAngleAlt{ABC} is equal to the angle \drawAngleAlt{ACB}, and the angle \drawAngleAlt{CBD} to the angle \drawAngleAlt{BCE}.
    \vskip \baselineskip
    Let a point  \drawPointAlt{F} be taken at random on \drawUnitLineAlt{BD}; from \drawUnitLineAlt{AE} the greater let \drawUnitLineAlt{AG} be cut off equal to \drawUnitLineAlt{AF} the less; \justification{I. 3}
    and let the straight lines \drawUnitLineAlt{FC},\drawUnitLineAlt{GB} be joined. \justification{Post. 1}
    Then, since \drawUnitLineAlt{AF} is equal to \drawUnitLineAlt{AG} and \drawUnitLineAlt{AB} to \drawUnitLineAlt{AC},\par
    {\centering the two sides \drawUnitLineAlt{FA},\drawUnitLineAlt{AC} are equal to the two sides \drawUnitLineAlt{GA},\drawUnitLineAlt{AB}, respectively; \par}
    and they contain a common angle, the angle \drawAngleAlt{FAG}.\par 
    {\centering Therefore the base \drawUnitLineAlt{FC} is equal to the base \drawUnitLineAlt{GB}, \par}
    {\centering and the triangle \drawLineAlt{A,F,C} is equal to the triangle \drawLineAlt{A,G,B}, \par}
    and the remaining angles will be equal to the remaining angles respectively, namely those which the equal sides subtend, \par
    {\centering that is, the angle \drawAngleAlt{ACF} to the angle \drawAngleAlt{ABG}, \par}
    \centeredjust{and the angle \drawAngleAlt{AFC} to the angle \drawAngleAlt{AGB}.}{I. 4}
    And, since the whole \drawUnitLineAlt{AF} is equal to the whole \drawUnitLineAlt{AG}, and in these \drawUnitLineAlt{AB} is equal to \drawUnitLineAlt{AC}, \par
    {\centering the remainder \drawUnitLineAlt{BF} is equal to the remainder \drawUnitLineAlt{CG}. \par}
    But \drawUnitLineAlt{FC} was also proved equal to \drawUnitLineAlt{GB}; \par
    {\centering therefore the two sides \drawUnitLineAlt{BF},\drawUnitLineAlt{FC} are equal to the two sides \drawUnitLineAlt{CG},\drawUnitLineAlt{GB} respectively; \par}
    and the angle \drawAngleAlt{BFC} is equal to the angle \drawAngleAlt{CGB}, \par
    \nolinenumbers
    \clearpage
    \linenumbers
    while the base \drawUnitLineAlt{BC} is common to them; \par
    therefore the triangle \drawLineAlt{B,F,C} is also equal to the triangle \drawLineAlt{C,G,B}, and the remaining angles will be equal to the remaining angles respectively, namely those which the equal sides subtend;\par
    {\centering therefore the angle \drawAngleAlt{FBC} is equal to the angle \drawAngleAlt{GCB}, \par}
    {\centering and the angle \drawAngleAlt{BCF} to the angle \drawAngleAlt{CBG}. \par}
    Accordingly, since the whole angle \drawAngleAlt{ABG} was proved equal to the angle \drawAngleAlt{ACF}, and in these the angle \drawAngleAlt{CBG} is equal to the angle \drawAngleAlt{BCF}, \par
    {\centering the remaining angle \drawAngleAlt{ABC} is equal to the remaining angle \drawAngleAlt{ACB}; \par}
    and they are at the base of the triangle \drawLineAlt{A,B,C}. \par
    But the angle \drawAngleAlt{FBC} was also proved equal to the angle \drawAngleAlt{GCB}; \par
    and they are under the base. \par
    Therefore etc. 
}{}

\iftoggle{raum}{
    Raum_Text
}{}
\nolinenumbers
\hfill Q. E. D.

}

\resetlinenumber

\startProposition{
    \markboth{Euclid in Color}{Proposition I.6}
    \defineNewPicture[1/4][4]{
    textLabels := true;
	pair A, B, C, D;
	B := (0, 0);
	C := B shifted (5/2u, 0);
	A := 1/2[B,C] shifted (0, 3u);
	D := 3/4[B, A];
	byAngleDefine(C, B, A, bygreen, SOLID_SECTOR);
	byAngleDefine(B, C, A, byblack, SOLID_SECTOR);
	draw byNamedAngleResized();
	byLineDefine(D, C, byyellow, SOLID_LINE, REGULAR_WIDTH);
	byLineDefine(B, C, byred, SOLID_LINE, REGULAR_WIDTH);
	byLineDefine(C, A, byblue, SOLID_LINE, REGULAR_WIDTH);
	byLineDefine(D, B, byblack, SOLID_LINE, REGULAR_WIDTH);
	byLineDefine(D, A, byblack, SOLID_LINE, REGULAR_WIDTH);
	draw byNamedLine(DC);
	draw byNamedLineSeq(0)(DB,DA,CA,BC);
	draw byLabelsOnPolygon(B, D, A, C)(ALL_LABELS, 0);
    }
    ~\hfill\drawCurrentPicture\hfill~
}{
    \chapter*{Proposition I.6}
    \renewcommand{\currentpropheader}{Proposition I.6}
    \addcontentsline{toc}{chapter}{Proposition I.6}
    \textit{If in a triangle two angles be equal to one another, the sides which subtend the equal angles will also be equal to one another.}
	\vskip 0.5\baselineskip

    % Check line numbering toggle
\iftoggle{lineNumbering}{\linenumbers}{}

% Prop text
\iftoggle{heath}{
    Let \drawLineAlt{A,B,C} be a triangle having the angle \drawAngleAlt{ABC} equal to the angle \drawAngleAlt{ACB}; 
    I say that the side \drawUnitLineAlt{AB} is also equal to the side \drawUnitLineAlt{AC}. 
    \vskip \baselineskip
    For, if \drawUnitLineAlt{AB} is unequal to \drawUnitLineAlt{AC}, one of them is greater. \par
    Let \drawUnitLineAlt{AB} be greater; \par
    and from \drawUnitLineAlt{AB} the greater let \drawUnitLineAlt{DB} be cut off equal to \drawUnitLineAlt{AC} the less; \par
    let \drawUnitLineAlt{DC} be joined. \par
    Then, since \drawUnitLineAlt{DB} is equal to \drawUnitLineAlt{AC}, \par
    and \drawUnitLineAlt{BC} is common, 
    \par the two sides \drawUnitLineAlt{DB},\drawUnitLineAlt{BC} are equal to the two sides \drawUnitLineAlt{AC},\drawUnitLineAlt{CB} respectively; \par
    and the angle \drawAngleAlt{DBC} is equal to the angle \drawAngleAlt{ACB}; \par
    {\centering therefore the base \drawUnitLineAlt{DC} is equal to the base \drawUnitLineAlt{AB}, \par}
    and the triangle \drawLineAlt{D,B,C} will be equal to the triangle \drawLineAlt{A,C,B}, the less to the greater: which is absurd. \par
    Therefore \drawUnitLineAlt{AB} is not unequal to \drawUnitLineAlt{AC}; \par
    {\centering it is therefore equal to it. \par}
    Therefore etc.
}{}

\iftoggle{raum}{
    Let ABC be a triangle having the angle ABC equal to the angle ACB; I say that the side AB is also equal to the side AC. For, if AB is unequal to AC, one of them is greater. Let AB be greater; and from AB the greater let DB be cut off equal to AC the less; let DC be joined. Then, since DB is equal to AC, and BC is common, the two sides DB, BC are equal to the two sides AC, CB respectively; and the angle DBC is equal to the angle ACB; therefore the base DC is equal to the base AB, and the triangle DBC will be equal to the triangle ACB, the less to the greater: which is absurd. Therefore AB is not unequal to AC; it is therefore equal to it. Therefore etc.
}{}
\nolinenumbers
\hfill Q. E. D.

}

\resetlinenumber

\startProposition{
    \markboth{Euclid in Color}{Proposition I.7}
    \defineNewPicture[1][3]{
    textLabels := true;
	pair A, B, C, D;
	A := (0, 0);
	B := A shifted (4u, 0);
	C := A shifted (11/4u, 3u);
	D := C shifted (u, -1/2u);
	byAngleDefine.C(B, C, A, byblack, SOLID_SECTOR);
	byAngleDefine(D, C, B, byyellow, SOLID_SECTOR);
	byAngleDefine.D(A, D, B, byred, SOLID_SECTOR);
	byAngleDefine(C, D, A, byblue, SOLID_SECTOR);
	draw byNamedAngleResized();
	draw byLine(C, D, byblack, SOLID_LINE, REGULAR_WIDTH);
	byLineDefine(A, B, byblack, SOLID_LINE, REGULAR_WIDTH);
	byLineDefine(B, C, byblue, SOLID_LINE, REGULAR_WIDTH);
	byLineDefine(C, A, byred, SOLID_LINE, REGULAR_WIDTH);
	byLineDefine(B, D, bygreen, SOLID_LINE, REGULAR_WIDTH);
	byLineDefine(D, A, byyellow, SOLID_LINE, REGULAR_WIDTH);
	draw byNamedLineSeq(0)(BC,CA);
	draw byNamedLineSeq(0)(AB,BD,DA);
	draw byLabelsOnPolygon(A, C, D, B)(ALL_LABELS, 0);
    }
    ~\hfill\drawCurrentPicture\hfill~
}{
\chapter*{Proposition I.7}
    \renewcommand{\currentpropheader}{Proposition I.7}
    \addcontentsline{toc}{chapter}{Proposition I.7}
    \textit{Given two straight lines constructed on a straight line (from its extremities) and meeting in a point, there cannot be constructed on the same straight line (from its extremities), and on the same side of it, two other straight lines meeting in another point and equal to the former two respectively, namely each to that which has the same extremity with it.}
	\vskip 0.5\baselineskip

    % Check line numbering toggle
\iftoggle{lineNumbering}{\linenumbers}{}

% Prop text
\iftoggle{heath}{
    For, if possible, given two straight lines \drawUnitLineAlt{AC},\drawUnitLineAlt{CB} constructed on the straight line \drawUnitLineAlt{AB} and meeting at the point \drawPointAlt{C}, let two other straight lines \drawUnitLineAlt{AD},\drawUnitLineAlt{DB} be constructed on the same straight line \drawUnitLineAlt{AB}, on the same side of it, meeting in another point \drawPointAlt{D} and equal to the former two respectively, namely each to that which has the same extremity with it, so that \drawUnitLineAlt{CA} is equal to \drawUnitLineAlt{DA} which has the same extremity \drawPointAlt{A} with it, \par
    and \drawUnitLineAlt{CB} to \drawUnitLineAlt{DB} which has the same extremity \drawPointAlt{B} with it; and let \drawUnitLineAlt{CD} be joined. \par
    Then, since \drawUnitLineAlt{AC} is equal to \drawUnitLineAlt{AD}, \par
    \centeredjust{the angle \drawAngleAlt{ACD} is also equal to the angle \drawAngleAlt{ADC};}{I. 5} 
    {\centering therefore the angle \drawAngleAlt{ADC} is greater than the angle \drawAngleAlt{DCB}; \par}
    {\centering therefore the angle \drawAngleAlt{CDB} is much greater than the angle \drawAngleAlt{DCB}. \par}
    Again, since \drawUnitLineAlt{CB} is equal to \drawUnitLineAlt{DB}, \par
    {\centering the angle \drawAngleAlt{CDB} is also equal to the angle \drawAngleAlt{DCB}. \par}
    But it was also proved much greater than it: \par
    which is impossible. \par
    Therefore etc.
}{}

\iftoggle{raum}{
    Raum text here
}{}
\nolinenumbers 
\hfill Q. E. D.

}

\resetlinenumber

\startProposition{
    \markboth{Euclid in Color}{Proposition I.8}
    \defineNewPicture[1/4][7/4]{
    textLabels := true;
	pair A, B, C, D, E, F, G, d;
	A := (0, 0);
	B := A shifted (-u, -4u);
	C := A shifted (3/2u, -3u);
	d := (4u, 0);
	D := A shifted d;
	E := B shifted d;
	F := C shifted d;
	G := D shifted (3/4u, 0);
	byAngleDefine(F, D, E, byred, SOLID_SECTOR);
	byAngleDefine(C, A, B, byyellow, SOLID_SECTOR);
	byAngleDefine(E, G, F, byblue, SOLID_SECTOR);
	draw byNamedAngleResized();
	byLineDefine(A, B, byred, SOLID_LINE, REGULAR_WIDTH);
	byLineDefine(B, C, byblue, SOLID_LINE, REGULAR_WIDTH);
	byLineDefine(C, A, byblack, SOLID_LINE, REGULAR_WIDTH);
	byLineDefine(D, E, byblack, SOLID_LINE, REGULAR_WIDTH);
	byLineDefine(E, F, bygreen, SOLID_LINE, REGULAR_WIDTH);
	byLineDefine(F, D, byblue, SOLID_LINE, REGULAR_WIDTH);
	byLineDefine(E, G, byyellow, SOLID_LINE, REGULAR_WIDTH);
	byLineDefine(F, G, byred, SOLID_LINE, REGULAR_WIDTH);
	draw byNamedLineSeq(0)(CA,BC,AB);
	draw byNamedLineSeq(0)(FD,DE);
	draw byNamedLineSeq(0)(EG,EF,FG);
	draw byLabelsOnPolygon(C, B, A)(ALL_LABELS, 0);
	draw byLabelsOnPolygon(F, E, D, G)(ALL_LABELS, 0);
    }
    ~\hfill\drawCurrentPicture\hfill~
}{
    \chapter*{Proposition I.8}
    \renewcommand{\currentpropheader}{Proposition I.8}
    \addcontentsline{toc}{chapter}{Proposition I.8}
    \textit{If two triangles have the two sides equal to two sides respectively, and have also the base equal to the base, they will also have the angles equal which are contained by the equal straight lines.}
	\vskip 0.5\baselineskip

    % Check line numbering toggle
\iftoggle{lineNumbering}{\linenumbers}{}

\iftoggle{heath}{
    Let \drawLineAlt{A,B,C}, \drawLineAlt{D,E,F} be two triangles having the two sides \drawUnitLineAlt{AB},\drawUnitLineAlt{AC} equal to the two sides \drawUnitLineAlt{DE},\drawUnitLineAlt{DF} respectively, namely \drawUnitLineAlt{AB} to \drawUnitLineAlt{DE}, and \drawUnitLineAlt{AC} to \drawUnitLineAlt{DF}; and let them have the base \drawUnitLineAlt{BC} equal to the base \drawUnitLineAlt{EF}; \par
    I say that the angle \drawAngleAlt{BAC} is also equal to the angle \drawAngleAlt{EDF}.
    \vskip \baselineskip
    For, if the triangle \drawLineAlt{A,B,C} be applied to the triangle \drawLineAlt{D,E,F}, and if the point \drawPointAlt{B} be placed on the point  \drawPointAlt{E} and the straight line \drawUnitLineAlt{BC} on \drawUnitLineAlt{EF}, \par
    {\centering the point  \drawPointAlt{C} will also coincide with  \drawPointAlt{F}, \par}
    because \drawUnitLineAlt{BC} is equal to \drawUnitLineAlt{EF}. \par
    Then,\drawUnitLineAlt{BC} coinciding with \drawUnitLineAlt{EF},\par
    {\centering  \drawUnitLineAlt{BA},\drawUnitLineAlt{AC} will also coincide with \drawUnitLineAlt{ED},\drawUnitLineAlt{DF}; \par}
    for, if the base \drawUnitLineAlt{BC} coincides with the base \drawUnitLineAlt{EF}, and the sides \drawUnitLineAlt{BA},\drawUnitLineAlt{AC} do not coincide with \drawUnitLineAlt{ED},\drawUnitLineAlt{DF} but fall beside them as \drawUnitLineAlt{EG},\drawUnitLineAlt{GF}, \par
    then, given two straight lines constructed on a straight line (from its extremities) and meeting in a point, there will have been constructed on the same straight line (from its extremities), and on the same side of it, two other straight lines meeting in another point and equal to the former two respectively, namely each to that which has the same extremity with it. \par
    But they cannot be so constructed. \justification{I. 7} 
    Therefore it is not possible that, if the base \drawUnitLineAlt{BC} be applied to the base \drawUnitLineAlt{EF}, the sides \drawUnitLineAlt{BA},\drawUnitLineAlt{AC} should not coincide with \drawUnitLineAlt{ED},\drawUnitLineAlt{DF}; \par
    {\centering they will therefore coincide, \par}
    so that the angle \drawAngleAlt{BAC} will also coincide with the angle \drawAngleAlt{EDF}, \par
    and will be equal to it. \par
    If therefore etc. 
}{}

\iftoggle{raum}{
    Raum_Text
}{}
\nolinenumbers
\hfill Q. E. D.

}

\resetlinenumber

\startProposition{
    \markboth{Euclid in Color}{Proposition I.9}
    \defineNewPicture[1/2][3]{
	textLabels := true;
	pair A, B, C, D, E, F;
	numeric r;
	A := (0, 10/3u);
	D := (-u, 0);
	E := D xscaled -1;
	r := abs(D-E);
	F := (subpath (4,8) of (fullcircle scaled 2r shifted D)) intersectionpoint (fullcircle scaled 2r shifted E);
	B := 3/2[A, D];
	C := 3/2[A, E];
	byAngleDefine(D, A, F, byblue, SOLID_SECTOR);
	byAngleDefine(E, A, F, byyellow, SOLID_SECTOR);
	draw byNamedAngleResized();
	byLineDefine(D, E, byyellow, SOLID_LINE, REGULAR_WIDTH);
	byLineDefine(A, F, byblack, SOLID_LINE, REGULAR_WIDTH);
	byLineDefine(F, D, byblue, SOLID_LINE, REGULAR_WIDTH);
	byLineDefine(E, F, byred, SOLID_LINE, REGULAR_WIDTH);
	byLineDefine(A, D, byred, SOLID_LINE, REGULAR_WIDTH);
	byLineDefine(E, A, bygreen, SOLID_LINE, REGULAR_WIDTH);
	byLineDefine(D, B, byred, SOLID_LINE, REGULAR_WIDTH);
	byLineDefine(E, C, bygreen, SOLID_LINE, REGULAR_WIDTH);
	draw byNamedLine(DE,AF);
	draw byNamedLineSeq(0)(FD,EF);
	draw byNamedLineSeq(0)(DB,AD,EA,EC);
	draw byLabelsOnPolygon(F, D, A, E)(ALL_LABELS, 0);
	draw byLabelPoint(B, 0, -2);
	draw byLabelPoint(C, 0, 2);
    }
    ~\hfill\drawCurrentPicture\hfill~
}{
    \chapter*{Proposition I.9}
    \renewcommand{\currentpropheader}{Proposition I.9}
    \addcontentsline{toc}{chapter}{Proposition I.9}
    \textit{To bisect a given rectilineal angle.}
	\vskip 0.5\baselineskip

    % Check line numbering toggle
\iftoggle{lineNumbering}{\linenumbers}{}

\iftoggle{heath}{
    Let the angle \drawAngleAlt{BAC} be the given rectilineal angle. \par
    Thus it is required to bisect it. 
    \vskip \baselineskip
    Let a point  \drawPointAlt{D} be taken at random on \drawUnitLineAlt{AB}; \par
    let \drawUnitLineAlt{AE} be cut off from \drawUnitLineAlt{AC} equal to \drawUnitLineAlt{AD};\justification{I. 3}
    let \drawUnitLineAlt{DE} be joined, \par
    and on \drawUnitLineAlt{DE} let the equilateral triangle \drawLineAlt{D,E,F} be constructed; \par
    let \drawUnitLineAlt{AF} be joined. \par
    I say that the angle \drawAngleAlt{BAC} has been bisected by the straight line \drawUnitLineAlt{AF}. \par
    For, since \drawUnitLineAlt{AD} is equal to \drawUnitLineAlt{AE}, \par
    and \drawUnitLineAlt{AF} is common, \par
    {\centering the two sides \drawUnitLineAlt{DA},\drawUnitLineAlt{AF} are equal to the two sides \drawUnitLineAlt{EA},\drawUnitLineAlt{AF} respectively. \par}
    And the base \drawUnitLineAlt{DF} is equal to the base \drawUnitLineAlt{EF}; \par
    \centeredjust{therefore the angle \drawAngleAlt{DAF} is equal to the angle \drawAngleAlt{EAF}.}{I. 8} 
    Therefore the given rectilineal angle \drawAngleAlt{BAC} has been bisected by the straight line \drawUnitLineAlt{AF}.
    
}{}

\iftoggle{raum}{
    Raum_Text
}{}
\nolinenumbers
\hfill Q. E. F.

}

\resetlinenumber

\startProposition{
    \markboth{Euclid in Color}{Proposition I.10}
    \defineNewPicture[1/4][3]{
    textLabels := true;
	pair A, B, C, D;
	numeric r;
	A := (-2u, 0);
	B := A xscaled -1;
	r := abs(A-B);
	C := (fullcircle scaled 2r shifted A) intersectionpoint (fullcircle scaled 2r shifted B);
	D := 1/2[A, B];
	byAngleDefine(A, C, D, bygreen, SOLID_SECTOR);
	byAngleDefine(B, C, D, byyellow, SOLID_SECTOR);
	draw byNamedAngleResized();
	draw byLine(C, D, byred, SOLID_LINE, REGULAR_WIDTH);
	byLineDefine(D, A, byblack, SOLID_LINE, REGULAR_WIDTH);
	byLineDefine(B, D, byblack, SOLID_LINE, REGULAR_WIDTH);
	byLineDefine(C, A, byyellow, SOLID_LINE, REGULAR_WIDTH);
	byLineDefine(B, C, byblue, SOLID_LINE, REGULAR_WIDTH);
	draw byNamedLineSeq(0)(CA,BC,BD,DA);
	draw byLabelsOnPolygon(A, C, B, D)(ALL_LABELS, 0);
    }
    ~\hfill\drawCurrentPicture\hfill~
}{
    \chapter*{Proposition I.10}
    \renewcommand{\currentpropheader}{Proposition I.10}
    \addcontentsline{toc}{chapter}{Proposition I.10}
    \textit{To bisect a given finite straight line.}

	\vskip 0.5\baselineskip
    % Check line numbering toggle
	\iftoggle{lineNumbering}{\linenumbers}{}
 
	\iftoggle{heath}{
    Let \drawUnitLineAlt{AB} be the given finite straight line.\par 
    Thus it is required to bisect the finite straight line \drawUnitLineAlt{AB}. 
    \vskip \baselineskip
    Let the equilateral triangle \drawLineAlt{A,B,C} be constructed on it, \justification{I. 1}
    and let the angle \drawAngleAlt{ACB} be bisected by the straight line \drawUnitLineAlt{CD}; \justification{I. 9}
    I say that the straight line \drawUnitLineAlt{AB} has been bisected at the point  \drawPointAlt{D}. \par
    For, since \drawUnitLineAlt{AC} is equal to \drawUnitLineAlt{CB}, and \drawUnitLineAlt{CD} is common, \par
    {\centering the two sides \drawUnitLineAlt{AC},\drawUnitLineAlt{CD} are equal to the two sides \drawUnitLineAlt{BC},\drawUnitLineAlt{CD} respectively; \par}
    and the angle \drawAngleAlt{ACD} is equal to the angle \drawAngleAlt{BCD}; \par
    \centeredjust{therefore the base \drawUnitLineAlt{AD} is equal to the base \drawUnitLineAlt{BD}.}{I. 4}
    Therefore the given finite straight line \drawUnitLineAlt{AB} has been bisected at  \drawPointAlt{D}. 
}{}

\iftoggle{raum}{
	Raum_Text
}{}
\nolinenumbers
\hfill Q. E. F.

}

\resetlinenumber

\resetlinenumber

\startProposition{
    \markboth{Euclid in Color}{Proposition I.11}
    \defineNewPicture[1/2][3]{
	pair A, B, C, D, E, F;
	numeric r;
	A := (-2u, 0);
	B := (2u, 0);
	C := 1/2[A, B];
	D := 2/3[C, A];
	E := 2/3[C, B];
	r := abs(D-E);
	F := (fullcircle scaled 2r shifted D) intersectionpoint (fullcircle scaled 2r shifted E);
	byAngleDefine(F, C, D, byred, SOLID_SECTOR);
	byAngleDefine(F, C, E, byblue, SOLID_SECTOR);
	draw byNamedAngleResized();
	draw byLine(F, C, byyellow, SOLID_LINE, REGULAR_WIDTH);
	byLineDefine(F, D, byblue, SOLID_LINE, REGULAR_WIDTH);
	byLineDefine(F, E, bygreen, SOLID_LINE, REGULAR_WIDTH);
	draw byNamedLineSeq(0)(FD,FE);
	byLineDefine(D, C, byblack, SOLID_LINE, REGULAR_WIDTH);
	byLineDefine(A, D, byblack, SOLID_LINE, REGULAR_WIDTH);
	byLineDefine(C, E, byblack, SOLID_LINE, REGULAR_WIDTH);
	byLineDefine(E, B, byblack, SOLID_LINE, REGULAR_WIDTH);
	draw byNamedLineSeq(0)(DC,AD,CE,EB);
	draw byLabelPoint(A, -90, 1);
	draw byLabelPoint(B, -90, 1);
	draw byLabelPoint(C, -90, 1);
	draw byLabelPoint(D, -90, 1);
	draw byLabelPoint(E, -90, 1);
	draw byLabelsOnPolygon(D, F, E)(OMIT_FIRST_LABEL+OMIT_LAST_LABEL, 0);
    }
    ~\hfill\drawCurrentPicture\hfill~
}{
    \chapter*{Proposition I.11}
    \renewcommand{\currentpropheader}{Proposition I.11}
    \addcontentsline{toc}{chapter}{Proposition 11}
    \textit{To draw a straight line at right angles to a given straight line from a given point on it.}
    \vskip \baselineskip

    % Check line numbering toggle
\iftoggle{lineNumbering}{\linenumbers}{}

\iftoggle{heath}{
    Let \drawUnitLineAlt{AB} be the given straight line, and  \drawPointAlt{C} the given point on it. 
    Thus it is required to draw from the point  \drawPointAlt{C} a straight line at right angles to the straight line \drawUnitLineAlt{AB}. 
    \vskip \baselineskip
    Let a point  \drawPointAlt{D} be taken at random on \drawUnitLineAlt{AC}; \par
    let \drawUnitLineAlt{CE} be made equal to \drawUnitLineAlt{CD}; \justification{I. 3}
    on \drawUnitLineAlt{DE} let the equilateral \drawLineAlt{F,D,E} be constructed, \justification{I. 1}
    and let \drawUnitLineAlt{FC} be joined; \par
    I say that the straight line \drawUnitLineAlt{FC} has been drawn at right angles to the given straight line \drawUnitLineAlt{AB} from  \drawPointAlt{C} the given point on it. \par
    For, since \drawUnitLineAlt{DC} is equal to \drawUnitLineAlt{CE}, and \drawUnitLineAlt{CF} is common, \par
    the two sides \drawUnitLineAlt{DC},\drawUnitLineAlt{CF} are equal to the two sides \drawUnitLineAlt{EC},\drawUnitLineAlt{CF} respectively; and the base \drawUnitLineAlt{DF} is equal to the base \drawUnitLineAlt{FE}; \par
    \centeredjust{therefore the angle \drawAngleAlt{DCF} is equal to the angle \drawAngleAlt{ECF};}{I. 8}
    and they are adjacent angles. \par
    But, when a straight line set up on a straight line makes the adjacent angles equal to one another, each of the equal angles is right; \justification{Def. 10}
    {\centering therefore each of the angles \drawAngleAlt{DCF}, \drawAngleAlt{FCE} is right. \par}
    Therefore the straight line \drawUnitLineAlt{CF} has been drawn at right angles to the given straight line \drawUnitLineAlt{AB} from the given point  \drawPointAlt{C} on it. 
}{}

\iftoggle{raum}{ 
}{}
\nolinenumbers
\hfill Q. E. F.

\resetlinenumber
}

\resetlinenumber

\startProposition{
    \markboth{Euclid in Color}{Proposition I.12}
    \defineNewPicture[1][2]{
    pair A, B, C, D, E, F, G, H;
	path c;
	numeric r, a[];
	C := (0, u);
	D := (sqrt(7)/2*u, -1/3u);
	%c := fullcircle scaled 2r shifted C;
	E := (7/4u, 0);
	F := C + (0, abs(C-E));
	G := E xscaled -1;
	H := 1/2[G, E];
	A := 5/3[H, G];
	B := 5/3[H, E];
	r := abs(C-G);
	byAngleDefine(C, H, G, byyellow, SOLID_SECTOR);
	byAngleDefine(E, H, C, byblue, SOLID_SECTOR);
	draw byNamedAngleResized();
	draw byLine(C, H, byred, SOLID_LINE, REGULAR_WIDTH);
	byLineDefine(C, G, byblue, SOLID_LINE, REGULAR_WIDTH);
	byLineDefine(E, C, bygreen, SOLID_LINE, REGULAR_WIDTH);
	draw byNamedLineSeq(0)(CG,EC);
	draw byCircleR.EFG(C, r, byred, 0, 0, 0);
	%draw byCircleABC.EFG(E, F, G, byred, 0, 0, 0);
	byLineDefine(H, G, byblack, SOLID_LINE, REGULAR_WIDTH);
	byLineDefine(G, A, byblack, SOLID_LINE, REGULAR_WIDTH);
	byLineDefine(E, H, byblack, SOLID_LINE, REGULAR_WIDTH);
	byLineDefine(B, E, byblack, SOLID_LINE, REGULAR_WIDTH);
	byLineDefine(C, D, byblack, SOLID_LINE, REGULAR_WIDTH);
	draw byNamedLineSeq(1)(GA, HG, EH, BE);
	draw byLabelsOnPolygon(G, C, E)(OMIT_FIRST_LABEL+OMIT_LAST_LABEL, 0);
	draw byLabelPoint(G, 240, 1.2);
	draw byLabelPoint(H, -90, 1);
	draw byLabelPoint(E, 300, 1.2);
	draw byLabelPoint(A, -90, 1);
	draw byLabelPoint(B, -90, 1);
	draw byLabelPoint(D, -90, 5);
	draw byLabelPoint(F, 90, 1);
    }
    ~\hfill\drawCurrentPicture\hfill~
}{
    \chapter*{Proposition I.12}
    \renewcommand{\currentpropheader}{Proposition I.12}
    \addcontentsline{toc}{chapter}{Proposition 12}
    \textit{To a given infinite straight line, from a given point which is not on it, to draw a perpendicular straight line.}
    \vskip \baselineskip

    % Check line numbering toggle
\iftoggle{lineNumbering}{\linenumbers}{}

\iftoggle{heath}{
    Let \drawUnitLineAlt{AB} be the given infinite straight line, and  \drawPointAlt{C} the given point which is not on it; \par
    thus it is required to draw to the given infinite straight line \drawUnitLineAlt{AB}, from the given point  \drawPointAlt{C} which is not on it, a perpendicular straight line. 
    \vskip \baselineskip
    For let a point  \drawPointAlt{D} be taken at random on the other side of the straight line \drawUnitLineAlt{AB}, and with center  \drawPointAlt{C} and distance  \drawPointAlt{CD} let the circle \drawCircleAlt[byred]{EFG} be described;\par 
    \justification{Post. 3}
    let the straight line \drawUnitLineAlt{EG} be bisected at  \drawPointAlt{H}, \justification{I. 10}
    and let the straight lines \drawUnitLineAlt{CG},\drawUnitLineAlt{CH},\drawUnitLineAlt{CE} be joined. \justification{Post. 1}
    I say that \drawUnitLineAlt{CH} has been drawn perpendicular to the given infinite straight line \drawUnitLineAlt{AB} from the given point  \drawPointAlt{C} which is not on it. \par
    For, since \drawUnitLineAlt{GH} is equal to \drawUnitLineAlt{HE}, and \drawUnitLineAlt{HC} is common, \par
    {\centering the two sides \drawUnitLineAlt{GH},\drawUnitLineAlt{HC} are equal to the two sides \drawUnitLineAlt{EH},\drawUnitLineAlt{HC} respectively; \par}
    and the base \drawUnitLineAlt{CG} is equal to the base \drawUnitLineAlt{CE}; \par
    \centeredjust{therefore the angle \drawAngleAlt{CHG} is equal to the angle \drawAngleAlt{EHC}.}{I. 8}
    And they are adjacent angles. \par
    But, when a straight line set up on a straight line makes the adjacent angles equal to one another, each of the equal angles is right, and the straight line standing on the other is called a perpendicular to that on which it stands. \justification{Def. 10}
    Therefore \drawUnitLineAlt{CH} has been drawn perpendicular to the given infinite straight line \drawUnitLineAlt{AB} from the given point  \drawPointAlt{C} which is not on it. 
}{}

\iftoggle{raum}{
    Raum_Text
}{}
\nolinenumbers
\hfill Q. E. F.

\resetlinenumber
}

\resetlinenumber

\startProposition{
    \markboth{Euclid in Color}{Proposition I.13}
    \defineNewPicture[1][2]{
	pair A, B, C, D, E;
	E := (0, 5/2u);
	D := (-3u, 0);
	C := D xscaled -1;
	B := (xpart(E), ypart(D));
	A := (1/3xpart(C), 9/10ypart(E));
	byAngleDefine(E, B, D, byyellow, SOLID_SECTOR);
	byAngleDefine(A, B, E, byred, SOLID_SECTOR);
	byAngleDefine(C, B, A, byblue, SOLID_SECTOR);
	draw byNamedAngleResized();
	draw byLine(E, B, byblack, SOLID_LINE, REGULAR_WIDTH);
	draw byLine(A, B, byyellow, SOLID_LINE, REGULAR_WIDTH);
	draw byLine(D, C, byred, SOLID_LINE, REGULAR_WIDTH);
	draw byLabelsOnPolygon(C, B, D, noPoint)(ALL_LABELS, 0);
	draw byLabelLineEnd(A, B, 0);
	draw byLabelLineEnd(E, B, 0);
    }
    ~\hfill\drawCurrentPicture\hfill~
}{
    \chapter*{Proposition I.13}
    \renewcommand{\currentpropheader}{Proposition I.13}
    \addcontentsline{toc}{chapter}{Proposition 13}
    \textit{If a straight line set up on a straight line make angles, it will make either two right angles or angles equal to two right angles.}
    \vskip \baselineskip

    % Check line numbering toggle
\iftoggle{lineNumbering}{\linenumbers}{}

\iftoggle{heath}{
    For let any straight line \drawUnitLineAlt{AB} set up on the straight line \drawUnitLineAlt{CD} make the angles \drawAngleAlt{CBA}, \drawAngleAlt{ABD}; \par
    I say that the angles \drawAngleAlt{CBA}, \drawAngleAlt{ABD} are either two right angles or equal to two right angles. 
    \vskip \baselineskip
    Now, if the angle \drawAngleAlt{CBA} is equal to the angle \drawAngleAlt{ABD}, they are two right angles.\justification{Def. 10}
    But, if not, let \drawUnitLineAlt{BE} be drawn from the point \drawPointAlt{B} at right angles to \drawUnitLineAlt{CD}; \justification{I. 11}
    {\centering therefore the angles \drawAngleAlt{CBE}, \drawAngleAlt{EBD} are two right angles. \par}
    Then, since the angle \drawAngleAlt{CBE} is equal to the two angles \drawAngleAlt{CBA}, \drawAngleAlt{ABE}, \par
    let the angle \drawAngleAlt{EBD} be added to each;\par 
    {\centering therefore the angles \drawAngleAlt{CBE}, \drawAngleAlt{EBD} are equal to the three angles \drawAngleAlt{CBA}, \drawAngleAlt{ABE}, \drawAngleAlt{EBD}.} 
    \hspace*{1em} \justification{C. N. 2}
    Again, since the angle \drawAngleAlt{DBA} is equal to the two angles \drawAngleAlt{DBE}, \drawAngleAlt{EBA}, \par
    let the angle \drawAngleAlt{ABC} be added to each; \par
    \centeredjust{therefore the angles \drawAngleAlt{DBA}, \drawAngleAlt{ABC} are equal to the three angles \drawAngleAlt{DBE}, \drawAngleAlt{EBA}, \drawAngleAlt{ABC}.}{C. N. 2}
    But the angles \drawAngleAlt{CBE}, \drawAngleAlt{EBD} were also proved equal to the same three angles; \par
    \centeredjust{and things which are equal to the same thing are also equal to one another;}{C. N. 1} 
    {\centering therefore the angles \drawAngleAlt{CBE}, \drawAngleAlt{EBD} are also equal to the angles \\ \drawAngleAlt{DBA}, \drawAngleAlt{ABC}. \par}
    But the angles \drawAngleAlt{CBE}, \drawAngleAlt{EBD} are two right angles; \par
    {\centering therefore the angles \drawAngleAlt{DBA}, \drawAngleAlt{ABC} are also equal to two right angles. \par}
    Therefore etc. 
}{}

\iftoggle{raum}{
    Raum_Text
}{}
\nolinenumbers
\hfill Q. E. D.

\resetlinenumber
}

\resetlinenumber

\startProposition{
    \markboth{Euclid in Color}{Proposition I.14}
    \defineNewPicture[1/4][5/2]{
	pair A, B, C, D, E, F;
	A := (-u, 3/2u);
	C := (-5/2u, 0);
	D := C xscaled -1;
	B := (0, 0);
	E := (9/10xpart(D), ypart(A));
	byAngleDefine(A, B, C, byyellow, SOLID_SECTOR);
	byAngleDefine(A, B, E, bygreen, SOLID_SECTOR);
	byAngleDefine(E, B, D, byred, SOLID_SECTOR);
	draw byNamedAngleResized();
	draw byLine(A, B, byred, SOLID_LINE, REGULAR_WIDTH);
	draw byLine(E, B, byyellow, SOLID_LINE, REGULAR_WIDTH);
	draw byLine(C, B, byblue, SOLID_LINE, REGULAR_WIDTH);
	draw byLine(D, B, byblack, SOLID_LINE, REGULAR_WIDTH);
	draw byLabelsOnPolygon(E, B, C, noPoint)(ALL_LABELS, 0);
	draw byLabelsOnPolygon(D, C, noPoint)(OMIT_LAST_LABEL, 0);
	draw byLabelLineEnd(A, B, 0);
    }
    ~\hfill\drawCurrentPicture\hfill~
}{
    \chapter*{Proposition I.14}
    \renewcommand{\currentpropheader}{Proposition I.14}
    \addcontentsline{toc}{chapter}{Proposition 14}
    \textit{If with any straight line, and at a point on it, two straight lines not lying on the same side make the adjacent angles equal to two right angles, the two straight lines will be in a straight line with one another.}
    \vskip \baselineskip

    % Check line numbering toggle
\iftoggle{lineNumbering}{\linenumbers}{}

\iftoggle{heath}{
    For with any straight line \drawUnitLineAlt{AB}, and at the point \drawPointAlt{B} on it, let the two straight lines \drawUnitLineAlt{BC},\drawUnitLineAlt{BD} not lying on the same side make the adjacent angles \drawAngleAlt{ABC}, \drawAngleAlt{ABD} equal to two right angles; \par
    I say that \drawUnitLineAlt{BD} is in a straight line with \drawUnitLineAlt{CB}. 
    \vskip \baselineskip
    For, if \drawUnitLineAlt{BD} is not in a straight line with \drawUnitLineAlt{BC}, let \drawUnitLineAlt{BE} be in a straight line with \drawUnitLineAlt{CB}. \par
    Then, since the straight line \drawUnitLineAlt{AB} stands on the straight line \drawAngleAlt{CBE}, \par
    \centeredjust{the angles \drawAngleAlt{ABC}, \drawAngleAlt{ABE} are equal to two right angles.}{I. 13}
    But the angles \drawAngleAlt{ABC}, \drawAngleAlt{ABD} are also equal to two right angles; \par
    {\centering therefore the angles \drawAngleAlt{CBA}, \drawAngleAlt{ABE} are equal to the angles \par}
    \centeredjust{\drawAngleAlt{CBA}, \drawAngleAlt{ABD}.}{Post. 4 and C. N. 1}
    Let the angle \drawAngleAlt{CBA} be subtracted from each;\par
    {\centering therefore the remaining angle \drawAngleAlt{ABE} is equal to the remaining angle \par}
    \centeredjust{\drawAngleAlt{ABD},}{C. N. 3}
    the less to the greater: which is impossible. \par
    Therefore \drawUnitLineAlt{BE} is not in a straight line with \drawUnitLineAlt{CB}. \par
    Similarly we can prove that neither is any other straight line except \drawUnitLineAlt{BD}. \par
   {\centering Therefore \drawUnitLineAlt{CB} is in a straight line with \drawUnitLineAlt{BD}. \par}
    Therefore etc. 
}{}

\iftoggle{raum}{
    Raum_Text
}{}
\nolinenumbers
\hfill Q. E. D.

\resetlinenumber
}

\resetlinenumber

\startProposition{
    \markboth{Euclid in Color}{Proposition I.15}
    \defineNewPicture[1][3]{
    pair A, B, C, D, E;
	A := (-2u, u);
	B := A scaled -1;
	C := (2u, 0);
	D := C xscaled -1;
	E := (A--B) intersectionpoint (C--D);
	byAngleDefine(B, E, C, byred, SOLID_SECTOR);
	byAngleDefine(C, E, A, byyellow, SOLID_SECTOR);
	byAngleDefine(A, E, D, byblue, SOLID_SECTOR);
	byAngleDefine(D, E, B, bygreen, SOLID_SECTOR);
	draw byNamedAngleResized();
	draw byLine(A, E, byred, SOLID_LINE, REGULAR_WIDTH);
	draw byLine(E, B, byred, SOLID_LINE, REGULAR_WIDTH);
	draw byLine(C, E, byblack, SOLID_LINE, REGULAR_WIDTH);
	draw byLine(E, D, byblack, SOLID_LINE, REGULAR_WIDTH);
	draw byLabelPoint(A, 90, 1);
	draw byLabelPoint(B, -90, 1);
	draw byLabelPoint(C, -90, 1);
	draw byLabelPoint(D, -90, 1);
	draw byLabelPoint(E, 60, 1);
    }
    ~\hfill\drawCurrentPicture\hfill~
}{
    \chapter*{Proposition I.15}
    \renewcommand{\currentpropheader}{Proposition I.15}
    \addcontentsline{toc}{chapter}{Proposition 15}
    \textit{If two straight lines cut one another, they make the vertical angles equal to one another.}
    \vskip \baselineskip

    % Check line numbering toggle
\iftoggle{lineNumbering}{\linenumbers}{}

\iftoggle{heath}{
    For let the straight lines \drawUnitLineAlt{AB},\drawUnitLineAlt{CD} cut one another at the point  \drawPointAlt{E}; \par
    I say that the angle \drawAngleAlt{AEC} is equal to the angle \drawAngleAlt{DEB}, and the angle \drawAngleAlt{CEB} to the angle \drawAngleAlt{AED}. 
    \vskip \baselineskip
    For, since the straight line \drawUnitLineAlt{AE} stands on the straight line \drawUnitLineAlt{CD}, making the angles \drawAngleAlt{CEA}, \drawAngleAlt{AED}, \par
    \centeredjust{the angles \drawAngleAlt{CEA}, \drawAngleAlt{AED} are equal to two right angles.}{I. 13}
    Again, since the straight line \drawUnitLineAlt{DE} stands on the straight line \drawUnitLineAlt{AB}, making the angles \drawAngleAlt{AED}, \drawAngleAlt{DEB}, \par
    \centeredjust{the angles \drawAngleAlt{AED}, \drawAngleAlt{DEB} are equal to two right angles.}{I. 13}
    But the angles \drawAngleAlt{CEA}, \drawAngleAlt{AED} were also proved equal to two right angles; \par
    {\centering therefore the angles \drawAngleAlt{CEA}, \drawAngleAlt{AED} are equal to the angles\par} 
    \centeredjust{\drawAngleAlt{AED} \drawAngleAlt{DEB}.}{Post. 4 and C. N. 1}
    Let the angle \drawAngleAlt{AED} be subtracted from each;\par
    {\centering therefore the remaining angle \drawAngleAlt{CEA} is equal to the remaining angle \par}
    \centeredjust{\drawAngleAlt{BED}.}{C. N. 3}
    Similarly it can be proved that the angles \drawAngleAlt{CEB}, \drawAngleAlt{DEA} are also equal.\par
    Therefore etc. 
}{}

\iftoggle{raum}{
    Raum_Text
}{}
\nolinenumbers
\hfill Q. E. D. 
\vskip \baselineskip
$\langle$\textbf{PORISM.} From this it is manifest that, if two straight lines cut one another, they will make the angles at the point of section equal to four right angles.$\rangle$

\resetlinenumber
}

\resetlinenumber

\startProposition{
    \markboth{Euclid in Color}{Proposition I.16}
    \defineNewPicture[1/4][9/4]{
    pair A, B, C, F, E, D, G;
	A := (0, 0);
	B := A shifted (-3/2u, -3u);
	C := B shifted (3u, 0);
	D := 5/3[B, C];
	E := 1/2[A, C];
	F := 2[B, E];
	G := 5/3[A, C];
	byAngleDefine(A, B, C, byblue, SOLID_SECTOR);
	byAngleDefine(C, A, B, byblack, SOLID_SECTOR);
	byAngleDefine(B, E, A, byyellow, SOLID_SECTOR);
	byAngleDefine(F, E, C, bygreen, SOLID_SECTOR);
	byAngleDefine(E, C, F, byblack, SOLID_SECTOR);
	byAngleDefine(G, C, B, byred, SOLID_SECTOR);
	byAngleDefine(F, C, D, bygreen, SOLID_SECTOR);
	draw byNamedAngleResized();
	byLineDefine(C, D, byblack, SOLID_LINE, REGULAR_WIDTH);
	byLineDefine(C, G, byblue, SOLID_LINE, REGULAR_WIDTH);
	byLineDefine(A, E, byblue, SOLID_LINE, REGULAR_WIDTH);
	byLineDefine(E, C, byblue, SOLID_LINE, REGULAR_WIDTH);
	byLineDefine(B, E, byred, SOLID_LINE, REGULAR_WIDTH);
	byLineDefine(E, F, byred, SOLID_LINE, REGULAR_WIDTH);
	byLineDefine(B, A, bygreen, SOLID_LINE, REGULAR_WIDTH);
	byLineDefine(B, C, byblack, SOLID_LINE, REGULAR_WIDTH);
	byLineDefine(C, F, byyellow, SOLID_LINE, REGULAR_WIDTH);
	draw byNamedLineSeq(0)(BE,EF,CF);
	draw byNamedLineSeq(0)(EC,CG,noLine,CD,BC,BA,AE);
	draw byLabelsOnPolygon(D, B, A, E, F, C)(OMIT_FIRST_LABEL+OMIT_LAST_LABEL, 0);
	draw byLabelsOnPolygon(D, C, G, noPoint)(ALL_LABELS, 0);
    }
    ~\hfill\drawCurrentPicture\hfill~
}{
    \chapter*{Proposition I.16}
    \renewcommand{\currentpropheader}{Proposition I.16}
    \addcontentsline{toc}{chapter}{Proposition 16}
    \textit{In any triangle, if one of the sides be produced, the exterior angle is greater than either of the interior and opposite angles.}
    \vskip \baselineskip

    % Check line numbering toggle
\iftoggle{lineNumbering}{\linenumbers}{}

\iftoggle{heath}{
    Let \drawLineAlt{A,B,C} be a triangle, and let one side of it \drawUnitLineAlt{BC} be produced to  \drawPointAlt{D}; \par
    I say that the exterior angle \drawAngleAlt{ACD} is greater than either of the interior and opposite angles \drawAngleAlt{CBA}, \drawAngleAlt{BAC}. 
    \vskip \baselineskip
    Let \drawUnitLineAlt{AC} be bisected at  \drawPointAlt{E}, \justification{I. 10}
    and let \drawUnitLineAlt{BE} be joined and produced in a straight line to  \drawPointAlt{F}; \par
    let \drawUnitLineAlt{EF} be made equal to \drawUnitLineAlt{BE}, \justification{I. 3}
    let \drawUnitLineAlt{FC} be joined, \justification{Post. 1}
    and let \drawUnitLineAlt{AC} be drawn through to  \drawPointAlt{G}. \justification{Post. 2}
    Then, since \drawUnitLineAlt{AE} is equal to \drawUnitLineAlt{EC}, and \drawUnitLineAlt{BE} to \drawUnitLineAlt{EF}, the two sides \drawUnitLineAlt{AE},\drawUnitLineAlt{EB} are equal to the two sides \drawUnitLineAlt{CE},\drawUnitLineAlt{EF} respectively; \par
    and the angle \drawAngleAlt{AEB} is equal to the angle \drawAngleAlt{FEC}, for they are vertical angles. \par\justification{I. 15}
    Therefore the base \drawUnitLineAlt{AB} is equal to the base \drawUnitLineAlt{FC}, and the triangle \drawLineAlt{A,B,E} is equal to the triangle \drawLineAlt{C,F,E}, \par
    and the remaining angles are equal to the remaining angles respectively, namely those which the equal sides subtend; 
    \justification{I. 4}
    {\centering therefore the angle \drawAngleAlt{BAE} is equal to the angle \drawAngleAlt{ECF}. \par}
    But the angle \drawAngleAlt{ECD} is greater than the angle \drawAngleAlt{ECF}; \justification{C. N. 5}
    therefore the angle \drawAngleAlt{ACD} is greater than the angle \drawAngleAlt{BAE}. \par
    Similarly also, if \drawUnitLineAlt{BC} be bisected, the angle \drawAngleAlt{BCG}, that is, the angle \drawAngleAlt{ACD} [\textit{I. 15}], can be proved greater than the angle \drawAngleAlt{ABC} as well. \par
    Therefore etc. 
}{}

\iftoggle{raum}{
    Raum_Text
}{}
\nolinenumbers
\hfill Q. E. D.

\vskip \baselineskip


\resetlinenumber
}

\resetlinenumber

\startProposition{
    \markboth{Euclid in Color}{Proposition I.17}
    \defineNewPicture[1/3][2]{
    pair A, B, C, D;
	B := (0, 0);
	A := B shifted (-u, 5/2u);
	C := B shifted (5/2u, 0);
	D := C shifted (5/2u, 0);
	byAngleDefine(A, B, C, byblue, SOLID_SECTOR);
	byAngleDefine(B, A, C, byblack, SOLID_SECTOR);
	byAngleDefine(B, C, A, byred, SOLID_SECTOR);
	byAngleDefine(A, C, D, byyellow, SOLID_SECTOR);
	draw byNamedAngleResized();
	byLineDefine(B, A, byred, SOLID_LINE, REGULAR_WIDTH);
	byLineDefine(A, C, byblue, SOLID_LINE, REGULAR_WIDTH);
	byLineDefine(B, C, byblack, SOLID_LINE, REGULAR_WIDTH);
	byLineDefine(C, D, byblack, SOLID_LINE, REGULAR_WIDTH);
	draw byNamedLineSeq(0)(noLine,AC,BA,BC,CD);
	draw byLabelsOnPolygon(D, C, B, A)(ALL_LABELS, 0);
    }
    ~\hfill\drawCurrentPicture\hfill~
}{
    \chapter*{Proposition I.17}
    \renewcommand{\currentpropheader}{Proposition I.17}
    \addcontentsline{toc}{chapter}{Proposition 17}
    \textit{In any triangle two angles taken together in any manner are less than two right angles.}
    \vskip \baselineskip

    % Check line numbering toggle
\iftoggle{lineNumbering}{\linenumbers}{}

\iftoggle{heath}{
    Let \drawLineAlt{A,B,C} be a triangle;\par
    I say that two angles of the triangle \drawLineAlt{A,B,C} taken together in any manner are less than two right angles. 
    \vskip \baselineskip
    For let \drawUnitLineAlt{BC} be produced to  \drawPointAlt{D}. \justification{Post. 2}
    Then, since the angle \drawAngleAlt{ACD} is an exterior angle of the triangle \drawLineAlt{A,B,C}, it is greater than the interior and opposite angle \drawAngleAlt{ABC}. \justification{I. 16}
    Let the angle \drawAngleAlt{ACB} be added to each; \par
    {\centering therefore the angles \drawAngleAlt{ACD}, \drawAngleAlt{ACB} are greater than the angles \drawAngleAlt{ABC}, \drawAngleAlt{BCA}.\par} 
    But the angles \drawAngleAlt{ACD}, \drawAngleAlt{ACB} are equal to two right angles. \justification{I. 13}
    {\centering Therefore the angles \drawAngleAlt{ABC}, \drawAngleAlt{BCA} are less than two right angles.\par}
    Similarly we can prove that the angles \drawAngleAlt{BAC}, \drawAngleAlt{ACB} are also less than two right angles, and so are the angles \drawAngleAlt{CAB}, \drawAngleAlt{ABC} as well. \par
    Therefore etc. 
}{}

\iftoggle{raum}{
    Raum_Text
}{}
\nolinenumbers
\hfill Q. E. D.

\resetlinenumber
}

\resetlinenumber

\startProposition{
    \markboth{Euclid in Color}{Proposition I.18}
    \defineNewPicture[1/3][5/2]{
    pair A, B, C, D;
	numeric a;
	A := (0, 0);
	B := A shifted (u, -2u);
	C := B shifted (7/2u, 0);
	D := A shifted (unitvector(C-A) scaled abs(A-B));
	a := angle(B-D);
	%forsuffixes i=A, B, C, D:
	%i := i rotated -a;
	%endfor;
	byAngleDefine(A, D, B, byblue, SOLID_SECTOR);
	byAngleDefine(D, B, C, byblack, SOLID_SECTOR);
	byAngleDefine(A, B, D, byred, SOLID_SECTOR);
	byAngleDefine(B, C, A, byyellow, SOLID_SECTOR);
	draw byNamedAngleResized();
	draw byLine(D, B, byyellow, SOLID_LINE, REGULAR_WIDTH);
	byLineDefine(D, C, byred, SOLID_LINE, REGULAR_WIDTH);
	byLineDefine(B, C, byblue, SOLID_LINE, REGULAR_WIDTH);
	byLineDefine(B, A, byblack, SOLID_LINE, REGULAR_WIDTH);
	byLineDefine(A, D, bygreen, SOLID_LINE, REGULAR_WIDTH);
	draw byNamedLineSeq(0)(DC,BC,BA,AD);
	draw byLabelsOnPolygon(D, C, B, A)(ALL_LABELS, 0);
    }
    ~\hfill\drawCurrentPicture\hfill~
}{
    \chapter*{Proposition I.18}
    \renewcommand{\currentpropheader}{Proposition I.18}
    \addcontentsline{toc}{chapter}{Proposition 18}
    \textit{In any triangle the greater side subtends the greater angle.}
    \vskip \baselineskip

    % Check line numbering toggle
\iftoggle{lineNumbering}{\linenumbers}{}

\iftoggle{heath}{
    For let \drawLineAlt{A,B,C} be a triangle having the side \drawUnitLineAlt{AC} greater than \drawUnitLineAlt{AB}; \par
    I say that the angle \drawAngleAlt{ABC} is also greater than the angle \drawAngleAlt{BCA}. 
    \vskip \baselineskip
    For, since \drawUnitLineAlt{AC} is greater than \drawUnitLineAlt{AB}, let \drawUnitLineAlt{AD} be made equal to \drawUnitLineAlt{AB} \textit{[I. 3]}, and let \drawUnitLineAlt{BD} be joined. \par
    Then, since the angle \drawAngleAlt{ADB} is an outside angle of the triangle \drawLineAlt{B,C,D}, \par
    \centeredjust{it is greater than the inside and opposite angle \drawAngleAlt{DCB}.}{I. 16}
    But the angle \drawAngleAlt{ADB} is equal to the angle \drawAngleAlt{ABD}, since the side \drawUnitLineAlt{AB} is equal to \drawUnitLineAlt{AD}; \par
    {\centering therefore the angle \drawAngleAlt{ABD} is also greater than the angle \drawAngleAlt{ACB};\par}
    {\centering therefore the angle \drawAngleAlt{ABC} is much greater than the angle \drawAngleAlt{ACB}.\par}
    Therefore etc. 
}{}

\iftoggle{raum}{
    Raum_Text
}{}
\nolinenumbers
\hfill Q. E. D.

\resetlinenumber
}

\resetlinenumber

\startProposition{
    \markboth{Euclid in Color}{Proposition I.19}
    \defineNewPicture[1/3][4]{
	pair A, B, C;
	A := (0, 0);
	B := A shifted (-3/2u, -u);
	C := A shifted (0, -3u);
	byAngleDefine(C, A, B, byblue, SOLID_SECTOR);
	byAngleDefine(A, B, C, byred, SOLID_SECTOR);
	byAngleDefine(B, C, A, byyellow, SOLID_SECTOR);
	draw byNamedAngleResized();
	byLineDefine(A, B, byyellow, SOLID_LINE, REGULAR_WIDTH);
	byLineDefine(B, C, byblue, SOLID_LINE, REGULAR_WIDTH);
	byLineDefine(C, A, byred, SOLID_LINE, REGULAR_WIDTH);
	draw byNamedLineSeq(0)(CA,BC,AB);
	draw byLabelsOnPolygon(B, A, C)(ALL_LABELS, 0);
    }
    ~\hfill\drawCurrentPicture\hfill~
}{
    \chapter*{Proposition I.19}
    \renewcommand{\currentpropheader}{Proposition I.19}
    \addcontentsline{toc}{chapter}{Proposition 19}
    \textit{In any triangle the greater angle is subtended by the greater side.}
    \vskip \baselineskip

    % Check line numbering toggle
\iftoggle{lineNumbering}{\linenumbers}{}

\iftoggle{heath}{
    Let \drawLineAlt{A,B,C} be a triangle having the angle \drawAngleAlt{ABC} greater than the angle \drawAngleAlt{BCA}; \par
    I say that the side \drawUnitLineAlt{AC} is also greater than the side \drawUnitLineAlt{AB}. 
    \vskip \baselineskip
    For, if not,\drawUnitLineAlt{AC} is either equal to \drawUnitLineAlt{AB} or less. \par
    Now \drawUnitLineAlt{AC} is not equal to \drawUnitLineAlt{AB};\par
    for then the angle \drawAngleAlt{ABC} would also have been equal to the angle \drawAngleAlt{ACB}; \par\justification{I. 5}
    but it is not; \par
    {\centering therefore \drawUnitLineAlt{AC} is not equal to \drawUnitLineAlt{AB}. \par}
    Neither is \drawUnitLineAlt{AC} less than \drawUnitLineAlt{AB}, for then the angle \drawAngleAlt{ABC} would also have been less than the angle \drawAngleAlt{ACB}; \justification{I. 18}
    but it is not; \par
    {\centering therefore \drawUnitLineAlt{AC} is not less than \drawUnitLineAlt{AB}.\par}
    And it was proved that it is not equal either.\par
    {\centering Therefore \drawUnitLineAlt{AC} is greater than \drawUnitLineAlt{AB}.\par}
    Therefore etc. 
}{}

\iftoggle{raum}{
    Raum_Text
}{}
\nolinenumbers
\hfill Q. E. D.

\resetlinenumber
}

\resetlinenumber

\startProposition{
    \markboth{Euclid in Color}{Proposition I.20}
    \defineNewPicture[1/3][3]{
	pair B, C, D, A;
	B := (0, 0);
	C := B shifted (7/2u, 0);
	A := B shifted (4/3u, 3/2u);
	D := ((fullcircle scaled 2arclength(A--C)) shifted A) intersectionpoint (A--10[B, A]);
	byAngleDefine(C, D, B, byred, SOLID_SECTOR);
	byAngleDefine(D, C, A, byblue, SOLID_SECTOR);
	byAngleDefine(A, C, B, byyellow, SOLID_SECTOR);
	draw byNamedAngleResized();
	byLineDefine(C, A, byred, SOLID_LINE, REGULAR_WIDTH);
	byLineDefine(B, C, byblack, SOLID_LINE, REGULAR_WIDTH);
	byLineDefine(C, D, byyellow, SOLID_LINE, REGULAR_WIDTH);
	byLineDefine(D, A, bygreen, SOLID_LINE, REGULAR_WIDTH);
	byLineDefine(A, B, byblue, SOLID_LINE, REGULAR_WIDTH);
	draw byNamedLineSeq(0)(CA);
	draw byNamedLineSeq(0)(AB,DA,CD,BC);
	draw byLabelsOnPolygon(A, D, C, B)(ALL_LABELS, 0);
    }
    ~\hfill\drawCurrentPicture\hfill~
}{
    \chapter*{Proposition I.20}
    \renewcommand{\currentpropheader}{Proposition I.20}
    \addcontentsline{toc}{chapter}{Proposition 20}
    \textit{In any triangle two sides taken together in any manner are greater than the remaining one.}
    \vskip \baselineskip

    % Check line numbering toggle
\iftoggle{lineNumbering}{\linenumbers}{}

\iftoggle{heath}{
    For let \drawLineAlt{A,B,C} be a triangle; \par
    I say that in the triangle \drawLineAlt{A,B,C} two sides taken together in any manner are greater than the remaining one, namely\par
    {\centering  \drawUnitLineAlt{BA},\drawUnitLineAlt{AC} greater than \drawUnitLineAlt{BC}; \par}
    {\centering  \drawUnitLineAlt{AB},\drawUnitLineAlt{BC} greater than \drawUnitLineAlt{AC}; \par}
    {\centering  \drawUnitLineAlt{BC},\drawUnitLineAlt{CA} greater than \drawUnitLineAlt{AB}. \par}
    \vskip \baselineskip
    For let \drawUnitLineAlt{BA} be drawn through to the point  \drawPointAlt{D}, let \drawUnitLineAlt{DA} be made equal to \drawUnitLineAlt{CA}, and let \drawUnitLineAlt{DC} be joined.\par
    Then, since \drawUnitLineAlt{DA} is equal to \drawUnitLineAlt{AC}, \par
    \centeredjust{the angle \drawAngleAlt{ADC} is also equal to the angle \drawAngleAlt{ACD};}{I. 5}
    \centeredjust{therefore the angle \drawAngleAlt{BCD} is greater than the angle \drawAngleAlt{ADC}.\hspace{1.5em}}{C. N. 5}
    And, since \drawLineAlt{D,C,B} is a triangle having the angle \drawAngleAlt{BCD} greater than the angle \drawAngleAlt{BDC}, and the greater angle is subtended by the greater side, \justification{I. 19}
    {\centering therefore \drawUnitLineAlt{DB} is greater than \drawUnitLineAlt{BC}. \par}
    But \drawUnitLineAlt{DA} is equal to \drawUnitLineAlt{AC}; \par
    {\centering therefore \drawUnitLineAlt{BA},\drawUnitLineAlt{AC} are greater than \drawUnitLineAlt{BC}. \par}
    Similarly we can prove that \drawUnitLineAlt{AB},\drawUnitLineAlt{BC} are also greater than \drawUnitLineAlt{CA}, and \drawUnitLineAlt{BC},\drawUnitLineAlt{CA} than \drawUnitLineAlt{AB}. \par
    Therefore etc. 
}{}

\iftoggle{raum}{
    Raum_Text
}{}
\nolinenumbers
\hfill Q. E. D.

\resetlinenumber
}

\resetlinenumber

\startProposition{
    \markboth{Euclid in Color}{Proposition I.21}
    \defineNewPicture[1/4][3]{
	pair B, C, A, D, E;
	B := (0, 0);
	C := B shifted (7/2u, 0);
	A := B shifted (u, 5/2u);
	D := 1/2[1/2[B, C], A];
	E = whatever[B, D] = whatever[C, A];
	byAngleDefine(C, D, B, byred, SOLID_SECTOR);
	byAngleDefine(C, E, D, byblue, SOLID_SECTOR);
	byAngleDefine(C, A, B, byyellow, SOLID_SECTOR);
	draw byNamedAngleResized();
	byLineDefine(C, D, byyellow, SOLID_LINE, REGULAR_WIDTH);
	byLineDefine(B, D, byblack, SOLID_LINE, REGULAR_WIDTH);
	byLineDefine(D, E, byblack, SOLID_LINE, REGULAR_WIDTH);
	byLineDefine(B, C, bygreen, SOLID_LINE, REGULAR_WIDTH);
	byLineDefine(C, E, byred, SOLID_LINE, REGULAR_WIDTH);
	byLineDefine(E, A, byred, SOLID_LINE, REGULAR_WIDTH);
	byLineDefine(A, B, byblue, SOLID_LINE, REGULAR_WIDTH);
	draw byNamedLine(CD);
	draw byNamedLineSeq(0)(BD,DE);
	draw byNamedLineSeq(0)(AB,EA,CE,BC);
	draw byLabelsOnPolygon(B, A, E, C)(ALL_LABELS, 0);
	draw byLabelsOnPolygon(B, D, E)(OMIT_FIRST_LABEL+OMIT_LAST_LABEL, 0);
    }
    ~\hfill\drawCurrentPicture\hfill~
}{
    \chapter*{Proposition I.21}
    \renewcommand{\currentpropheader}{Proposition I.21}
    \addcontentsline{toc}{chapter}{Proposition 21}
    \textit{If on one of the sides of a triangle, from its extremities, there be constructed two straight lines meeting within the triangle, the straight lines so constructed will be less than the remaining two sides of the triangle, but will contain a greater angle.}
    \vskip \baselineskip

    % Check line numbering toggle
\iftoggle{lineNumbering}{\linenumbers}{}

\iftoggle{heath}{
    On \drawUnitLineAlt{BC}, one of the sides of the triangle \drawLineAlt{A,B,C}, from its extremities \drawPointAlt{B},  \drawPointAlt{C}, let the two straight lines \drawUnitLineAlt{BD},\drawUnitLineAlt{DC} be constructed meeting within the triangle; \par
    I say that \drawUnitLineAlt{BD},\drawUnitLineAlt{DC} are less than the remaining two sides of the triangle,\drawUnitLineAlt{BA},\drawUnitLineAlt{AC}, but contain an angle \drawAngleAlt{BDC} greater than the angle \drawAngleAlt{BAC}. 
    \vskip \baselineskip
    For let \drawUnitLineAlt{BD} be drawn through to  \drawPointAlt{E}. \par
    Then, since in any triangle two sides are greater than the remaining one, \justification{I. 20}
    {\centering therefore, in the triangle \hspace{-0.5em} \drawLineAlt{A,B,E}, the two sides \drawUnitLineAlt{AB},\drawUnitLineAlt{AE} are greater than \drawUnitLineAlt{BE}. \par}
    Let \drawUnitLineAlt{EC} be added to each; \par
    {\centering therefore \drawUnitLineAlt{BA},\drawUnitLineAlt{AC} are greater than \drawUnitLineAlt{BE},\drawUnitLineAlt{EC}. \par}
    Again, since, in the triangle \drawLineAlt{C,E,D}, the two sides \drawUnitLineAlt{CE},\drawUnitLineAlt{ED} are greater than \drawUnitLineAlt{CD}, let \drawUnitLineAlt{DB} be added to each; \par
    {\centering therefore \drawUnitLineAlt{CE},\drawUnitLineAlt{EB} are greater than \drawUnitLineAlt{CD},\drawUnitLineAlt{DB}. \par}
    But \drawUnitLineAlt{BA},\drawUnitLineAlt{AC} were proved greater than \drawUnitLineAlt{BE},\drawUnitLineAlt{EC};\par
    {\centering therefore \drawUnitLineAlt{BA},\drawUnitLineAlt{AC} are much greater than \drawUnitLineAlt{BD},\drawUnitLineAlt{DC}. \par}
    Again, since in any triangle the exterior angle is greater than the interior and opposite angle,\justification{I. 16}
    therefore, in the triangle \drawLineAlt{C,D,E}, \par
    {\centering the exterior angle \drawAngleAlt{BDC} is greater than the angle \drawAngleAlt{CED}. \par}
    For the same reason, moreover, in the triangle \drawLineAlt{A,B,E} also, \par
    {\centering the exterior angle \drawAngleAlt{CEB} is greater than the angle \drawAngleAlt{BAC}. \par}
    But the angle \drawAngleAlt{BDC} was proved greater than the angle \drawAngleAlt{CEB}; \par
    {\centering therefore the angle \drawAngleAlt{BDC} is much greater than the angle \drawAngleAlt{BAC}. \par}
    Therefore etc. 
}{}

\iftoggle{raum}{
    Raum_Text
}{}
\nolinenumbers
\hfill Q. E. D.

\vskip \baselineskip


\resetlinenumber
}

\resetlinenumber

\startProposition{
    \markboth{Euclid in Color}{Proposition I.22}
    \defineNewPicture[1/2][3/2]{
	numeric r[], d;
	pair E, F, G, K, D, H, L, Ai, Aii, Bi, Bii, Ci, Cii;
	path q[];
	r1 := 4/3u; % length of C
	r2 := 2u; % length of A
	r3 := (1/2)*(r1+r2); % length of B
	d := 2/3u; % y separation between A, B, C
	F := (0, 0);
	G := F shifted (r3, 0);
	q1 := (fullcircle scaled 2r2) shifted F;
	q2 := (fullcircle scaled 2r1) shifted G;
	K := q1 intersectionpoint q2;
	L:= (xpart(K), -ypart(K));
	D := point 4 of q1;
	H := point 0 of q2;
	E := H xshifted u;
	Bi := (xpart(point 0 of q2) + 1.5u, 0);
	Bii := Bi shifted (r3, 0);
	Ai := Bi shifted (0, d);
	Aii := Ai shifted (r2, 0);
	Ci := Bi shifted (0, -d);
	Cii := Ci shifted (r1, 0);
	draw byCircleABC(D, K, L, byblue, 0, 0, 0);
	byLineDefine(F, D, byblue, SOLID_LINE, REGULAR_WIDTH);
	byLineDefine(G, H, byred, SOLID_LINE, REGULAR_WIDTH);
	byLineDefine(F, G, byblack, SOLID_LINE, REGULAR_WIDTH);
	byLineDefine(G, K, byyellow, SOLID_LINE, REGULAR_WIDTH);
	byLineDefine(K, F, bygreen, SOLID_LINE, REGULAR_WIDTH);
	byLineDefine(H, E, bygreen, SOLID_LINE, REGULAR_WIDTH);
	draw byNamedLineSeq(0)(GK,KF, HE);
	draw byNamedLineSeq(0)(FD,FG,GH);
	draw byLineWithName (Ai, Aii, bygreen, SOLID_LINE, 0)(A);
	draw byLineWithName (Bi, Bii, byblack, SOLID_LINE, 0)(B);
	draw byLineWithName (Ci, Cii, byyellow, SOLID_LINE, 0)(C);
	draw byCircleABC(L, K, H, byred, 0, 0, 0);
	draw byLabelsOnPolygon(F, K, G)(OMIT_FIRST_LABEL+OMIT_LAST_LABEL, 0);
	draw byLabelsOnCircle(D)(DKL);
	draw byLabelLine(0)(A, B, C);
	draw byLabelPoint(F, -90, 2);
	draw byLabelPoint(G, -90, 2);
	draw byLabelPoint(L, -90, 2);
	draw byLabelPoint(H, -45, 3);
	draw byLabelPoint(E, -90, 2);
    }
    ~\hfill\drawCurrentPicture\hfill~
}{
    \chapter*{Proposition I.22}
    \renewcommand{\currentpropheader}{Proposition I.22}
    \addcontentsline{toc}{chapter}{Proposition 22}
    \textit{Out of three straight lines, which are equal to three given straight lines, to construct a triangle: thus it is necessary that two of the straight lines taken together in any manner should be greater than the remaining one. [I. 20]}
    \vskip \baselineskip

    % Check line numbering toggle
\iftoggle{lineNumbering}{\linenumbers}{}

\iftoggle{heath}{
    Let the three given straight lines be \drawPointAlt{A}, \drawPointAlt{B}, \drawPointAlt{C}, and of these let two taken together in any manner be greater than the remaining one, namely \par
    {\centering \drawPointAlt{A}, \drawPointAlt{B} greater than \drawPointAlt{C}; \par}
    {\centering \drawPointAlt{A}, \drawPointAlt{C} greater than \drawPointAlt{B}; \par}
    {\centering and \drawPointAlt{B}, \drawPointAlt{C} greater than \drawPointAlt{A}; \par}
    thus it is required to construct a triangle out of straight lines equal to \drawPointAlt{A}, \drawPointAlt{B}, \drawPointAlt{C}. 
    \vskip \baselineskip
    Let there be set out a straight line \drawUnitLineAlt{DE}, terminated at \drawPointAlt{D} but of infinite length in the direction of \drawPointAlt{E}, \par
    and let \drawUnitLineAlt{DF} be made equal to \drawPointAlt{A},\drawUnitLineAlt{FG} equal to \drawPointAlt{B}, and \drawUnitLineAlt{GH} equal to \drawPointAlt{C}. \justification{I. 3}
    With center \drawPointAlt{F} and distance \drawUnitLineAlt{FD} let the circle \drawCircleAlt[byblue]{DKL} be described; \par
    again, with center \drawPointAlt{G} and distance \drawUnitLineAlt{GH} let the circle \drawCircleAlt[byred]{KLH} be described; \par
    and let \drawUnitLineAlt{KF},\drawUnitLineAlt{KG} be joined; \par
    I say that the triangle \drawLineAlt{K,F,G} has been constructed out of three straight lines equal to  \drawPointAlt{A}, \drawPointAlt{B}, \drawPointAlt{C}. \par
    For, since the point \drawPointAlt{F} is the center of the circle \drawCircleAlt[byblue]{DKL},\par
    {\centering  \drawUnitLineAlt{FD} is equal to \drawUnitLineAlt{FK}. \par}
    But \drawUnitLineAlt{FD} is equal to \drawPointAlt{A}; \par
    {\centering therefore \drawUnitLineAlt{KF} is also equal to \drawPointAlt{A}. \par}
    Again, since the point \drawPointAlt{G} is the center of the circle \drawCircleAlt[byred]{LKH},\par
    {\centering  \drawUnitLineAlt{GH} is equal to \drawUnitLineAlt{GK}. \par}
    But \drawUnitLineAlt{GH} is equal to \drawPointAlt{C}; \par
    {\centering therefore \drawUnitLineAlt{KG} is also equal to \drawPointAlt{C}. \par}
    And \drawUnitLineAlt{FG} is also equal to \drawPointAlt{B}; therefore the three straight lines \drawUnitLineAlt{KF},\drawUnitLineAlt{FG},\drawUnitLineAlt{GK} are equal to the three straight lines  \drawPointAlt{A}, \drawPointAlt{B}, \drawPointAlt{C}. \par
    Therefore out of the three straight lines \drawUnitLineAlt{KF},\drawUnitLineAlt{FG},\drawUnitLineAlt{GK}, which are equal to the three given straight lines  \drawPointAlt{A}, \drawPointAlt{B}, \drawPointAlt{C}, the triangle \drawLineAlt{K,F,G} has been constructed. 
}{}

\iftoggle{raum}{
    Raum_Text
}{}
\nolinenumbers
\hfill Q. E. F.

\vskip \baselineskip


\resetlinenumber
}

\resetlinenumber

\startProposition{
    \markboth{Euclid in Color}{Proposition I.23}
    \defineNewPicture[1/2][5/4]{
	pair C, E, D, J, H, A, G, F, B, d;
	C := (0, 0) rotated -30;
	E := C shifted (7/2u, 0) rotated -30;
	D := C shifted (3u, 11/5u) rotated -30;
	J := 5/4[C, E];
	H := 7/6[C, D];
	d := (-23/4u, -2u);
	A := C rotated 30 shifted d;
	G := E rotated 30 shifted d;
	F := D rotated 30 shifted d;
	B := J rotated 30 shifted d;
	byAngleDefine(E, C, D, byred, SOLID_SECTOR);
	byAngleDefine(G, A, F, byblue, SOLID_SECTOR);
	draw byNamedAngleResized();
	byLineDefine(E, J, byblack, SOLID_LINE, REGULAR_WIDTH);
	byLineDefine(D, H, bygreen, SOLID_LINE, REGULAR_WIDTH);
	byLineDefine(C, E, byblack, SOLID_LINE, REGULAR_WIDTH);
	byLineDefine(D, C, bygreen, SOLID_LINE, REGULAR_WIDTH);
	draw byLine(E, D, byred, SOLID_LINE, REGULAR_WIDTH);
	draw byNamedLineSeq(0)(DH,DC,CE,ED, EJ);
	byLineDefine(G, B, byblack, SOLID_LINE, REGULAR_WIDTH);
	byLineDefine(A, G, byblack, SOLID_LINE, REGULAR_WIDTH);
	byLineDefine(G, F, byred, SOLID_LINE, REGULAR_WIDTH);
	byLineDefine(F, A, byyellow, SOLID_LINE, REGULAR_WIDTH);
	draw byNamedLineSeq(0)(noLine,GF,FA,AG,GB);
	draw byLabelsOnPolygon(J, E, C, D, H)(OMIT_FIRST_LABEL+OMIT_LAST_LABEL, 0);
	draw byLabelsOnPolygon(noPoint,B, G, A, F, G)(OMIT_FIRST_LABEL+OMIT_LAST_LABEL, 0);
    }
    ~\hfill\drawCurrentPicture\hfill~
}{
    \chapter*{Proposition I.23}
    \renewcommand{\currentpropheader}{Proposition I.23}
    \addcontentsline{toc}{chapter}{Proposition 23}
    \textit{On a given straight line and at a point on it to construct a rectilineal angle equal to a given rectilineal angle.}
    \vskip \baselineskip

    % Check line numbering toggle
\iftoggle{lineNumbering}{\linenumbers}{}

\iftoggle{heath}{
    Let \drawUnitLineAlt{AB} be the given straight line, \drawPointAlt{A} the point on it, and the angle \drawAngleAlt{DCE} the given rectilineal angle; \par
    thus it is required to construct on the given straight line \drawUnitLineAlt{AB}, and at the point \drawPointAlt{A} on it, a rectilineal angle equal to the given rectilineal angle \drawAngleAlt{DCE}. 
    \vskip \baselineskip
    On the straight lines \drawUnitLineAlt{CD},\drawUnitLineAlt{CE} respectively let the points \drawPointAlt{D}, \drawPointAlt{E} be taken at random; let \drawUnitLineAlt{DE} be joined, \par
    and out of three straight lines which are equal to the three straight lines \drawUnitLineAlt{CD},\drawUnitLineAlt{DE},\drawUnitLineAlt{CE} let the triangle \drawLineAlt{A,F,G} be constructed in such a way that \drawUnitLineAlt{CD} is equal to \drawUnitLineAlt{AF},\drawUnitLineAlt{CE} to \drawUnitLineAlt{AG}, and further \drawUnitLineAlt{DE} to \drawUnitLineAlt{FG}. \justification{I.22}
    Then, since the two sides \drawUnitLineAlt{DC},\drawUnitLineAlt{CE} are equal to the two sides \drawUnitLineAlt{FA},\drawUnitLineAlt{AG} respectively, and the base \drawUnitLineAlt{DE} is equal to the base \drawUnitLineAlt{FG}, \par
    \centeredjust{the angle \drawAngleAlt{DCE} is equal to the angle \drawAngleAlt{FAG}.}{I. 8}
    Therefore on the given straight line \drawUnitLineAlt{AB}, and at the point \drawPointAlt{A} on it, the rectilineal angle \drawAngleAlt{FAG} has been constructed equal to the given rectilineal angle \drawAngleAlt{DCE}. 
}{}

\iftoggle{raum}{
    Raum_Text
}{}
\nolinenumbers
\hfill Q. E. F.

\resetlinenumber
}

\resetlinenumber

\startProposition{
    \markboth{Euclid in Color}{Proposition I.24}
    \defineNewPicture[1/3][3/2]{
	pair D, E, F, G, A, B, C, d;
	D := (0, 0);
	E := D shifted (u, -5/2u);
	F := D shifted (-u, -7/2u);
	G := ((fullcircle scaled 2abs(F-D)) shifted D) intersectionpoint (E--E shifted (-2abs(F-D), -u));
	d := (-5u, 0);
	A := D shifted d;
	B := E shifted d;
	C := G shifted d;
	byAngleDefine(E, D, F, bygreen, SOLID_SECTOR);
	byAngleDefine(F, D, G, byblack, SOLID_SECTOR);
	byAngleDefine(B, A, C, byyellow, SOLID_SECTOR);
	byAngleDefine(E, G, D, byblue, SOLID_SECTOR);
	byAngleDefine(F, G, E, byred, SOLID_SECTOR);
	byAngleDefine(G, F, D, byyellow, SOLID_SECTOR);
	byAngleDefine(D, F, E, byblack, SOLID_SECTOR);
	draw byNamedAngleResized();
	byLineDefine(D, E, byblue, SOLID_LINE, REGULAR_WIDTH);
	byLineDefine(E, F, byyellow, SOLID_LINE, REGULAR_WIDTH);
	draw byLine(F, D, byred, SOLID_LINE, REGULAR_WIDTH);
	draw byLine(E, G, byblack, SOLID_LINE, REGULAR_WIDTH);
	byLineDefine(D, G, bygreen, SOLID_LINE, REGULAR_WIDTH);
	byLineDefine(F, G, byblue, SOLID_LINE, REGULAR_WIDTH);
	draw byNamedLineSeq(0)(DG,FG,EF,DE);
	byLineDefine(A, B, byblue, SOLID_LINE, REGULAR_WIDTH);
	byLineDefine(B, C, byyellow, SOLID_LINE, REGULAR_WIDTH);
	byLineDefine(C, A, byred, SOLID_LINE, REGULAR_WIDTH);
	draw byNamedLineSeq(0)(AB,BC,CA);
	draw byLabelsOnPolygon(G, D, E, F)(ALL_LABELS, 0);
	draw byLabelsOnPolygon(C, A, B)(ALL_LABELS, 0);
    }
    ~\hfill\drawCurrentPicture\hfill~
}{
    \chapter*{Proposition I.24}
    \renewcommand{\currentpropheader}{Proposition I.24}
    \addcontentsline{toc}{chapter}{Proposition 24}
    \textit{If two triangles have the two sides equal to two sides respectively, but have the one of the angles contained by the equal straight lines greater than the other, they will also have the base greater than the base.}
    \vskip \baselineskip

    % Check line numbering toggle
\iftoggle{lineNumbering}{\linenumbers}{}

\iftoggle{heath}{
    Let \drawLineAlt{A,B,C}, \drawLineAlt{D,E,F} be two triangles having the two sides \drawUnitLineAlt{AB},\drawUnitLineAlt{AC} equal to the two sides \drawUnitLineAlt{DE},\drawUnitLineAlt{DF} respectively, namely \drawUnitLineAlt{AB} to \drawUnitLineAlt{DE}, and \drawUnitLineAlt{AC} to \drawUnitLineAlt{DF}, and let the angle at \drawPointAlt{A} be greater than the angle at \drawPointAlt{D}; \par
    I say that the base \drawUnitLineAlt{BC} is also greater than the base \drawUnitLineAlt{EF}. 
    \vskip \baselineskip
    For, since the angle \drawAngleAlt{BAC} is greater than the angle \drawAngleAlt{EDF}, let there be constructed, on the straight line \drawUnitLineAlt{DE}, and at the point \drawPointAlt{D} on it, the angle \drawAngleAlt{EDG} equal to the angle \drawAngleAlt{BAC}; \justification{I. 23}
    let \drawUnitLineAlt{DG} be made equal to either of the two straight lines \drawUnitLineAlt{AC},\drawUnitLineAlt{DF}, and let \drawUnitLineAlt{EG},\drawUnitLineAlt{FG} be joined. \par
    Then, since \drawUnitLineAlt{AB} is equal to \drawUnitLineAlt{DE}, and \drawUnitLineAlt{AC} to \drawUnitLineAlt{DG}, \par
    {\centering the two sides \drawUnitLineAlt{BA},\drawUnitLineAlt{AC} are equal to the two sides \drawUnitLineAlt{ED},\drawUnitLineAlt{DG}, respectively; \par}
    and the angle \drawAngleAlt{BAC} is equal to the angle \drawAngleAlt{EDG}; \par
    \centeredjust{therefore the base \drawUnitLineAlt{BC} is equal to the base \drawUnitLineAlt{EG}.}{I. 4}
    Again, since \drawUnitLineAlt{DF} is equal to \drawUnitLineAlt{DG}, \par
    \centeredjust{the angle \drawAngleAlt{DGF} is also equal to the angle \drawAngleAlt{DFG};}{I. 5}
    {\centering therefore the angle \drawAngleAlt{DFG} is greater than the angle \drawAngleAlt{EGF}. \par}
    {\centering Therefore the angle \drawAngleAlt{EFG} is much greater than the angle \drawAngleAlt{EGF}. \par}
    And, since \drawLineAlt{E,F,G} is a triangle having the angle \drawAngleAlt{EFG} greater than the angle \drawAngleAlt{EGF}, and the greater angle is subtended by the greater side, \par\justification{I. 19}
    {\centering the side \drawUnitLineAlt{EG} is also greater than \drawUnitLineAlt{EF}. \par}
    But \drawUnitLineAlt{EG} is equal to \drawUnitLineAlt{BC}.\par
    {\centering Therefore \drawUnitLineAlt{BC} is also greater than \drawUnitLineAlt{EF}. \par}
    Therefore etc. 
}{}

\iftoggle{raum}{
    Raum_Text
}{}
\nolinenumbers
\hfill Q. E. D.

\vskip \baselineskip


\resetlinenumber
}

\resetlinenumber

\startProposition{
    \markboth{Euclid in Color}{Proposition I.25}
    \defineNewPicture[1/4][1]{
	pair A, B, C, D, E, F, d;
	numeric x, y;
	d := (8u, 0);
	x := 4u;
	y := 5u;
	A := (0, 0);
	B := A + x*dir(240);
	C := A + y*dir(315);
	E := B shifted d;
	D := E + length(A-B)*dir(86);
	F := D + y*dir(306);
	byAngleDefine(B, A, C, byyellow, SOLID_SECTOR);
	byAngleDefine(E, D, F, byblack, SOLID_SECTOR);
	draw byNamedAngleResized();
	byLineDefine(A, B, byblue, SOLID_LINE, REGULAR_WIDTH);
	byLineDefine(B, C, byblack, SOLID_LINE, REGULAR_WIDTH);
	byLineDefine(C, A, byred, SOLID_LINE, REGULAR_WIDTH);
	draw byNamedLineSeq(0)(AB,BC,CA);
	byLineDefine(D, E, bygreen, SOLID_LINE, REGULAR_WIDTH);
	byLineDefine(E, F, byred, SOLID_LINE, REGULAR_WIDTH);
	byLineDefine(F, D, byyellow, SOLID_LINE, REGULAR_WIDTH);
	draw byNamedLineSeq(0)(DE,EF,FD);
	draw byLabelsOnPolygon(B, A, C)(ALL_LABELS, 0);
	draw byLabelsOnPolygon(E, D, F)(ALL_LABELS, 0);
    }
    ~\hfill\drawCurrentPicture\hfill~
}{
    \chapter*{Proposition I.25}
    \renewcommand{\currentpropheader}{Proposition I.25}
    \addcontentsline{toc}{chapter}{Proposition 25}
    \textit{If two triangles have the two sides equal to two sides respectively, but have the base greater than the base, they will also have the one of the angles contained by the equal straight lines greater than the other.}
    \vskip \baselineskip

    % Check line numbering toggle
\iftoggle{lineNumbering}{\linenumbers}{}

\iftoggle{heath}{
    Let \drawLineAlt{A,B,C}, \drawLineAlt{D,E,F} be two triangles having the two sides \drawUnitLineAlt{AB},\drawUnitLineAlt{AC} equal to the two sides \drawUnitLineAlt{DE},\drawUnitLineAlt{DF} respectively, namely \drawUnitLineAlt{AB} to \drawUnitLineAlt{DE}, and \drawUnitLineAlt{AC} to \drawUnitLineAlt{DF}; and let the base \drawUnitLineAlt{BC} be greater than the base \drawUnitLineAlt{EF}; \par
    I say that the angle \drawAngleAlt{BAC} is also greater than the angle \drawAngleAlt{EDF}. 
    \vskip \baselineskip
    For, if not, it is either equal to it or less. \par
    Now the angle \drawAngleAlt{BAC} is not equal to the angle \drawAngleAlt{EDF}; for then the base \drawUnitLineAlt{BC} would also have been equal to the base \drawUnitLineAlt{EF},\justification{I. 4}
    but it is not; \par
    {\centering therefore the angle \drawAngleAlt{BAC} is not equal to the angle \drawAngleAlt{EDF}. \par}
    Neither again is the angle \drawAngleAlt{BAC} less than the angle \drawAngleAlt{EDF}; for then the base \drawUnitLineAlt{BC} would also have been less than the base \drawUnitLineAlt{EF},\justification{I. 24}
    but it is not; \par
    {\centering therefore the angle \drawAngleAlt{BAC} is not less than the angle \drawAngleAlt{EDF}. \par}
    But it was proved that it is not equal either; \par
    {\centering therefore the angle \drawAngleAlt{BAC} is greater than the angle \drawAngleAlt{EDF}. \par}
    Therefore etc. 
}{}

\iftoggle{raum}{
    Raum_Text
}{}
\nolinenumbers
\hfill Q. E. D.

\resetlinenumber
}

\resetlinenumber

\startProposition{
    \markboth{Euclid in Color}{Proposition I.26}
    \defineNewPicture[1/3][3/2]{
	pair A, B, C, D, E, F, G, H, d;
	B := (0, 0);
	C := B shifted (3u, 0);
	A := B shifted (1/2u, 3u);
	d := (5u, 0);
	E := B shifted d;
	F := C shifted d;
	D := A shifted d;
	G := 3/4[B, A];
	H := 3/4[B, C];
	byAngleDefine(C, B, A, bygreen, SOLID_SECTOR);
	byAngleDefine(A, C, G, byred, SOLID_SECTOR);
	byAngleDefine(G, C, B, byblue, SOLID_SECTOR);
	byAngleDefine(H, A, C, byyellow, SOLID_SECTOR);
	byAngleDefine(B, A, H, byred, SOLID_SECTOR);
	byAngleDefine(A, H, B, byblack, SOLID_SECTOR);
	draw byNamedAngleResized(CBA, ACG, BAH, GCB, AHB, HAC);
	byLineDefine(B, C, byred, SOLID_LINE, REGULAR_WIDTH);
	byLineDefine(C, A, byblack, SOLID_LINE, REGULAR_WIDTH);
	byLineDefine(G, B, byyellow, SOLID_LINE, REGULAR_WIDTH);
	byLineDefine(A, G, byyellow, SOLID_LINE, REGULAR_WIDTH);
	byLineDefine(C, G, bygreen, SOLID_LINE, REGULAR_WIDTH);
	byLineDefine(A, H, byblue, SOLID_LINE, REGULAR_WIDTH);
	byLineDefine(B, H, byred, SOLID_LINE, REGULAR_WIDTH);
	byLineDefine(H, C, byred, SOLID_LINE, REGULAR_WIDTH);
	draw byNamedLineSeq(0)(AH,CG);
	draw byNamedLineSeq(0)(AG,GB,BC,CA);
	byAngleDefine(F, E, D, bygreen, SOLID_SECTOR);
	byAngleDefine(D, F, E, byblue, SOLID_SECTOR);
	byAngleDefine(E, D, F, byyellow, SOLID_SECTOR);
	draw byNamedAngleResized(FED, DFE, EDF);
	byLineDefine(E, F, byblack, SOLID_LINE, REGULAR_WIDTH);
	byLineDefine(F, D, bygreen, SOLID_LINE, REGULAR_WIDTH);
	byLineDefine(D, E, byred, SOLID_LINE, REGULAR_WIDTH);
	draw byNamedLineSeq(0)(DE,FD,EF);
	draw byLabelsOnPolygon(B, G, A, C, H)(ALL_LABELS, 0);
	draw byLabelsOnPolygon(E, D, F)(ALL_LABELS, 0);
    }
    ~\hfill\drawCurrentPicture\hfill~
}{
    \chapter*{Proposition I.26}
    \renewcommand{\currentpropheader}{Proposition I.26}
    \addcontentsline{toc}{chapter}{Proposition 26}
    \textit{If two triangles have the two angles equal to two angles respectively, and one side equal to one side, namely, either the side adjoining the equal angles, or that subtending one of the equal angles, they will also have the remaining sides equal to the remaining sides and the remaining angle to the remaining angle.}
    \vskip \baselineskip

    % Check line numbering toggle
\iftoggle{lineNumbering}{\linenumbers}{}

\iftoggle{heath}{
    Let \drawLineAlt{A,B,C}, \drawLineAlt{D,E,F} be two triangles having the two angles \drawAngleAlt{ABC}, \drawAngleAlt{BCA} equal to the two angles \drawAngleAlt{DEF}, \drawAngleAlt{EFD} respectively, namely the angle \drawAngleAlt{ABC} to the angle \drawAngleAlt{DEF}, and the angle \drawAngleAlt{BCA} to the angle \drawAngleAlt{EFD}; and let them also have one side equal to one side, first that adjoining the equal angles, namely \drawUnitLineAlt{BC} to \drawUnitLineAlt{EF}; \par
    I say that they will also have the remaining sides equal to the remaining sides respectively, namely \drawUnitLineAlt{AB} to \drawUnitLineAlt{DE} and \drawUnitLineAlt{AC} to \drawUnitLineAlt{DF}, and the remaining angle to the remaining angle, namely the angle \drawAngleAlt{BAC} to the angle \drawAngleAlt{EDF}. 
    \vskip \baselineskip
    For, if \drawUnitLineAlt{AB} is unequal to \drawUnitLineAlt{DE}, one of them is greater. \par
    Let \drawUnitLineAlt{AB} be greater, and let \drawUnitLineAlt{BG} be made equal to \drawUnitLineAlt{DE}; and let \drawUnitLineAlt{GC} be joined. \par
    Then, since \drawUnitLineAlt{BG} is equal to \drawUnitLineAlt{DE}, and \drawUnitLineAlt{BC} to \drawUnitLineAlt{EF}, \par
    {\centering the two sides \drawUnitLineAlt{GB},\drawUnitLineAlt{BC} are equal to the two sides \drawUnitLineAlt{DE},\drawUnitLineAlt{EF} respectively; \par}
    and the angle \drawAngleAlt{GBC} is equal to the angle \drawAngleAlt{DEF}; \par
    {\centering therefore the base \drawUnitLineAlt{GC} is equal to the base \drawUnitLineAlt{DF}, \par}
    {\centering and the triangle \drawLineAlt{G,B,C} is equal to the triangle \drawLineAlt{D,E,F}, \par}
    and the remaining angles will be equal to the remaining angles, namely those which the equal sides subtend;\justification{I. 4}
    {\centering therefore the angle \drawAngleAlt{GCB} is equal to the angle \drawAngleAlt{DFE}. \par}
    But the angle \drawAngleAlt{DFE} is by hypothesis equal to the angle \drawAngleAlt{BCA}; \par
    {\centering therefore the angle \drawAngleAlt{BCG} is equal to the angle \drawAngleAlt{BCA}, \par}
    the less to the greater: which is impossible.\par
    Therefore \drawUnitLineAlt{AB} is not unequal to \drawUnitLineAlt{DE}, \par
    {\centering and is therefore equal to it. \par}
    But \drawUnitLineAlt{BC} is also equal to \drawUnitLineAlt{EF}; \par
therefore the two sides \drawUnitLineAlt{AB},\drawUnitLineAlt{BC} are equal to the two sides \drawUnitLineAlt{DE},\drawUnitLineAlt{EF} respectively, \par
	\nolinenumbers
    \clearpage
    \linenumbers
    \enlargethispage{2\baselineskip}
    and the angle \drawAngleAlt{ABC} is equal to the angle \drawAngleAlt{DEF}; \par
    {\centering therefore the base \drawUnitLineAlt{AC} is equal to the base \drawUnitLineAlt{DF},\par}
    \centeredjust{and the remaining angle \drawAngleAlt{BAC} is equal to the remaining angle \drawAngleAlt{EDF}.}{I. 4}
    Again, let sides subtending equal angles be equal, as \drawUnitLineAlt{AB} to \drawUnitLineAlt{DE}; \par
    I say again that the remaining sides will be equal to the remaining sides, namely \drawUnitLineAlt{AC} to \drawUnitLineAlt{DF} and \drawUnitLineAlt{BC} to \drawUnitLineAlt{EF}, and further the remaining angle \drawAngleAlt{BAC} is equal to the remaining angle \drawAngleAlt{EDF}. \par
    For, if \drawUnitLineAlt{BC} is unequal to \drawUnitLineAlt{EF}, one of them is greater. \par
    Let \drawUnitLineAlt{BC} be greater, if possible, and let \drawUnitLineAlt{BH} be made equal to \drawUnitLineAlt{EF}; let \drawUnitLineAlt{AH} be joined. \par
    Then, since \drawUnitLineAlt{BH} is equal to \drawUnitLineAlt{EF}, and \drawUnitLineAlt{AB} to \drawUnitLineAlt{DE}, \par
    {\centering the two sides \drawUnitLineAlt{AB},\drawUnitLineAlt{BH} are equal to the two sides \drawUnitLineAlt{DE},\drawUnitLineAlt{EF} respectively, \par}
    and they contain equal angles; \par
    {\centering therefore the base \drawUnitLineAlt{AH} is equal to the base \drawUnitLineAlt{DF}, \par}
    {\centering and the triangle \drawLineAlt{A,B,H} is equal to the triangle \drawLineAlt{D,E,F}, \par}
    and the remaining angles will be equal to the remaining angles, namely those which the equal sides subtend;\justification{I. 4}
    {\centering therefore the angle \drawAngleAlt{BHA} is equal to the angle \drawAngleAlt{EFD}. \par}
    But the angle \drawAngleAlt{EFD} is equal to the angle \drawAngleAlt{BCA}; \par
    therefore, in the triangle \drawLineAlt{A,H,C}, the exterior angle \drawAngleAlt{BHA} is equal to the interior and opposite angle \drawAngleAlt{BCA}: which is impossible.\justification{I. 16}
    Therefore \drawUnitLineAlt{BC} is not unequal to \drawUnitLineAlt{EF}, \par
    {\centering and is therefore equal to it. \par}
    But \drawUnitLineAlt{AB} is also equal to \drawUnitLineAlt{DE}; \par
    therefore the two sides \drawUnitLineAlt{AB},\drawUnitLineAlt{BC} are equal to the two sides \drawUnitLineAlt{DE},\drawUnitLineAlt{EF} respectively, and they contain equal angles; \par
    {\centering therefore the base \drawUnitLineAlt{AC} is equal to the base \drawUnitLineAlt{DF}, \par}
    {\centering the triangle \drawLineAlt{A,B,C} equal to the triangle \drawLineAlt{D,E,F}, \par}
\ and the remaining angle \drawAngleAlt{BAC} equal to the remaining angle \drawAngleAlt{EDF}.\justification{I. 4}
    Therefore etc. 
}{}

\iftoggle{raum}{
    Raum_Text
}{}
\nolinenumbers
\hfill Q. E. D.

\vskip \baselineskip

}

\resetlinenumber

\startProposition{
    \markboth{Euclid in Color}{Proposition I.27}
    \defineNewPicture[1/2][2]{
	pair A, B, C, D, E, F, G, H, I, d;
	A := (0, 0);
	B := A shifted (4u, 0);
	d := (0, -3/2u);
	C := A shifted d;
	D := B shifted d;
	E := 2/3[A, B];
	F := 1/3[C, D];
	G := 1/2[B, D] shifted (2u, 0);
	H := 3/2[E, F];
	I := 3/2[F, E];
	byAngleDefine(A, E, H, byyellow, SOLID_SECTOR);
	byAngleDefine(H, E, B, byred, SOLID_SECTOR);
	byAngleDefine(C, F, I, byblue, SOLID_SECTOR);
	byAngleDefine(I, F, D, byyellow, SOLID_SECTOR);
	draw byNamedAngleResized();
	byLineDefine(B, G, byblue, SOLID_LINE, REGULAR_WIDTH);
	byLineDefine(A, B, byblue, SOLID_LINE, REGULAR_WIDTH);
	byLineDefine(E, B, byblue, SOLID_LINE, REGULAR_WIDTH);
	byLineDefine(D, G, byred, SOLID_LINE, REGULAR_WIDTH);
	byLineDefine(C, D, byred, SOLID_LINE, REGULAR_WIDTH);
	byLineDefine(F, D, byred, SOLID_LINE, REGULAR_WIDTH);
	byLineDefine(I, H, byblack, SOLID_LINE, REGULAR_WIDTH);
	byLineDefine(E, F, byblack, SOLID_LINE, REGULAR_WIDTH);
	draw byNamedLineSeq(0)(BG,AB,DG,CD,IH);
	draw byLabelsOnPolygon(A, B, G, D, C)(ALL_LABELS, 0);
	draw byLabelsOnPolygon(A, E, I)(OMIT_FIRST_LABEL+OMIT_LAST_LABEL, 0);
	draw byLabelsOnPolygon(D, F, H)(OMIT_FIRST_LABEL+OMIT_LAST_LABEL, 0);
    }
    ~\hfill\drawCurrentPicture\hfill~
}{
    \chapter*{Proposition I.27}
    \renewcommand{\currentpropheader}{Proposition I.27}
    \addcontentsline{toc}{chapter}{Proposition 27}
    \textit{If a straight line falling on two straight lines make the alternate angles equal to one another, the straight lines will be parallel to one another.}
    \vskip \baselineskip

    % Check line numbering toggle
\iftoggle{lineNumbering}{\linenumbers}{}

\iftoggle{heath}{
    For let the straight line \drawUnitLineAlt{EF} falling on the two straight lines \drawUnitLineAlt{AB},\drawUnitLineAlt{CD} make the alternate angles \drawAngleAlt{AEF}, \drawAngleAlt{EFD} equal to one another; \par
    I say that \drawUnitLineAlt{AB} is parallel to \drawUnitLineAlt{CD}. 
    \vskip \baselineskip
    For, if not,\drawUnitLineAlt{AB},\drawUnitLineAlt{CD} when produced will meet either in the direction of \drawPointAlt{B}, \drawPointAlt{D} or towards \drawPointAlt{A}, \drawPointAlt{C}. \par
    Let them be produced and meet, in the direction of \drawPointAlt{B}, \drawPointAlt{D}, at \drawPointAlt{G}.\par
    Then, in the triangle \drawLineAlt[GEF]{G,B,E,F,D}, \par
    the exterior angle \drawAngleAlt{AEF} is equal to the interior and opposite angle \drawAngleAlt[EFG]{EFD}: which is impossible.
    \hspace*{1em}\justification{I. 16}
    {\centering Therefore \drawUnitLineAlt{AB},\drawUnitLineAlt{CD} when produced will not meet in the direction of \drawPointAlt{B}, \drawPointAlt{D}. \par}
    Similarly it can be proved that neither will they meet towards \drawPointAlt{A}, \drawPointAlt{C}. \par
    But straight lines which do not meet in either direction are parallel; \justification{Def. 23}
    {\centering therefore \drawUnitLineAlt{AB} is parallel to \drawUnitLineAlt{CD}. \par}
    Therefore etc. 
}{}

\iftoggle{raum}{
    Raum_Text
}{}
\nolinenumbers
\hfill Q. E. D.

\resetlinenumber
}

\resetlinenumber

\startProposition{
    \markboth{Euclid in Color}{Proposition I.28}
    \defineNewPicture[1][5/2]{
	pair A, B, C, D, E, F, G, H, d;
	A := (0, 0);
	B := A shifted (5u, 0);
	d := (0, -3/2u);
	C := A shifted d;
	D := B shifted d;
	G := 9/20[A, B];
	H := 11/20[C, D];
	E := 5/3[H, G];
	F := 5/3[G, H];
	byAngleDefine(E, G, A, byblack, SOLID_SECTOR);
	byAngleDefine(E, G, B, byyellow, SOLID_SECTOR);
	byAngleDefine(A, G, H, bygreen, SOLID_SECTOR);%
	byAngleDefine(B, G, H, byblue, SOLID_SECTOR);
	byAngleDefine(C, H, E, byblue, SOLID_SECTOR);
	byAngleDefine(E, H, D, byred, SOLID_SECTOR);
	draw byNamedAngleResized();
	draw byLine(A, B, byred, SOLID_LINE, REGULAR_WIDTH);
	draw byLine(C, D, byyellow, SOLID_LINE, REGULAR_WIDTH);
	draw byLine(E, F, byblack, SOLID_LINE, REGULAR_WIDTH);
	draw byLabelLine(0)(AB, CD, EF);
	draw byLabelsOnPolygon(E, G, B)(OMIT_FIRST_LABEL+OMIT_LAST_LABEL, 0);
	draw byLabelsOnPolygon(G, H, D)(OMIT_FIRST_LABEL+OMIT_LAST_LABEL, 0);
    }
    ~\hfill\drawCurrentPicture\hfill~
}{
    \chapter*{Proposition I.28}
    \renewcommand{\currentpropheader}{Proposition I.28}
    \addcontentsline{toc}{chapter}{Proposition 28}
    \textit{If a straight line falling on two straight lines make the exterior angle equal to the interior and opposite angle on the same side, or the interior angles on the same side equal to two right angles, the straight lines will be parallel to one another.}
    \vskip \baselineskip

    % Check line numbering toggle
\iftoggle{lineNumbering}{\linenumbers}{}

\iftoggle{heath}{
    For let the straight line \drawUnitLineAlt{EF} falling on the two straight lines \drawUnitLineAlt{AB},\drawUnitLineAlt{CD} make the exterior angle \drawAngleAlt{EGB} equal to the interior and opposite angle \drawAngleAlt{GHD}, or the interior angles on the same side, namely \drawAngleAlt{BGH}, \drawAngleAlt{GHD}, equal to two right angles; \par
    I say that \drawUnitLineAlt{AB} is parallel to \drawUnitLineAlt{CD}. 
    \vskip \baselineskip
    For, since the angle \drawAngleAlt{EGB} is equal to the angle \drawAngleAlt{GHD}, while the angle \drawAngleAlt{EGB} is equal to the angle \drawAngleAlt{AGH}, \justification{I. 15}
    {\centering the angle \drawAngleAlt{AGH} is also equal to the angle \drawAngleAlt{GHD}; \par}
    and they are alternate; \par
    \centeredjust{therefore \drawUnitLineAlt{AB} is parallel to \drawUnitLineAlt{CD}.}{I. 27}
    Again, since the angles \drawAngleAlt{BGH}, \drawAngleAlt{GHD} are equal to two right angles, and the angles \drawAngleAlt{AGH}, \drawAngleAlt{BGH} are also equal to two right angles, \justification{I. 13}
    {\centering the angles \drawAngleAlt{AGH}, \drawAngleAlt{BGH} are equal to the angles \drawAngleAlt{BGH}, \drawAngleAlt{GHD}. \par}
    Let the angle \drawAngleAlt{BGH} be subtracted from each;\par
    {\centering therefore the remaining angle \drawAngleAlt{AGH} is equal to the remaining angle \drawAngleAlt{GHD};\par}
    and they are alternate; \par
    \centeredjust{therefore \drawUnitLineAlt{AB} is parallel to \drawUnitLineAlt{CD}.}{I. 27}
    Therefore etc. 
}{}

\iftoggle{raum}{
    Raum_Text
}{}
\nolinenumbers
\hfill Q. E. D.

\resetlinenumber
}

\resetlinenumber

\startProposition{
    \markboth{Euclid in Color}{Proposition I.29}
    \defineNewPicture[1][5/2]{
	pair A, B, C, D, E, F, G, H, d;
	A := (0, 0);
	B := A shifted (5u, 0);
	d := (0, -3/2u);
	C := A shifted d;
	D := B shifted d;
	G := 9/20[A, B];
	H := 11/20[C, D];
	E := 5/3[H, G];
	F := 5/3[G, H];
	byAngleDefine(E, G, B, byyellow, SOLID_SECTOR);
	byAngleDefine(A, G, H, bygreen, SOLID_SECTOR);
	byAngleDefine(B, G, H, byblue, SOLID_SECTOR);
	byAngleDefine(E, H, D, byred, SOLID_SECTOR);
	draw byNamedAngleResized();
	draw byLine(A, B, byred, SOLID_LINE, REGULAR_WIDTH);
	draw byLine(C, D, byyellow, SOLID_LINE, REGULAR_WIDTH);
	draw byLine(E, F, byblack, SOLID_LINE, REGULAR_WIDTH);
	draw byLabelLine(0)(AB, CD, EF);
	draw byLabelsOnPolygon(E, G, B)(OMIT_FIRST_LABEL+OMIT_LAST_LABEL, 0);
	draw byLabelsOnPolygon(G, H, D)(OMIT_FIRST_LABEL+OMIT_LAST_LABEL, 0);
    }
    ~\hfill\drawCurrentPicture\hfill~
}{
    \chapter*{Proposition I.29}
    \renewcommand{\currentpropheader}{Proposition I.29}
    \addcontentsline{toc}{chapter}{Proposition 29}
    \textit{A straight line falling on parallel straight lines makes the alternate angles equal to one another, the exterior angle equal to the interior and opposite angle, and the interior angles on the same side equal to two right angles.}
    \vskip \baselineskip

    % Check line numbering toggle
\iftoggle{lineNumbering}{\linenumbers}{}

\iftoggle{heath}{
    For let the straight line \drawUnitLineAlt{EF} fall on the parallel straight lines \drawUnitLineAlt{AB},\drawUnitLineAlt{CD}; \par
    I say that it makes the alternate angles \drawAngleAlt{AGH}, \drawAngleAlt{GHD} equal, the exterior angle \drawAngleAlt{EGB} equal to the interior and opposite angle \drawAngleAlt{GHD}, and the interior angles on the same side, namely \drawAngleAlt{BGH}, \drawAngleAlt{GHD}, equal to two right angles. 
    \vskip \baselineskip
    For, if the angle \drawAngleAlt{AGH} is unequal to the angle \drawAngleAlt{GHD}, one of them is greater. \par
    Let the angle \drawAngleAlt{AGH} be greater. \par
    Let the angle \drawAngleAlt{BGH} be added to each;\par
    {\centering therefore the angles \drawAngleAlt{AGH}, \drawAngleAlt{BGH} are greater than the angles \drawAngleAlt{BGH}, \drawAngleAlt{GHD}. \par}
    But the angles \drawAngleAlt{AGH}, \drawAngleAlt{BGH} are equal to two right angles; \justification{I. 13}
    {\centering therefore the angles \drawAngleAlt{BGH}, \drawAngleAlt{GHD} are less than two right angles. \par}
    But straight lines produced indefinitely from angles less than two right angles meet; \justification{Post. 5}
    {\centering therefore \drawUnitLineAlt{AB},\drawUnitLineAlt{CD}, if produced indefinitely, will meet; \par}
    but they do not meet, because they are by hypothesis parallel. \par
    {\centering Therefore the angle \drawAngleAlt{AGH} is not unequal to the angle \drawAngleAlt{GHD}, \par}
    {\centering and is therefore equal to it. \par}
    Again, the angle \drawAngleAlt{AGH} is equal to the angle \drawAngleAlt{EGB}; \justification{I. 15}
    \centeredjust{therefore the angle \drawAngleAlt{EGB} is also equal to the angle \drawAngleAlt{GHD}.\hspace{1em}}{C. N. 1}
    Let the angle \drawAngleAlt{BGH} be added to each; \par
    {\centering therefore the angles \drawAngleAlt{EGB}, \drawAngleAlt{BGH} are equal to the angles\par} 
    \centeredjust{\drawAngleAlt{BGH}, \drawAngleAlt{GHD}.}{C. N. 2}
    But the angles \drawAngleAlt{EGB}, \drawAngleAlt{BGH} are equal to two right angles; \justification{I. 13}
    {\centering therefore the angles \drawAngleAlt{BGH}, \drawAngleAlt{GHD} are also equal to two right angles. \par}
    Therefore etc. 
}{}

\iftoggle{raum}{
    Raum_Text
}{}
\nolinenumbers
\hfill Q. E. D.

\vskip \baselineskip


\resetlinenumber
}

\resetlinenumber

\startProposition{
    \markboth{Euclid in Color}{Proposition I.30}
    \defineNewPicture[1][7/2]{
	pair A, B, C, D, E, F, G, H, I, J, K, d;
	A := (0, 0);
	B := A shifted (7/2u, 0);
	d := (0, -3/4u);
	E := A shifted d;
	F := B shifted d;
	C := E shifted 4/3d;
	D := F shifted 4/3d;
	G := 13/20[A, B];
	K := 7/20[C, D];
	H := (G--K) intersectionpoint (E--F);
	J := 5/4[K, G];
	I := 5/4[G, K];
	byAngleDefine(A, G, K, byyellow, SOLID_SECTOR);
	byAngleDefine(F, H, J, byblue, SOLID_SECTOR);
	byAngleDefine(D, K, J, byred, SOLID_SECTOR);
	draw byNamedAngleResized();
	draw byLine(A, B, byred, SOLID_LINE, REGULAR_WIDTH);
	draw byLine(E, F, byyellow, SOLID_LINE, REGULAR_WIDTH);
	draw byLine(C, D, byblue, SOLID_LINE, REGULAR_WIDTH);
	draw byLine(J, I, byblack, SOLID_LINE, REGULAR_WIDTH);
	byLineDefine(G, K, byblack, SOLID_LINE, REGULAR_WIDTH);
	draw byLabelLine(0)(AB, CD, EF);
	draw byLabelsOnPolygon(A, G, J)(OMIT_FIRST_LABEL+OMIT_LAST_LABEL, 0);
	draw byLabelsOnPolygon(E, H, J)(OMIT_FIRST_LABEL+OMIT_LAST_LABEL, 0);
	draw byLabelsOnPolygon(C, K, J)(OMIT_FIRST_LABEL+OMIT_LAST_LABEL, 0);
    }
    ~\hfill\drawCurrentPicture\hfill~
}{
    \chapter*{Proposition I.30}
    \renewcommand{\currentpropheader}{Proposition I.30}
    \addcontentsline{toc}{chapter}{Proposition 30}
    \textit{Straight lines parallel to the same straight line are also parallel to one another.}
    \vskip \baselineskip

    % Check line numbering toggle
\iftoggle{lineNumbering}{\linenumbers}{}

\iftoggle{heath}{
    Let each of the straight lines \drawUnitLineAlt{AB},\drawUnitLineAlt{CD} be parallel to \drawUnitLineAlt{EF}; \par
    I say that \drawUnitLineAlt{AB} is also parallel to \drawUnitLineAlt{CD}. 
    \vskip \baselineskip
    For let the straight line \drawUnitLineAlt{GK} fall upon them; \par
    Then, since the straight line \drawUnitLineAlt{GK} has fallen on the parallel straight lines \drawUnitLineAlt{AB},\drawUnitLineAlt{EF}, \par
    \centeredjust{the angle \drawAngleAlt{AGK} is equal to the angle \drawAngleAlt{GHF}.}{I. 29}
    Again, since the straight line \drawUnitLineAlt{GK} has fallen on the parallel straight lines \drawUnitLineAlt{EF},\drawUnitLineAlt{CD}, \par
    \centeredjust{the angle \drawAngleAlt{GHF} is equal to the angle \drawAngleAlt{GKD}.}{I. 29}
    But the angle \drawAngleAlt{AGK} was also proved equal to the angle \drawAngleAlt{GHF}; \par
    \centeredjust{therefore the angle \hspace{-.2em}\drawAngleAlt{AGK} is also equal to the angle \hspace{-.2em}\drawAngleAlt{GKD};\hspace{1.5em}}{C. N. 1}
    and they are alternate. \par
    Therefore \drawUnitLineAlt{AB} is parallel to \drawUnitLineAlt{CD}. 
}{}

\iftoggle{raum}{
    Raum_Text
}{}
\nolinenumbers
\hfill Q. E. D.

\resetlinenumber
}

\resetlinenumber

\startProposition{
    \markboth{Euclid in Color}{Proposition I.31}
    \defineNewPicture[1][2]{
	pair A, B, C, D, E, F, d;
	E := (0, 0);
	F := E shifted (6u, 0);
	d := (0, -2u);
	B := E shifted d;
	C := F shifted d;
	A := 5/9[E, F];
	D := 4/9[B, C];
	byAngleDefine(D, A, E, byyellow, SOLID_SECTOR);
	byAngleDefine(A, D, C, byred, SOLID_SECTOR);
	draw byNamedAngleResized();
	draw byLine(A, D, byblack, SOLID_LINE, REGULAR_WIDTH);
	byLineDefine(E, A, byred, SOLID_LINE, REGULAR_WIDTH);
	byLineDefine(A, F, bygreen, SOLID_LINE, REGULAR_WIDTH);
	byLineDefine(B, D, byblue, SOLID_LINE, REGULAR_WIDTH);
	byLineDefine(D, C, byblue, SOLID_LINE, REGULAR_WIDTH);
	draw byNamedLineSeq(0)(EA, AF);
	draw byNamedLineSeq(0)(BD, DC);
	draw byLabelsOnPolygon(E, A, F, noPoint, C, D, B, noPoint)(ALL_LABELS, 0);
    }
    ~\hfill\drawCurrentPicture\hfill~
}{
    \chapter*{Proposition I.31}
    \renewcommand{\currentpropheader}{Proposition I.31}
    \addcontentsline{toc}{chapter}{Proposition 31}
    \textit{Through a given point to draw a straight line parallel to a given straight line.}
    \vskip \baselineskip

    % Check line numbering toggle
\iftoggle{lineNumbering}{\linenumbers}{}

\iftoggle{heath}{
    Let \drawPointAlt{A} be the given point, and \drawUnitLineAlt{BC} the given straight line; \par
    thus it is required to draw through the point \drawPointAlt{A} a straight line parallel to the straight line \drawUnitLineAlt{BC}. 
    \vskip \baselineskip
    Let a point \drawPointAlt{D} be taken at random on \drawUnitLineAlt{BC}, and let \drawUnitLineAlt{AD} be joined; \par
    on the straight line \drawUnitLineAlt{DA}, and at the point \drawPointAlt{A} on it, let the angle \drawAngleAlt{DAE} be constructed equal to the angle \drawAngleAlt{ADC}; \justification{I. 23}
    and let the straight line \drawUnitLineAlt{AF} be produced in a straight line with \drawUnitLineAlt{EA}. \par
    Then, since the straight line \drawUnitLineAlt{AD} falling on the two straight lines \drawUnitLineAlt{BC},\drawUnitLineAlt{EF} has made the alternate angles \drawAngleAlt{EAD}, \drawAngleAlt{ADC} equal to one another, \par
    \centeredjust{therefore \drawUnitLineAlt{EA,AF} is parallel to \drawUnitLineAlt{BC}.}{I. 27}
    Therefore through the given point \drawPointAlt{A} the straight line \drawUnitLineAlt{EA,AF} has been drawn parallel to the given straight line \drawUnitLineAlt{BC}. 
}{}

\iftoggle{raum}{
    Raum_Text
}{}
\nolinenumbers
\hfill Q. E. F.

\resetlinenumber
}

\resetlinenumber

\startProposition{
    \markboth{Euclid in Color}{Proposition I.32}
    \defineNewPicture[1/3][7/4]{
	pair A, B, C, D, E;
	A := (0, 0);
	B := A shifted (-1u, -3u);
	C := A shifted (5/2u, -3u);
	D := 2[B, C];
	E := C shifted (unitvector(A-B) scaled 7/2u);
	byAngleDefine(D, C, E, byred, SOLID_SECTOR);
	byAngleDefine(E, C, A, bygreen, SOLID_SECTOR);
	byAngleDefine(C, A, B, byblack, SOLID_SECTOR);
	byAngleDefine(A, B, C, byyellow, SOLID_SECTOR);
	byAngleDefine(B, C, A, byblue, SOLID_SECTOR);
	draw byNamedAngleResized();
	draw byLine(C, E, byblue, SOLID_LINE, REGULAR_WIDTH);
	byLineDefine(C, D, byblack, SOLID_LINE, REGULAR_WIDTH);
	byLineDefine(A, B, bygreen, SOLID_LINE, REGULAR_WIDTH);
	byLineDefine(B, C, byred, SOLID_LINE, REGULAR_WIDTH);
	byLineDefine(C, A, byyellow, SOLID_LINE, REGULAR_WIDTH);
	draw byNamedLineSeq(0)(CD,CA,AB,BC);
	draw byLabelsOnPolygon(B, A, noPoint, E, B)(ALL_LABELS, 0);
	draw byLabelPoint(C, 90, -3);
	draw byLabelPoint(D, 90, -3);
    }
    ~\hfill\drawCurrentPicture\hfill~
}{
    \chapter*{Proposition I.32}
    \renewcommand{\currentpropheader}{Proposition I.32}
    \addcontentsline{toc}{chapter}{Proposition 32}
    \textit{In any triangle, if one of the sides be produced, the exterior angle is equal to the two interior and opposite angles, and the three interior angles of the triangle are equal to two right angles.}
    \vskip \baselineskip

    % Check line numbering toggle
\iftoggle{lineNumbering}{\linenumbers}{}

\iftoggle{heath}{
    Let \drawLineAlt{A,B,C} be a triangle, and let one side of it \drawUnitLineAlt{BC} be produced to \drawPointAlt{D}; \par
    I say that the exterior angle \drawAngleAlt{ACD} is equal to the two interior and opposite angles \drawAngleAlt{CAB}, \drawAngleAlt{ABC}, and the three interior angles of the triangle, \drawAngleAlt{ABC}, \drawAngleAlt{BCA}, \drawAngleAlt{CAB} are equal to two right angles. 
    \vskip \baselineskip
    For let \drawUnitLineAlt{CE} be drawn through the point \drawPointAlt{C} parallel to the straight line \drawUnitLineAlt{AB}. \par\justification{I. 31}
    Then, since \drawUnitLineAlt{AB} is parallel to \drawUnitLineAlt{CE}, and \drawUnitLineAlt{AC} has fallen upon them, \par
    \centeredjust{the alternate angles \drawAngleAlt{BAC}, \drawAngleAlt{ACE} are equal to one another.}{I. 29}
    Again, since \drawUnitLineAlt{AB} is parallel to \drawUnitLineAlt{CE}, and the straight line \drawUnitLineAlt{BD} has fallen upon them, \par
    \centeredjust{the exterior angle \drawAngleAlt{ECD} is equal to the interior and opposite angle \drawAngleAlt{ABC}.}{I. 29}
    But the angle \drawAngleAlt{ACE} was also proved equal to the angle \drawAngleAlt{BAC}; \par
    {\centering therefore the whole angle \drawAngleAlt{ACD} is equal to the two interior and opposite angles \drawAngleAlt{BAC}, \drawAngleAlt{ABC}. \par}
    Let the angle \drawAngleAlt{ACB} be added to each; \par
    {\centering therefore the angles \drawAngleAlt{ACD}, \drawAngleAlt{ACB} are equal to the three angles \drawAngleAlt{ABC}, \drawAngleAlt{BCA}, \drawAngleAlt{CAB}. \par}
    But the angles \drawAngleAlt{ACD}, \drawAngleAlt{ACB} are equal to two right angles; \justification{I. 13}
    {\centering therefore the angles \drawAngleAlt{ABC}, \drawAngleAlt{BCA}, \drawAngleAlt{CAB} are also equal to two right angles. \par}
    Therefore etc. 
}{}

\iftoggle{raum}{
    Raum_Text
}{}
\nolinenumbers
\hfill Q. E. D.

\resetlinenumber
}

\resetlinenumber

\startProposition{
    \markboth{Euclid in Color}{Proposition I.33}
    \defineNewPicture[1/4][3]{
	pair A, B, C, D, d[];
	d1 := (-5/2u, 0);
	d2 := (-7/8u, -3u);
	A := (0, 0);
	B := A shifted d1;
	C := A shifted d2;
	D := C shifted d1;
	byAngleDefine(A, B, C, byyellow, SOLID_SECTOR);
	byAngleDefine(A, C, B, byred, SOLID_SECTOR);
	byAngleDefine(B, C, D, byblack, SOLID_SECTOR);
	byAngleDefine(C, B, D, byblue, SOLID_SECTOR);
	draw byNamedAngleResized();
	draw byLine(B, C, byblack, SOLID_LINE, REGULAR_WIDTH);
	byLineDefine(A, B, byred, SOLID_LINE, REGULAR_WIDTH);
	byLineDefine(C, D, bygreen, SOLID_LINE, REGULAR_WIDTH);
	byLineDefine(A, C, byblue, SOLID_LINE, REGULAR_WIDTH);
	byLineDefine(B, D, byyellow, SOLID_LINE, REGULAR_WIDTH);
	draw byNamedLineSeq(0)(AB,BD,CD,AC);
	draw byLabelsOnPolygon(B, A, C, D)(ALL_LABELS, 0);
    }
    ~\hfill\drawCurrentPicture\hfill~
}{
    \chapter*{Proposition I.33}
    \renewcommand{\currentpropheader}{Proposition I.33}
    \addcontentsline{toc}{chapter}{Proposition 33}
    \textit{The straight lines joining equal and parallel straight lines (at the extremities which are) in the same directions (respectively) are themselves also equal and parallel.}
    \vskip \baselineskip

    % Check line numbering toggle
\iftoggle{lineNumbering}{\linenumbers}{}

\iftoggle{heath}{
    Let \drawUnitLineAlt{AB},\drawUnitLineAlt{CD} be equal and parallel, and let the straight lines \drawUnitLineAlt{AC},\drawUnitLineAlt{BD} join them (at the extremities which are) in the same directions (respectively); \par
    I say that \drawUnitLineAlt{AC},\drawUnitLineAlt{BD} are also equal and parallel. 
    \vskip \baselineskip
    Let \drawUnitLineAlt{BC} be joined. \par
    Then, since \drawUnitLineAlt{AB} is parallel to \drawUnitLineAlt{CD}, and \drawUnitLineAlt{BC} has fallen upon them, \par
    \centeredjust{the alternate angles \drawAngleAlt{ABC}, \drawAngleAlt{BCD} are equal to one another.}{I. 29}
    And, since \drawUnitLineAlt{AB} is equal to \drawUnitLineAlt{CD}, and \drawUnitLineAlt{BC} is common, \par
    the two sides \drawUnitLineAlt{AB},\drawUnitLineAlt{BC} are equal to the two sides \drawUnitLineAlt{DC},\drawUnitLineAlt{CB}; \hfill~\par
    and the angle \drawAngleAlt{ABC} is equal to the angle \drawAngleAlt{BCD}; \par
    {\centering therefore the base \drawUnitLineAlt{AC} is equal to the base \drawUnitLineAlt{BD}, \par}
    {\centering and the triangle \drawLineAlt{A,B,C} is equal to the triangle \drawLineAlt{D,C,B}, \par}
    and the remaining angles will be equal to the remaining angles respectively, namely those which the equal sides subtend;\justification{I. 4}
    {\centering therefore the angle \drawAngleAlt{ACB} is equal to the angle \drawAngleAlt{CBD}. \par}
    And, since the straight line \drawUnitLineAlt{BC} falling on the two straight lines \drawUnitLineAlt{AC},\drawUnitLineAlt{BD} has made the alternate angles equal to one another,\par
    \centeredjust{ \drawUnitLineAlt{AC} is parallel to \drawUnitLineAlt{BD}.}{I. 27}
    And it was also proved equal to it. \par
    Therefore etc.
}{}

\iftoggle{raum}{
    Raum_Text
}{}
\nolinenumbers
\hfill Q. E. D.

\resetlinenumber
}

\resetlinenumber

\startProposition{
    \markboth{Euclid in Color}{Proposition I.34}
    \defineNewPicture[1/6][2]{
	pair A, B, C, D, d[];
	d1 := (9/2u, 0);
	d2 := (-7/8u, -5/2u);
	A := (0, 0);
	B := A shifted d1;
	C := A shifted d2;
	D := C shifted d1;
	byAngleDefine(A, B, C, byblue, SOLID_SECTOR);
	byAngleDefine(B, C, D, byblue, SOLID_SECTOR);
	byAngleDefine(A, C, B, bygreen, SOLID_SECTOR);
	byAngleDefine(C, B, D, byred, SOLID_SECTOR);
	byAngleDefine(C, D, B, byblack, SOLID_SECTOR);
	byAngleDefine(B, A, C, byyellow, SOLID_SECTOR);
	draw byNamedAngleResized();
	draw byLine(B, C, byblack, SOLID_LINE, REGULAR_WIDTH);
	byLineDefine(A, B, bygreen, SOLID_LINE, REGULAR_WIDTH);
	byLineDefine(C, D, byred, SOLID_LINE, REGULAR_WIDTH);
	byLineDefine(A, C, byblue, SOLID_LINE, REGULAR_WIDTH);
	byLineDefine(B, D, byyellow, SOLID_LINE, REGULAR_WIDTH);
	draw byNamedLineSeq(0)(AB,BD,CD,AC);
	draw byLabelsOnPolygon(A, B, D, C)(ALL_LABELS, 0);
    }
    ~\hfill\drawCurrentPicture\hfill~
}{
    \chapter*{Proposition I.34}
    \renewcommand{\currentpropheader}{Proposition I.34}
    \addcontentsline{toc}{chapter}{Proposition 34}
    \textit{In parallelogrammic areas the opposite sides and angles are equal to one another, and the diameter bisects the areas.}
    \vskip \baselineskip

    % Check line numbering toggle
\iftoggle{lineNumbering}{\linenumbers}{}

\iftoggle{heath}{
    Let \drawLineAlt{A,C,D,B} be a parallelogrammic area, and \drawUnitLineAlt{BC} its diameter; \par
    I say that the opposite sides and angles of the parallelogram \drawLineAlt{A,C,D,B} are equal to one another, and the diameter \drawUnitLineAlt{BC} bisects it. 
    \vskip \baselineskip
    For, since \drawUnitLineAlt{AB} is parallel to \drawUnitLineAlt{CD}, and the straight line \drawUnitLineAlt{BC} has fallen upon them, \par
    \centeredjust{the alternate angles \drawAngleAlt{ABC}, \drawAngleAlt{BCD} are equal to one another.}{I. 29}
    Again, since \drawUnitLineAlt{AC} is parallel to \drawUnitLineAlt{BD}, and \drawUnitLineAlt{BC} has fallen upon them, \par
    \centeredjust{the alternate angles \drawAngleAlt{ACB}, \drawAngleAlt{CBD} are equal to one another.}{I. 29}
    Therefore \drawLineAlt{A,B,C}, \drawLineAlt{D,C,B} are two triangles having the two angles \drawAngleAlt{ABC}, \drawAngleAlt{BCA} equal to the two angles \drawAngleAlt{DCB}, \drawAngleAlt{CBD} respectively, and one side equal to one side, namely that adjoining the equal angles and common to both of them,\drawUnitLineAlt{BC}; \par
    therefore they will also have the remaining sides equal to the remaining sides respectively, and the remaining angle to the remaining angle;\justification{I. 26}
    {\centering therefore the side \drawUnitLineAlt{AB} is equal to \drawUnitLineAlt{CD}, \par}
    {\centering and \drawUnitLineAlt{AC} to \drawUnitLineAlt{BD}, \par}
    {\centering and further the angle \drawAngleAlt{BAC} is equal to the angle \drawAngleAlt{CDB}. \par}
    And, since the angle \drawAngleAlt{ABC} is equal to the angle \drawAngleAlt{BCD}, \par
    and the angle \drawAngleAlt{CBD} to the angle \drawAngleAlt{ACB}, \par
    \centeredjust{the whole angle \drawAngleAlt{ABD} is equal to the whole angle \drawAngleAlt{ACD}.\hspace{1.2em}}{C. N. 2}
    And the angle \drawAngleAlt{BAC} was also proved equal to the angle \drawAngleAlt{CDB}. \par
    Therefore in parallelogrammic areas the opposite sides and angles are equal to one another. \par
    I say, next, that the diameter also bisects the areas.\par
    For, since \drawUnitLineAlt{AB} is equal to \drawUnitLineAlt{CD}, and \drawUnitLineAlt{BC} is common, \par
    {\centering the two sides \drawUnitLineAlt{AB},\drawUnitLineAlt{BC} are equal to the two sides \drawUnitLineAlt{DC},\drawUnitLineAlt{CB} respectively; \par}
    and the angle \drawAngleAlt{ABC} is equal to the angle \drawAngleAlt{BCD}; \par
    {\centering therefore the base \drawUnitLineAlt{AC} is also equal to \drawUnitLineAlt{DB}, \par}
    \centeredjust{and the triangle \drawLineAlt{A,B,C} is equal to the triangle \drawLineAlt{D,C,B}.}{I. 4}
    Therefore the diameter \drawUnitLineAlt{BC} bisects the parallelogram \drawLineAlt{A,C,D,B}. 
}{}

\iftoggle{raum}{
    Raum_Text
}{}
\nolinenumbers
\hfill Q. E. D.

\vskip \baselineskip


\resetlinenumber
}

\resetlinenumber

\startProposition{
    \markboth{Euclid in Color}{Proposition I.35}
    \defineNewPicture[1/4][2]{
	pair A, B, C, D, E, F, G, d[];
	d1 := (9/4u, 0);
	d2 := (u, -3u);
	d3 := (-5/2u, -3u);
	A := (0, 0);
	D := A shifted d1;
	B := A shifted d2;
	C := B shifted d1;
	E := B shifted -d3;
	F := C shifted -d3;
	G := (D--C) intersectionpoint (B--E);
	draw byPolygon(A,B,G,D)(byyellow); % Keep solid polygon colors?
	draw byPolygon(E,G,C,F)(byyellow); %
	draw byPolygon(D,G,E)(byyellow); %
	draw byPolygon(B,C,G)(byred); %
	byAngleDefine(E, A, B, byblue, SOLID_SECTOR);
	byAngleDefine(F, D, C, byred, SOLID_SECTOR);
	byAngleDefine(A, E, B, byblack, SOLID_SECTOR);
	byAngleDefine(D, F, C, bygreen, SOLID_SECTOR);
	draw byNamedAngleResized();
	draw byNamedAngleDummySides(DFC);
	byLineDefine(A, B, byred, SOLID_LINE, REGULAR_WIDTH);
	byLineDefine(B, C, byblack, SOLID_LINE, REGULAR_WIDTH);
	byLineDefine(B, E, byblue, SOLID_LINE, REGULAR_WIDTH);
	byLineDefine(B, G, byblue, SOLID_LINE, REGULAR_WIDTH);
	byLineDefine(G, E, byblue, SOLID_LINE, REGULAR_WIDTH);
	byLineDefine(A, D, bygreen, SOLID_LINE, REGULAR_WIDTH);
	byLineDefine(D, E, byblue, SOLID_LINE, REGULAR_WIDTH);
	byLineDefine(E, F, byred, SOLID_LINE, REGULAR_WIDTH);
	byLineDefine(D, C, bygreen, SOLID_LINE, REGULAR_WIDTH);
	byLineDefine(D, G, bygreen, SOLID_LINE, REGULAR_WIDTH);
	byLineDefine(G, C, bygreen, SOLID_LINE, REGULAR_WIDTH);
	byLineDefine(C, F, byblue, SOLID_LINE, REGULAR_WIDTH);
	draw byNamedLineSeq(0)(AB,BC,DC,AD,DE,EF,BE,CF);
	draw byLabelsOnPolygon(B, A, D, E, F, C, G)(ALL_LABELS, 0);
    }
    ~\hfill\drawCurrentPicture\hfill~
}{
    \chapter*{Proposition I.35}
    \renewcommand{\currentpropheader}{Proposition I.35}
    \addcontentsline{toc}{chapter}{Proposition 35}
    \textit{Parallelograms which are on the same base and in the same parallels are equal to one another.}
    \vskip \baselineskip

    % Check line numbering toggle
\iftoggle{lineNumbering}{\linenumbers}{}

\iftoggle{heath}{
    Let \drawPolygonAlt[ABCD]{ABGD, BCG}, \drawPolygonAlt[EBCF]{EGCF, BCG} be parallelograms on the same base \drawUnitLineAlt{BC} and in the same parallels \drawUnitLineAlt{AF},\drawUnitLineAlt{BC}; \par
    I say that \drawPolygonAlt[ABCD]{ABGD, BCG} is equal to the parallelogram \drawPolygonAlt[EBCF]{EGCF, BCG}. 
    \vskip \baselineskip
    For, since \drawPolygonAlt[ABCD]{ABGD, BCG} is a parallelogram,\par
    \centeredjust{ \drawUnitLineAlt{AD} is equal to \drawUnitLineAlt{BC}.}{I. 34}
    For the same reason also\par
    {\centering  \drawUnitLineAlt{EF} is equal to \drawUnitLineAlt{BC}, \par}
    \centeredjust{so that \drawUnitLineAlt{AD} is also equal to \drawUnitLineAlt{EF};}{C. N. 1}
    and \drawUnitLineAlt{DE} is common; \par
    \centeredjust{therefore the whole \drawUnitLineAlt{AE} is equal to the whole \drawUnitLineAlt{DF}.}{C. N. 2}
    But \drawUnitLineAlt{AB} is also equal to \drawUnitLineAlt{DC}; \justification{I. 34}
    therefore the two sides \drawUnitLineAlt{EA},\drawUnitLineAlt{AB} are equal to the two sides \drawUnitLineAlt{FD},\drawUnitLineAlt{DC} respectively, and the angle \drawAngleAlt{FDC} is equal to the angle \drawAngleAlt{EAB}, the exterior to the interior; \par\justification{I. 29}
    {\centering therefore the base \drawUnitLineAlt{EB} is equal to the base \drawUnitLineAlt{FC}, \par}
    \centeredjust{and the triangle \hspace{-0.25em}\drawPolygonAlt[EAB]{ABGD, DGE} will be equal to the triangle \hspace{-0.25em}\drawPolygonAlt[FDC]{EGCF, DGE}.}{I. 4}
    Let \drawPolygonAlt{DGE} be subtracted from each; \par
    therefore the trapezium \drawPolygonAlt{ABGD} which remains is equal to the trapezium \par \drawPolygonAlt{EGCF} which remains.\justification{C. N. 3}
    Let the triangle \drawPolygonAlt[GCB]{BCG} be added to each; \par
    {\centering therefore the whole parallelogram \drawPolygonAlt[ABCD]{ABGD, BCG} is equal to the whole \par}
    \centeredjust{parallelogram \drawPolygonAlt[EBCF]{EGCF, BCG}.}{C. N. 2}
    Therefore etc. 
}{}

\iftoggle{raum}{
    Raum_Text
}{}
\nolinenumbers
\hfill Q. E. D.

\vskip \baselineskip


\resetlinenumber
}

\resetlinenumber

\startProposition{
    \markboth{Euclid in Color}{Proposition I.36}
    \defineNewPicture[1/3][7/4]{
	pair A, D, B, C, E, H, F, G, J, I, d[];
	numeric h;
	h := 5/2u;
	d1 := (2u, 0);
	d2 := (-1/3u, -h);
	d3 := (4u, h);
	d4 := (2/3u, -h);
	A := (0, 0);
	D := A shifted d1;
	B := A shifted d2;
	C := B shifted d1;
	E := B shifted d3;
	H := C shifted d3;
	F := E shifted d4;
	G := H shifted d4;
	I := (D--C) intersectionpoint (B--E);
	J := (E--F) intersectionpoint (C--H);
	draw byPolygon(A,D,I,B)(byred);
	draw byPolygon(B,C,I)(byred);
	draw byPolygon(I,C,J,E)(byblue);
	draw byPolygon(E,H,J)(byyellow);
	draw byPolygon(J,H,G,F)(byyellow);
	byLineDefine(B, E, byyellow, SOLID_LINE, REGULAR_WIDTH);
	byLineDefine(C, H, bygreen, SOLID_LINE, REGULAR_WIDTH);
	byLineDefine(B, C, byblue, SOLID_LINE, REGULAR_WIDTH);
	byLineDefine(E, H, byred, SOLID_LINE, REGULAR_WIDTH);
	byLineDefine(F, G, byblue, SOLID_LINE, REGULAR_WIDTH);
	byLineDefine(A, D, bygreen, SOLID_LINE, REGULAR_WIDTH);
	byLineDefine(C, F, byyellow, SOLID_LINE, REGULAR_WIDTH);
	byLineDefine(D, E, byblue, SOLID_LINE, REGULAR_WIDTH);
	byLineDefine(A, B, byblack, SOLID_LINE, REGULAR_WIDTH);
	byLineDefine(D, C, byblack, SOLID_LINE, REGULAR_WIDTH);
	byLineDefine(E, F, byblack, SOLID_LINE, REGULAR_WIDTH);
	byLineDefine(H, G, byblack, SOLID_LINE, REGULAR_WIDTH);
	draw byNamedLineSeq(0)(DC,AD,AB,BC,CH,BE,CF,DE,EF,FG,HG,EH);
	%draw byLineFull(F, G, byblue, 0, 0)(E, H, 1, 1, 0);
	draw byLabelsOnPolygon(A, D, C, B)(ALL_LABELS, 0);
	draw byLabelsOnPolygon(E, H, G, F)(ALL_LABELS, 0);
    }
    ~\hfill\drawCurrentPicture\hfill~
}{
    \chapter*{Proposition I.36}
    \renewcommand{\currentpropheader}{Proposition I.36}
    \addcontentsline{toc}{chapter}{Proposition 36}
    \textit{Parallelograms which are on equal bases and in the same parallels are equal to one another.}
    \vskip \baselineskip

    % Check line numbering toggle
\iftoggle{lineNumbering}{\linenumbers}{}

\iftoggle{heath}{
    Let \drawPolygonAlt[ABCD]{ADIB, BCI}, \drawPolygonAlt[EFGH]{EHJ, JHGF} be parallelograms which are on equal bases \drawUnitLineAlt{BC},\drawUnitLineAlt{FG} and in the same parallels \drawUnitLineAlt{AH},\drawUnitLineAlt{BG}; \par
    I say that the parallelogram \drawPolygonAlt[ABCD]{ADIB, BCI} is equal to \drawPolygonAlt[EFGH]{EHJ, JHGF}. 
    \vskip \baselineskip
    For let \drawUnitLineAlt{BE},\drawUnitLineAlt{CH} be joined. \par
    Then, since \drawUnitLineAlt{BC} is equal to \drawUnitLineAlt{FG} while \drawUnitLineAlt{FG} is equal to \drawUnitLineAlt{EH},\par
    \centeredjust{ \drawUnitLineAlt{BC} is also equal to \drawUnitLineAlt{EH}.}{C. N. 1}
    But they are also parallel. \par
    And \drawUnitLineAlt{EB},\drawUnitLineAlt{HC} join them; \par
    but straight lines joining equal and parallel straight lines [at the extremities which are] in the same directions [respectively] are equal and parallel.\footnote{\noindent Square brackets are used to indicate content that Heath himself added, which was not literally in the Greek text.} \justification{I. 33}
    \centeredjust{Therefore \drawPolygonAlt[EBCH]{BCI,EHJ, ICJE} is a parallelogram.}{I. 34}
    And it is equal to \drawPolygonAlt[ABCD]{ADIB, BCI}; \par
    for it has the same base \drawUnitLineAlt{BC} with it, and is in the same parallels \drawUnitLineAlt{BC},\drawUnitLineAlt{AH} with it. \justification{I. 35}
    For the same reason also \drawPolygonAlt[EFGH]{EHJ, JHGF} is equal to the same \drawPolygonAlt[EBCH]{BCI,EHJ, ICJE}; \par\justification{I. 35}
    \centeredjust{so that the parallelogram \drawPolygonAlt[ABCD]{ADIB, BCI} is also equal to \drawPolygonAlt[EFGH]{EHJ, JHGF}.}{C. N. 1}
    Therefore etc. 
}{}

\iftoggle{raum}{
    Raum_Text
}{}
\nolinenumbers
\hfill Q. E. D.

\resetlinenumber
}

\resetlinenumber

\startProposition{
    \markboth{Euclid in Color}{Proposition I.37}
    \defineNewPicture[1/2][9/4]{
	pair E, A, B, C, D, F, G, H, I, d[];
	d1 := (5/2u, 0);
	d2 := (u, -5/2u);
	d3 := (-2u, -5/2u);
	E := (0, 0);
	A := E shifted d1;
	B := E shifted d2;
	C := B shifted d1;
	D := B shifted -d3;
	F := C shifted -d3;
	G := (A--C) intersectionpoint (B--D);
	draw byPolygon(E,A,B)(byblue);
	draw byPolygon(A,B,G)(byyellow);
	draw byPolygon(B,C,G)(byyellow);
	draw byPolygon(C,G,D)(byblack);
	draw byPolygon(D,F,C)(byred);
	draw byLine(A, C, byred, SOLID_LINE, REGULAR_WIDTH);
	draw byLine(D, B, byblue, SOLID_LINE, REGULAR_WIDTH);
	byLineDefine(E, B, byyellow, SOLID_LINE, REGULAR_WIDTH);
	byLineDefine(F, C, bygreen, SOLID_LINE, REGULAR_WIDTH);
	byLineDefine(B, C, byblack, SOLID_LINE, REGULAR_WIDTH);
	byLineDefine(D, C, byyellow, SOLID_LINE, REGULAR_WIDTH);
	byLineDefine(A, B, bygreen, SOLID_LINE, REGULAR_WIDTH);
	byLineDefine(E, F, byblack, SOLID_LINE, REGULAR_WIDTH);
	byLineDefine(E, A, byblack, SOLID_LINE, REGULAR_WIDTH);
	byLineDefine(A, D, byblack, SOLID_LINE, REGULAR_WIDTH);
	byLineDefine(D, F, byblack, SOLID_LINE, REGULAR_WIDTH);
	draw byNamedLineSeq(0)(EB,BC,FC,DC,AB,EF);
	draw byLabelsOnPolygon(E, A, D, F, noPoint)(ALL_LABELS, 0);
	draw byLabelsOnPolygon(F, C, B, E)(OMIT_FIRST_LABEL+OMIT_LAST_LABEL, 0);
    }
    ~\hfill\drawCurrentPicture\hfill~
}{
    \chapter*{Proposition I.37}
    \renewcommand{\currentpropheader}{Proposition I.37}
    \addcontentsline{toc}{chapter}{Proposition 37}
    \textit{Triangles which are on the same base and in the same parallels are equal to one another.}
    \vskip \baselineskip

    % Check line numbering toggle
\iftoggle{lineNumbering}{\linenumbers}{}

\iftoggle{heath}{
    Let \drawPolygonAlt[ABC]{ABG, BCG}, \drawPolygonAlt[DBC]{CGD, BCG} be triangles on the same base \drawUnitLineAlt{BC} and in the same parallels \drawUnitLineAlt{AD},\drawUnitLineAlt{BC}; \par
    I say that the triangle \drawPolygonAlt[ABC]{ABG, BCG} is equal to the triangle \drawPolygonAlt[DBC]{CGD, BCG}. 
    \vskip \baselineskip
    Let \drawUnitLineAlt{AD} be produced in both directions to \drawPointAlt{E}, \drawPointAlt{F}; \par
    through \drawPointAlt{B} let \drawUnitLineAlt{BE} be drawn parallel to \drawUnitLineAlt{CA}, \justification{I. 31}
    and through \drawPointAlt{C} let \drawUnitLineAlt{CF} be drawn parallel to \drawUnitLineAlt{BD}. \justification{I. 31}
    Then each of the figures \drawPolygonAlt[EBCA]{EAB,ABG, BCG}, \drawPolygonAlt[DBCF]{DFC,CGD, BCG} is a parallelogram; and they are equal, for they are on the same base \drawUnitLineAlt{BC} and in the same parallels \drawUnitLineAlt{BC},\drawUnitLineAlt{EF}. \justification{I. 35}
    Moreover the triangle \drawPolygonAlt[ABC]{ABG, BCG} is half of the parallelogram \drawPolygonAlt[EBCA]{EAB,ABG, BCG}; for the diameter \drawUnitLineAlt{AB} bisects it. \justification{I. 34}
    And the triangle \drawPolygonAlt[DBC]{CGD, BCG} is half of the parallelogram \drawPolygonAlt[DBCF]{DFC,CGD, BCG}; for the diameter \drawUnitLineAlt{DC} bisects it. \justification{I. 34}
    $\langle{}$But the halves of equal things are equal to one another.$\rangle{}$ \par
    {\centering Therefore the triangle \drawPolygonAlt[ABC]{ABG, BCG} is equal to the triangle \drawPolygonAlt[DBC]{CGD, BCG}. \par}
    Therefore etc. 
}{}

\iftoggle{raum}{
    Raum_Text
}{}
\nolinenumbers
\hfill Q. E. D.

\resetlinenumber
}

\resetlinenumber

\startProposition{
    \markboth{Euclid in Color}{Proposition I.38}
    \defineNewPicture[1/3][5/4]{
	pair G, A, B, C, D, H, E, F, J, I, d[];
	numeric h;
	h := 3u;
	d1 := (3u, 0);
	d2 := (-1/2u, -h);
	d3 := (9/2u, h);
	d4 := (3/2u, -h);
	G := (0, 0);
	A := G shifted d1;
	B := G shifted d2;
	C := B shifted d1;
	D := B shifted d3;
	H := C shifted d3;
	E := D shifted d4;
	F := H shifted d4;
	draw byPolygon(G,A,B)(byyellow);
	draw byPolygon(A,B,C)(byred);
	draw byPolygon(D,H,F)(bygreen);
	draw byPolygon(D,E,F)(byblue);
	draw byLine(A, C, byblue, SOLID_LINE, REGULAR_WIDTH);
	draw byLine(D, E, byred, SOLID_LINE, REGULAR_WIDTH);
	draw byLine(A, B, bygreen, SOLID_LINE, REGULAR_WIDTH);
	draw byLine(D, F, byyellow, SOLID_LINE, REGULAR_WIDTH);
	byLineDefine(G, B, byred, SOLID_LINE, REGULAR_WIDTH);
	byLineDefine(H, F, byblue, SOLID_LINE, REGULAR_WIDTH);
	byLineDefine(G, H, byblack, SOLID_LINE, REGULAR_WIDTH);
	byLineDefine(G, A, byblack, SOLID_LINE, REGULAR_WIDTH);
	byLineDefine(A, D, byblack, SOLID_LINE, REGULAR_WIDTH);
	byLineDefine(D, H, byblack, SOLID_LINE, REGULAR_WIDTH);
	byLineDefine(B, F, byblack, SOLID_LINE, REGULAR_WIDTH);
	byLineDefine(B, C, byblack, SOLID_LINE, REGULAR_WIDTH);
	byLineDefine(C, E, byblack, SOLID_LINE, REGULAR_WIDTH);
	byLineDefine(E, F, byblack, SOLID_LINE, REGULAR_WIDTH);
	draw byNamedLineSeq(0)(GB,GH,HF,BF);
	draw byLabelsOnPolygon(G, A, D, H, noPoint)(ALL_LABELS, 0);
	draw byLabelsOnPolygon(F, E, C, B, noPoint)(ALL_LABELS, 0);
    }
    ~\hfill\drawCurrentPicture\hfill~
}{
    \chapter*{Proposition I.38}
    \renewcommand{\currentpropheader}{Proposition I.38}
    \addcontentsline{toc}{chapter}{Proposition 38}
    \textit{Triangles which are on equal bases and in the same parallels are equal to one another.}
    \vskip \baselineskip

    % Check line numbering toggle
\iftoggle{lineNumbering}{\linenumbers}{}

\iftoggle{heath}{
    Let \drawPolygonAlt{ABC}, \drawPolygonAlt{DEF} be triangles on equal bases \drawUnitLineAlt{BC},\drawUnitLineAlt{EF} and in the same parallels \drawUnitLineAlt{BF},\drawUnitLineAlt{AD}; \par
    I say that the triangle \drawPolygonAlt{ABC} is equal to the triangle \drawPolygonAlt{DEF}. 
    \vskip \baselineskip
    For let \drawUnitLineAlt{AD} be produced in both directions to \drawPointAlt{G}, \drawPointAlt{H}; \par
    through \drawPointAlt{B} let \drawUnitLineAlt{BG} be drawn parallel to \drawUnitLineAlt{CA}, \justification{I. 31}
    and through \drawPointAlt{F} let \drawUnitLineAlt{FH} be drawn parallel to \drawUnitLineAlt{DE}. \par
    {\centering Then each of the figures \drawPolygonAlt[GBCA]{ABC,GAB}, \drawPolygonAlt[DEFH]{DEF,DHF} is a parallelogram; \par}
    and \drawPolygonAlt[GBCA]{ABC,GAB} is equal to \drawPolygonAlt[DEFH]{DEF,DHF}; \par
    for they are on equal bases \drawUnitLineAlt{BC},\drawUnitLineAlt{EF} and in the same parallels \drawUnitLineAlt{BF},\drawUnitLineAlt{GH}. \justification{I. 36}
    Moreover the triangle \drawPolygonAlt{ABC} is half of the parallelogram \drawPolygonAlt[GBCA]{ABC,GAB}; for the diameter \drawUnitLineAlt{AB} bisects it. \justification{I. 34}
    And the triangle \drawPolygonAlt[FED]{DEF} is half of the parallelogram \drawPolygonAlt[DEFH]{DEF,DHF}; \par
    for the diameter \drawUnitLineAlt{DF} bisects it. \justification{I. 34}
    $\langle$But the halves of equal things are equal to one another.$\rangle$ \par
    {\centering Therefore the triangle \drawPolygonAlt{ABC} is equal to the triangle \drawPolygonAlt{DEF}. \par}
    Therefore etc. 
}{}

\iftoggle{raum}{
    Raum_Text
}{}
\nolinenumbers
\hfill Q. E. D.

\resetlinenumber
}

\resetlinenumber

\startProposition{
    \markboth{Euclid in Color}{Proposition I.39}
    \defineNewPicture[1/3][5/2]{
	pair A, B, C, D, G, E;
	A := (0, 0);
	D := A shifted (3u, 0);
	B := A shifted (-1/2u, -5/2u);
	C := B shifted (9/2u, 0);
	E := 4/5[B, D];
	G = whatever[A, C] = whatever[D, B];
	draw byPolygon(A,D,G)(byred);
	draw byPolygon(A,B,G)(byyellow);
	draw byPolygon(E,G,C)(byblue);
	draw byPolygon(G,B,C)(byyellow);
	draw byPolygon(D,E,C)(byblack);
	byLineDefine(A, E, byblue, SOLID_LINE, REGULAR_WIDTH);
	byLineDefine(C, E, byyellow, SOLID_LINE, REGULAR_WIDTH);
	byLineDefine(A, D, byblue, SOLID_LINE, REGULAR_WIDTH);
	byLineDefine(B, C, byblack, SOLID_LINE, REGULAR_WIDTH);
	byLineDefine(B, D, bygreen, SOLID_LINE, REGULAR_WIDTH);
	byLineDefine(D, C, byred, SOLID_LINE, REGULAR_WIDTH);
	byLineDefine(A, B, byblack, SOLID_LINE, REGULAR_WIDTH);	
	byLineDefine(A, C, byblack, SOLID_LINE, REGULAR_WIDTH);	
	draw byNamedLineSeq(-4/5)(AC,AE,BD,CE);
	draw byNamedLineSeq(0)(AD, AB, BC, DC);
	draw byLabelsOnPolygon(E, C, B, A)(OMIT_FIRST_LABEL+OMIT_LAST_LABEL, 0);
	draw byLabelsOnPolygon(A, D, E)(OMIT_FIRST_LABEL+OMIT_LAST_LABEL, 0);
	draw byLabelsOnPolygon(A, E, noPoint)(ALL_LABELS, 0);
    }
    ~\hfill\drawCurrentPicture\hfill~
}{
    \chapter*{Proposition I.39}
    \renewcommand{\currentpropheader}{Proposition I.39}
    \addcontentsline{toc}{chapter}{Proposition 39}
    \textit{Equal triangles which are on the same base and on the same side are also in the same parallels.}
    \vskip \baselineskip

    % Check line numbering toggle
\iftoggle{lineNumbering}{\linenumbers}{}

\iftoggle{heath}{
    Let \drawPolygonAlt[ABC]{ABG, GBC}, \drawPolygonAlt[DBC]{DEC, EGC, GBC} be equal triangles which are on the same base \drawUnitLineAlt{BC} and on the same side of it; \par
    $\langle$I say that they are also in the same parallels.$\rangle$\par
    And $\langle$For$\rangle$ let \drawUnitLineAlt{AD} be joined; \par
    I say that \drawUnitLineAlt{AD} is parallel to \drawUnitLineAlt{BC}.
    \vskip \baselineskip
    For, if not, let \drawUnitLineAlt{AE} be drawn through the point \drawPointAlt{A} parallel to the straight line \drawUnitLineAlt{BC}, \justification{I. 31}
    and let \drawUnitLineAlt{EC} be joined. \par
    {\centering Therefore the triangle \drawPolygonAlt[ABC]{ABG, GBC} is equal to the triangle \drawPolygonAlt[EBC]{EGC, GBC}; \par}
    for it is on the same base \drawUnitLineAlt{BC} with it and in the same parallels. \justification{I. 37}
    But \drawPolygonAlt[ABC]{ABG, GBC} is equal to \drawPolygonAlt[DBC]{DEC, EGC, GBC}; \par
    \centeredjust{therefore \drawPolygonAlt[DBC]{DEC, EGC, GBC} is also equal to \drawPolygonAlt[EBC]{EGC, GBC},}{C. N. 1}
    the greater to the less: which is impossible. \par
    {\centering Therefore \drawUnitLineAlt{AE} is not parallel to \drawUnitLineAlt{BC}. \par}
    Similarly we can prove that neither is any other straight line except \drawUnitLineAlt{AD}; \par
    {\centering therefore \drawUnitLineAlt{AD} is parallel to \drawUnitLineAlt{BC}. \par}
    Therefore etc. 
}{}

\iftoggle{raum}{
    Raum_Text
}{}
\nolinenumbers
\hfill Q. E. D.

\resetlinenumber
}

\resetlinenumber

\startProposition{
    \markboth{Euclid in Color}{Proposition I.40}
    \defineNewPicture[1/3][2]{
	pair A, B, C, D, E, F, d;
	A := (0, 0);
	D := A shifted (5u, 0);
	B := A shifted (-1/4u, -5/2u);
	d := (3u, 0);
	C := B shifted d;
	E := C shifted d;
	F := 13/16[C, D];
	draw byPolygon(A,B,C)(byyellow);
	draw byPolygon(F,C,E)(byred);
	draw byPolygon(F,D,E)(byblue);
	byLineDefine(C, D, bygreen, SOLID_LINE, REGULAR_WIDTH);
	byLineDefine(A, F, byred, SOLID_LINE, REGULAR_WIDTH);
	byLineDefine(E, F, byyellow, SOLID_LINE, REGULAR_WIDTH);
	byLineDefine(A, D, byblue, SOLID_LINE, REGULAR_WIDTH);
	byLineDefine(C, E, byblack, SOLID_LINE, REGULAR_WIDTH);
	byLineDefine(B, C, byblack, SOLID_LINE, REGULAR_WIDTH);
	byLineDefine(A, B, byblack, SOLID_LINE, REGULAR_WIDTH);
	byLineDefine(A, C, byblack, SOLID_LINE, REGULAR_WIDTH);
	byLineDefine(D, E, byblack, SOLID_LINE, REGULAR_WIDTH);
	draw byNamedLineSeq(0)(AC,AB,BC);
	draw byNamedLineSeq(0)(DE,CD,CE);
	draw byNamedLineSeq(0)(AD,AF,EF);
	draw byLabelsOnPolygon(E, C, B, A)(OMIT_LAST_LABEL, 0);
	draw byLabelsOnPolygon(A, D, F)(OMIT_FIRST_LABEL+OMIT_LAST_LABEL, 0);
	draw byLabelsOnPolygon(F, noPoint, A)(OMIT_FIRST_LABEL, 0);
	draw byLabelPoint(F, 30, -6);
    }
    ~\hfill\drawCurrentPicture\hfill~
}{
    \chapter*{$\langle$Proposition I.40}
    \addcontentsline{toc}{chapter}{Proposition 40}
    \textit{Equal triangles which are on equal bases and on the same side are also in the same parallels.}
    \vskip \baselineskip

    % Check line numbering toggle
\iftoggle{lineNumbering}{\linenumbers}{}

\iftoggle{heath}{
    Let \drawPolygonAlt{ABC}, \drawPolygonAlt[CDE]{FCE,FDE} be equal triangles on equal bases \drawUnitLineAlt{BC},\drawUnitLineAlt{CE} and on the same side. \par
    I say that they are also in the same parallels. For let \drawUnitLineAlt{AD} be joined; I say that \drawUnitLineAlt{AD} is parallel to \drawUnitLineAlt{BE}. 
    \vskip \baselineskip
    For, if not, let \drawUnitLineAlt{AF} be drawn through \drawPointAlt{A} parallel to \drawUnitLineAlt{BE}, \justification{I. 31}
    and let \drawUnitLineAlt{FE} be joined. \par
    {\centering Therefore the triangle \drawPolygonAlt{ABC} is equal to the triangle \drawPolygonAlt{FCE}; \par}
    for they are on equal bases \drawUnitLineAlt{BC},\drawUnitLineAlt{CE} and in the same parallels \drawUnitLineAlt{BE},\drawUnitLineAlt{AF}. \justification{I. 38}
    But the triangle \drawPolygonAlt{ABC} is equal to the triangle \drawPolygonAlt[DCE]{FCE,FDE};\par
    \centeredjust{therefore the triangle \drawPolygonAlt[DCE]{FCE,FDE} is also equal to the triangle \drawPolygonAlt{FCE},}{C. N. 1}
    the greater to the less: which is impossible. \par 
    {\centering Therefore \drawUnitLineAlt{AF} is not parallel to \drawUnitLineAlt{BE}. \par}
    Similarly we can prove that neither is any other straight line except \drawUnitLineAlt{AD};\par
    {\centering therefore \drawUnitLineAlt{AD} is parallel to \drawUnitLineAlt{BE}. \par}
    Therefore etc. 
}{}

\iftoggle{raum}{
    Raum_Text
}{}
\nolinenumbers
\hfill Q. E. D.$\rangle$ \footnote{\noindent Heath and others believe that proposition 40 was not originally part of Euclid's \textit{Elements}, but was added later by an editor.}

\resetlinenumber
}

\resetlinenumber

\startProposition{
    \markboth{Euclid in Color}{Proposition I.41}
    \defineNewPicture[1/3][5/2]{
	pair A, B, C, D, E, F, G, d;
	A := (0, 0);
	d := (3u, 0);
	D := A shifted d;
	E := D shifted (2u, 0);
	B := A shifted (1/2u, -2u);
	C := B shifted d;
	F = whatever[D, C] = whatever[B, E];
	G = whatever[A, C] = whatever[B, E];
	draw byPolygon(A,D,F,G)(byyellow);
	draw byPolygon(G,C,F)(byyellow);
	draw byPolygon(A,G,B)(byblue);
	draw byPolygon(B,C,G)(byblue);
	draw byPolygon(E,F,C)(byred);
	draw byLine(B, E, byblack, SOLID_LINE, THIN_WIDTH);
	draw byLine(A, C, byred, SOLID_LINE, REGULAR_WIDTH);
	draw byLine(C, E, byblack, SOLID_LINE, THIN_WIDTH);
	draw byLine(A, B, byblack, SOLID_LINE, THIN_WIDTH);
	draw byLine(A, E, bygreen, SOLID_LINE, REGULAR_WIDTH);
	draw byLine(B, C, byyellow, SOLID_LINE, REGULAR_WIDTH);
	draw byLine(D, C, byblack, SOLID_LINE, THIN_WIDTH);
	draw byLabelsOnPolygon(C, B, A, D, E)(ALL_LABELS, 0);
    }
    ~\hfill\drawCurrentPicture\hfill~
}{
    \chapter*{Proposition I.41}
    \renewcommand{\currentpropheader}{Proposition I.41}
    \addcontentsline{toc}{chapter}{Proposition 41}
    \textit{If a parallelogram have the same base with a triangle and be in the same parallels, the parallelogram is double of the triangle.}
    \vskip \baselineskip

    % Check line numbering toggle
\iftoggle{lineNumbering}{\linenumbers}{}

\iftoggle{heath}{
    For let the parallelogram \drawPolygonAlt[ABCD]{AGB,ADFG,GCF,BCG} have the same base \drawUnitLineAlt{BC} with the triangle \par \drawPolygonAlt[EBC]{EFC,GCF,BCG}, and let it be in the same parallels \drawUnitLineAlt{BC},\drawUnitLineAlt{AE}; \par
    I say that the parallelogram \drawPolygonAlt[ABCD]{AGB, BCG} is double of the triangle \drawPolygonAlt[EBC]{EFC,GCF,BCG}. 
    \vskip \baselineskip
    For let \drawUnitLineAlt{AC} be joined. \par
    {\centering Then the triangle \drawPolygonAlt[ABC]{AGB, BCG} is equal to the triangle \drawPolygonAlt[EBC]{EFC,GCF,BCG}; \par}
    for it is on the same base \drawUnitLineAlt{BC} with it and in the same parallels \drawUnitLineAlt{BC},\drawUnitLineAlt{AE}. \justification{I. 37}
    {\centering But the parallelogram \drawPolygonAlt[ABCD]{AGB,ADFG,GCF,BCG} is double of the triangle \drawPolygonAlt[ABC]{AGB, BCG}; \par}
    for the diameter \drawUnitLineAlt{AC} bisects it; \justification{I. 34}
    {\centering so that the parallelogram \drawPolygonAlt[ABCD]{AGB,ADFG,GCF,BCG} is also double of the triangle \drawPolygonAlt[EBC]{EFC,GCF,BCG}. \par}
    Therefore etc. 
}{}

\iftoggle{raum}{
    Raum_Text
}{}
\nolinenumbers
\hfill Q. E. D.

\resetlinenumber
}

\resetlinenumber

\startProposition{
    \markboth{Euclid in Color}{Proposition I.42}
    \defineNewPicture[1/3][3/2]{
	pair A, B, C, D, E, F, G, H, I, J, K, d[];
	numeric a, b;
	G := (0, 0);
	d1 := (2u, 0);
	F := G shifted -d1;
	A := F shifted (-7/2u, 0);
	C := G shifted (-5/2u, -5/2u);
	E := C shifted -d1;
	K := (F--E) intersectionpoint (C--A);
	B := 2[C, E];
	d2 := (-xpart(B)+xpart(E)-1/2u, 0);
	a := 40/abs(E-F);
	b := 40/abs(E-C);
	H := a[E,F] rotated 100 shifted (-10u, 3u);
	D := E rotated 100 shifted (-10u, 3u);
	J := b[E,C] rotated 100 shifted (-10u, 3u);
	draw byPolygon(G,F,K,C)(byyellow);
	draw byPolygon(C,E,K)(byyellow);
	draw byPolygon(A,K,E)(byblue);
	draw byPolygon(A,B,E)(byblack);
	byAngleDefine(F, E, C, bygreen, SOLID_SECTOR);
	byAngleDefine(H, D, J, byyellow, SOLID_SECTOR);
	setAttribute("angle", "Standalone", "HDJ", 1);
	draw byNamedAngleResized();
	draw byLine(A, B, byred, SOLID_LINE, REGULAR_WIDTH);
	draw byLine(C, A, byblack, SOLID_LINE, THIN_WIDTH);
	draw byLine(F, E, byred, SOLID_LINE, REGULAR_WIDTH);
	byLineDefine(G, C, bygreen, SOLID_LINE, REGULAR_WIDTH);
	byLineDefine(A, E, byyellow, SOLID_LINE, REGULAR_WIDTH);
	byLineDefine(G, A, byblue, SOLID_LINE, REGULAR_WIDTH);
	byLineDefine(C, E, byblack, SOLID_LINE, REGULAR_WIDTH);
	byLineDefine(E, B, bygreen, SOLID_LINE, REGULAR_WIDTH);
	draw byNamedLineSeq(0)(AE,GA,GC,CE,EB);
	byLineDefine(H, D, byblack, SOLID_LINE, REGULAR_WIDTH);
	byLineDefine(D, J, byblack, SOLID_LINE, REGULAR_WIDTH);
	draw byNamedLineSeq(0)(HD,DJ);
	draw byLabelsOnPolygon(A, F, G, C, E, B)(ALL_LABELS, 0);
	startAutoLabeling;
	draw byNamedAngleSidesFull(HDJ)();
	stopAutoLabeling;
    }
    ~\hfill\drawCurrentPicture\hfill~
}{
    \chapter*{Proposition I.42}
    \renewcommand{\currentpropheader}{Proposition I.42}
    \addcontentsline{toc}{chapter}{Proposition 42}
    \textit{To construct, in a given rectilineal angle, a parallelogram equal to a given triangle.}
    \vskip \baselineskip

    % Check line numbering toggle
\iftoggle{lineNumbering}{\linenumbers}{}

\iftoggle{heath}{
    Let \drawPolygonAlt[ABC]{ABE, AKE, CEK} be the given triangle, and \drawAngleAlt[D]{HDJ} the given rectilineal angle; \par
    thus it is required to construct in the rectilineal angle \drawAngleAlt[D]{HDJ} a parallelogram equal to the triangle \drawPolygonAlt[ABC]{ABE, AKE, CEK}. 
    \vskip \baselineskip
    Let \drawUnitLineAlt{BC} be bisected at \drawPointAlt{E}, and let \drawUnitLineAlt{AE} be joined; \par
    on the straight line \drawUnitLineAlt{EC}, and at the point \drawPointAlt{E} on it, let the angle \drawAngleAlt{CEF} be constructed equal to the angle \drawAngleAlt[D]{HDJ}; \justification{I. 23}
    through \drawPointAlt{A} let \drawUnitLineAlt{AG} be drawn parallel to \drawUnitLineAlt{EC}, \par
    and through \drawPointAlt{C} let \drawUnitLineAlt{CG} be drawn parallel to \drawUnitLineAlt{EF}. \justification{I. 31}
    {\centering Then \drawPolygonAlt[FECG]{GFKC, CEK} is a parallelogram. \par}
    And, since \drawUnitLineAlt{BE} is equal to \drawUnitLineAlt{EC}, \par
    {\centering the triangle \drawPolygonAlt{ABE} is also equal to the triangle \drawPolygonAlt[AEC]{AKE, CEK}, \par}
    for they are on equal bases \drawUnitLineAlt{BE},\drawUnitLineAlt{EC} and in the same parallels \drawUnitLineAlt{BC},\drawUnitLineAlt{AG}; \justification{I. 38}
    {\centering therefore the triangle \drawPolygonAlt[ABC]{ABE, AKE, CEK} is double of the triangle \drawPolygonAlt[AEC]{AKE, CEK}.\par}
    {\centering But the parallelogram \drawPolygonAlt[FECG]{GFKC, CEK} is also double of the triangle \drawPolygonAlt[AEC]{AKE, CEK}, \par}
    for it has the same base with it and is in the same parallels with it; \justification{I. 41}
    {\centering therefore the parallelogram \drawPolygonAlt[FECG]{GFKC, CEK} is equal to the triangle \drawPolygonAlt[ABC]{ABE, AKE, CEK}.\par}
    And it has the angle \drawAngleAlt{CEF} equal to the given angle \drawAngleAlt[D]{HDJ}. \par
    {\centering Therefore the parallelogram \drawPolygonAlt[FECG]{GFKC, CEK} has been constructed equal to the given \drawPolygonAlt[ABC]{ABE, AKE, CEK}, in the angle \drawAngleAlt{CEF} which is equal to \drawAngleAlt[D]{HDJ}. \par}
}{}

\iftoggle{raum}{
    Raum_Text
}{}
\nolinenumbers
\hfill Q. E. F.

\vskip \baselineskip


\resetlinenumber
}

\resetlinenumber

\startProposition{
    \markboth{Euclid in Color}{Proposition I.43}
    \defineNewPicture[1/3][2]{
	pair A, D, B, C, E, F, G, H, K, d[];
	path q[];
	d1 := (9/2u, 0);
	d2 := (-u, -3u);
	A := (0, 0);
	D := A shifted d1;
	B := A shifted d2;
	C := B shifted d1;
	K := 1/4[A, C];
	q1 := (K shifted d1) -- (K shifted -d1);
	q2 := (K shifted d2) -- (K shifted -d2);
	E := q1 intersectionpoint (A--B);
	F := q1 intersectionpoint (D--C);
	H := q2 intersectionpoint (A--D);
	G := q2 intersectionpoint (B--C);
	draw byPolygon(A,H,K)(byblue);
	draw byPolygon(A,E,K)(byblue);
	draw byPolygon(H,D,F,K)(bygreen);
	draw byPolygon(E,B,G,K)(byyellow);
	draw byPolygon(K,G,C)(byred);
	draw byPolygon(K,F,C)(byred);
	draw byLine(H, K, byblack, SOLID_LINE, REGULAR_WIDTH);
	draw byLine(E, K, byred, SOLID_LINE, REGULAR_WIDTH);
	draw byLine(K, F, byyellow, SOLID_LINE, REGULAR_WIDTH);
	draw byLine(K, G, byblue, SOLID_LINE, REGULAR_WIDTH);
	draw byLine(A, C, bygreen, SOLID_LINE, REGULAR_WIDTH);
	draw byLine(A, K, bygreen, SOLID_LINE, REGULAR_WIDTH);
	draw byLine(A, B, byred, SOLID_LINE, THIN_WIDTH);
	draw byLine(A, D, byblack, SOLID_LINE, THIN_WIDTH);
	draw byLine(B, G, byblack, SOLID_LINE, THIN_WIDTH);	
	draw byLine(G, C, byyellow, SOLID_LINE, THIN_WIDTH);
	draw byLine(D, C, byblue, SOLID_LINE, THIN_WIDTH);
	draw byLabelsOnPolygon(B, E, A, H, D, F, C, G)(ALL_LABELS, 0);
	draw byLabelPoint(K, 40, 3);
    }
    ~\hfill\drawCurrentPicture\hfill~
}{
    \chapter*{Proposition I.43}
    \renewcommand{\currentpropheader}{Proposition I.43}
    \addcontentsline{toc}{chapter}{Proposition 43}
    \textit{In any parallelogram the complements of the parallelograms about the diameter are equal to one another.}
    \vskip \baselineskip

    % Check line numbering toggle
\iftoggle{lineNumbering}{\linenumbers}{}

\iftoggle{heath}{
    Let \drawPolygonAlt[ABCD]{AEK,AHK,HDFK,EBGK,KGC,KFC} be a parallelogram, and \drawUnitLineAlt{AC} its diameter; and about \drawUnitLineAlt{AC} let \drawPolygonAlt[EH]{AEK,AHK}, \drawPolygonAlt[FG]{KGC,KFC} be parallelograms, and \drawPolygonAlt[BK]{EBGK}, \drawPolygonAlt[KD]{HDFK} the so-called complements; \par
    I say that the complement \drawPolygonAlt[BK]{EBGK} is equal to the complement \drawPolygonAlt[KD]{HDFK}. 
    \vskip \baselineskip
    For, since \drawPolygonAlt[ABCD]{AEK,AHK,HDFK,EBGK,KGC,KFC} is a parallelogram, and \drawUnitLineAlt{AC} its diameter, \par
    \centeredjust{the triangle \drawPolygonAlt[ABC]{AEK,EBGK,KGC} is equal to the triangle \drawPolygonAlt[ACD]{AHK,HDFK,KFC} \hspace{0.5em}.}{I. 34}
    Again, since \drawPolygonAlt[EH]{AEK,AHK} is a parallelogram, and \drawUnitLineAlt{AK} is its diameter, \par
    {\centering the triangle \drawPolygonAlt{AEK} is equal to the triangle \drawPolygonAlt{AHK}. \par}
    For the same reason \par
    {\centering the triangle \drawPolygonAlt{KFC} is also equal to \drawPolygonAlt{KGC}. \par}
    Now, since the triangle \drawPolygonAlt{AEK} is equal to the triangle \drawPolygonAlt{AHK},\par
    and \drawPolygonAlt{KFC} to \drawPolygonAlt{KGC},\par
    {\centering the triangle \drawPolygonAlt{AEK} together with \drawPolygonAlt{KGC} is equal to the triangle \drawPolygonAlt{AHK} together with \drawPolygonAlt{KFC}.} \justification{C. N. 2}
    And the whole \drawPolygonAlt[ABC]{AEK,EBGK,KGC} is also equal to the whole \drawPolygonAlt[ADC]{AHK,HDFK,KFC}; \par
    {\centering therefore the complement \drawPolygonAlt[BK]{EBGK} which remains is equal to the complement \drawPolygonAlt[KD]{HDFK} which remains.} \justification{C. N. 3}
    Therefore etc. 
}{}

\iftoggle{raum}{
    Raum_Text
}{}
\nolinenumbers
\hfill Q. E. D.

\vskip \baselineskip


\resetlinenumber
}

\resetlinenumber

\startProposition{
    \markboth{Euclid in Color}{Proposition I.44}
    \defineNewPicture[1/3][3/2]{
	pair F, H, L, B, G, E, A, C, J, K, R, S, M, Q, X, D, Z, d[];
	path q[];
	d1 := (3u, 0);
	d2 := (-3/2u, -3u);
	d3 := (-3u, -9/5u);
	d4 := (-3u, u);
	F := (0, 0);
	K := F shifted d1;
	H := F shifted d2;
	L := H shifted d1;
	B := 2/5[H, K];
	q1 := (B shifted d1) -- (B shifted -d1);
	q2 := (B shifted d2) -- (B shifted -d2);
	G := q1 intersectionpoint (F--H);
	M := q1 intersectionpoint (K--L);
	E := q2 intersectionpoint (F--K);
	A := q2 intersectionpoint (H--L);
	J := F shifted d3; %vertex of C
	S := J shifted (2(xpart(F)-xpart(E)), 0); %vertex of C
	R := (xpart(1/3[J, S]), (ypart(F)-ypart(G)+ypart(J))); %vertex of C
	D := B rotated 30 shifted d4;
	X := 3/5[D, E rotated 30 shifted d4];
	Z := G rotated 30 shifted d4;
	C := ((xpart(J)+xpart(S)+xpart(R))/3, (ypart(J)+ypart(S)+ypart(R))/3); %
	draw byPolygon.C(J,S,R)(byred); %
	draw byPolygon(B,E,F,G)(bygreen);
	draw byPolygon(B,M,L,A)(byblue);
	draw byPolygon(G,B,A,H)(byyellow);
	draw byPolygon(E,K,M,B)(byred);
	byAngleDefine.D(X, D, Z, byyellow, SOLID_SECTOR);
	byAngleDefine(E, B, G, byyellow, SOLID_SECTOR);
	byAngleDefine(A, H, B, byblue, SOLID_SECTOR);
	byAngleDefine(B, H, G, byred, SOLID_SECTOR);
	byAngleDefine(H, F, E, byblue, SOLID_SECTOR);
	byAngleDefine(A, B, M, bygreen, SOLID_SECTOR);
	draw byNamedAngleResized();
	byLineDefine(X, D, byblack, SOLID_LINE, REGULAR_WIDTH);
	byLineDefine(D, Z, byblack, SOLID_LINE, REGULAR_WIDTH);
	draw byNamedLineSeq(0)(XD,DZ);
	byLineDefine(G, B, byred, SOLID_LINE, REGULAR_WIDTH);
	byLineDefine(B, M, byyellow, SOLID_LINE, REGULAR_WIDTH);
	byLineDefine(F, E, byblack, SOLID_LINE, REGULAR_WIDTH);
	byLineDefine(F, G, byyellow, SOLID_LINE, REGULAR_WIDTH);
	byLineDefine(G, H, byblue, SOLID_LINE, REGULAR_WIDTH);
	byLineDefine(A, H, bygreen, SOLID_LINE, REGULAR_WIDTH);
	byLineDefine(A, L, byred, SOLID_LINE, REGULAR_WIDTH);
	byLineDefine(L, M, bygreen, SOLID_LINE, REGULAR_WIDTH);
	byLineDefine(K, M, byblack, SOLID_LINE, REGULAR_WIDTH);
	byLineDefine(E, K, bygreen, SOLID_LINE, REGULAR_WIDTH);
	draw byLine(K, B, byblue, SOLID_LINE, REGULAR_WIDTH);
	draw byLine(B, H, bygreen, SOLID_LINE, REGULAR_WIDTH);
	draw byLine(E, B, byblue, SOLID_LINE, REGULAR_WIDTH);
	draw byLine(A, B, byblack, SOLID_LINE, REGULAR_WIDTH);
	draw byLine(J, S, byblack, SOLID_LINE, REGULAR_WIDTH);
	draw byLine(S, R, byblack, SOLID_LINE, REGULAR_WIDTH);
	draw byLine(R, J, byblack, SOLID_LINE, REGULAR_WIDTH);
	draw byNamedLineSeq(0)(JS, SR, RJ);
	draw byNamedLineSeq(0)(GB,BM,FE,FG,GH,AH,AL,LM,KM,EK);
	draw byLabelsOnPolygon(G, F, E, K, M, L, A, H)(ALL_LABELS, 0);
	draw byLabelsOnPolygon(G, B, E)(OMIT_FIRST_LABEL+OMIT_LAST_LABEL, 0);
	draw byLabelPoint(C, 0, 0);
	draw byLabelPoint(D, 150, 2);
    }
    ~\hfill\drawCurrentPicture\hfill~
}{
    \chapter*{Proposition I.44}
    \renewcommand{\currentpropheader}{Proposition I.44}
    \addcontentsline{toc}{chapter}{Proposition 44}
    \textit{To a given straight line to apply, in a given rectilineal angle, a parallelogram equal to a given triangle.}
    \vskip 0.5\baselineskip

    % Check line numbering toggle
\iftoggle{lineNumbering}{\linenumbers}{}

\iftoggle{heath}{
    Let \drawUnitLineAlt{AB} be the given straight line, \drawPolygonAlt{C} the given triangle and \drawAngleAlt{D} the given rectilineal angle; \par
    thus it is required to apply to the given straight line \drawUnitLineAlt{AB}, in an angle equal to the angle \drawAngleAlt{D}, a parallelogram equal to the given \drawPolygonAlt{C}. 
    \vskip \baselineskip
    Let the parallelogram \drawPolygonAlt{BEFG} be constructed equal to the triangle \drawPolygonAlt{C}, in the angle \drawAngleAlt{EBG} which is equal to \drawAngleAlt{D}; \justification{I. 42}
    let it be placed so that \drawUnitLineAlt{BE} is in a straight line with \drawUnitLineAlt{AB}; \par
    let \drawUnitLineAlt{FG} be drawn through to \drawPointAlt{H},\par
    and let \drawUnitLineAlt{AH} be drawn through \drawPointAlt{A} parallel to either \drawUnitLineAlt{BG} or \drawUnitLineAlt{EF}. \justification{I. 31}
    Let \drawUnitLineAlt{HB} be joined.\par
    Then, since the straight line \drawUnitLineAlt{HF} falls upon the parallels \drawUnitLineAlt{AH},\drawUnitLineAlt{EF}, \par
    \centeredjust{the angles \drawAngleAlt{AHF}, \drawAngleAlt{HFE} are equal to two right angles.}{I. 29}
    Therefore the angles \drawAngleAlt{BHG}, \drawAngleAlt{GFE} are less than two right angles; \par
    and straight lines produced indefinitely from angles less than two right angles meet; \justification{Post. 5}
    {\centering therefore \drawUnitLineAlt{HB},\drawUnitLineAlt{FE}, when produced, will meet. \par}
    Let them be produced and meet at \drawPointAlt{K}; \par
    through the point \drawPointAlt{K} let \drawUnitLineAlt{KL} be drawn parallel to either \drawUnitLineAlt{EA} or \drawUnitLineAlt{FH}, \justification{I. 31}
    and let \drawUnitLineAlt{HA},\drawUnitLineAlt{GB} be produced to the points \drawPointAlt{L}, \drawPointAlt{M}.\par
    Then \drawPolygonAlt[HLKF]{BEFG,BMLA,GBAH,EKMB} is a parallelogram,\par
     \drawUnitLineAlt{HK} is its diameter, \par
    and \drawPolygonAlt[AG]{GBAH}, \drawPolygonAlt[ME]{EKMB} are parallelograms, \par
    and \drawPolygonAlt[LB]{BMLA}, \drawPolygonAlt[BF]{BEFG} the so-called complements, about \drawUnitLineAlt{HK}; \par
    \centeredjust{therefore \drawPolygonAlt[LB]{BMLA} is equal to \drawPolygonAlt[BF]{BEFG}.}{I. 43}
    But \drawPolygonAlt[BF]{BEFG} is equal to the triangle \drawPolygonAlt{C}; \par
    \nolinenumbers
    \clearpage
    \linenumbers
    \centeredjust{therefore \drawPolygonAlt[LB]{BMLA} is also equal to \drawPolygonAlt{C}.}{C. N. 1}
    And, since the angle \drawAngleAlt{GBE} is equal to the angle \drawAngleAlt{ABM}, \justification{I. 15}
    while the angle \drawAngleAlt{GBE} is equal to \drawAngleAlt{D}, \par
    {\centering the angle \drawAngleAlt{ABM} is also equal to the angle \drawAngleAlt{D}. \par}
    Therefore the parallelogram \drawPolygonAlt[LB]{BMLA} equal to the given \drawPolygonAlt{C} has been applied to the given straight line \drawUnitLineAlt{AB}, in the angle \drawAngleAlt{ABM} which is equal to \drawAngleAlt{D}. 
}{}

\iftoggle{raum}{
    Raum_Text
}{}
\nolinenumbers
\hfill Q. E. F.

\vskip \baselineskip


\resetlinenumber
}

\resetlinenumber

\startProposition{
    \markboth{Euclid in Color}{Proposition I.45}
    \defineNewPicture[1/3][11/8]{
	pair A, B, C, D, E, F, G, H, K, L, M, N, P, d[];
	numeric a, h[], b[], s[];
	% h is the height of each triangle, b is the base, and s is the area
	a := 15;
	B := (0, 0);
	C := B shifted (3/2u, 3/2u);
	D := B shifted (0, 2u);
	A := B shifted (-5/2u, 3/2u);
	b1 := arclength(A--B);
	h1 := distanceToLine(D, A--B);
	s1 := (b1 * h1)/2;
	b2 := arclength(B--C);
	h2 := distanceToLine(D, B--C);
	s2 := (b2 * h2)/2;
	d1 := (0, ypart(A)-u);
	d2 := (0, -b1/2) rotated -a; % height of parallelogram, as vector
	d3 := (h1*(1/cosd(a)), 0); % base of parallelogram 1, as vector
	d4 := (2*(s2/b1)*(1/cosd(a)), 0); % base of parallelogram 2, as vector
	d5 := (-3u/2, u/2);
	F := (4u, 3u/2);
	G := F shifted d3;
	H := G shifted d2;
	K := F shifted d2;
	L := G shifted d4;
	M := L shifted d2;
	N := F shifted d5;
	E := K shifted d5;
	P := H shifted d5;
	N := E + 1/2*(N-E);
	P := E + arclength(N--E)/arclength(P--E)*(P-E);
	draw byPolygon(D,B,C)(byred);
	draw byPolygon(A,B,D)(byyellow);
	byLineDefine(D, B, byblue, SOLID_LINE, REGULAR_WIDTH);
	byLineDefine(D, A, bygreen, SOLID_LINE, REGULAR_WIDTH);
	byLineDefine(A, B, byblack, SOLID_LINE, REGULAR_WIDTH);
	byLineDefine(B, C, byblack, SOLID_LINE, REGULAR_WIDTH);
	byLineDefine(C, D, byblack, SOLID_LINE, REGULAR_WIDTH);
	draw byNamedLineSeq(0)(DB);	
	draw byNamedLineSeq(0)(DA, AB, BC, CD);
	draw byPolygon(F,G,H,K)(byyellow);
	draw byPolygon(G,L,M,H)(byred);
	byAngleDefine(H, K, F, byred, SOLID_SECTOR);
	byAngleDefine(G, H, M, byblue, SOLID_SECTOR);
	byAngleDefine(K, H, G, byblack, SOLID_SECTOR);
	byAngleDefine(H, G, L, byyellow, SOLID_SECTOR);
	byAngleDefine(H, G, F, byblue, SOLID_SECTOR);
	byAngleDefine.E(N, E, P, byred, SOLID_SECTOR);
	setAttribute("angle", "Standalone", "NEP", 1);
	draw byNamedAngleResized();
	byLineDefine(A, B, byblack, SOLID_LINE, REGULAR_WIDTH);
	byLineDefine(B, C, byblack, SOLID_LINE, REGULAR_WIDTH);
	byLineDefine(C, D, byblack, SOLID_LINE, REGULAR_WIDTH);
	byLineDefine(F, K, bygreen, SOLID_LINE, REGULAR_WIDTH);
	byLineDefine(G, H, bygreen, SOLID_LINE, REGULAR_WIDTH);
	byLineDefine(K, H, byblue, SOLID_LINE, REGULAR_WIDTH);
	byLineDefine(H, M, byyellow, SOLID_LINE, REGULAR_WIDTH);
	byLineDefine(F, G, byred, SOLID_LINE, REGULAR_WIDTH);
	byLineDefine(G, L, byblue, SOLID_LINE, REGULAR_WIDTH);
	byLineDefine(M, L, bygreen, SOLID_LINE, REGULAR_WIDTH);
	byLineDefine(E, N, byblack, SOLID_LINE, REGULAR_WIDTH);
	byLineDefine(E, P, byblack, SOLID_LINE, REGULAR_WIDTH);
	draw byNamedLineSeq(0)(EN, EP);	
	draw byNamedLineSeq(0)(GH);
	draw byNamedLineSeq(0)(ML, GL, FG, FK, KH, HM);
	% draw byLine(E, P, byblack, SOLID_LINE, REGULAR_WIDTH);
	% draw byLine(E, N, byblack, SOLID_LINE, REGULAR_WIDTH);
	draw byLabelsOnPolygon(D, C, B, A)(ALL_LABELS, 0);
	draw byLabelsOnPolygon(F, G, L, M, H, K)(ALL_LABELS, 0);
	draw byLabelPoint(E, 45, 2);
    }
    ~\hfill\drawCurrentPicture\hfill~
}{
    \chapter*{Proposition I.45}
    \renewcommand{\currentpropheader}{Proposition I.45}
    \addcontentsline{toc}{chapter}{Proposition 45}
    \textit{To construct, in a given rectilineal angle, a parallelogram equal to a given rectilineal figure.}
    \vskip \baselineskip

    % Check line numbering toggle
\iftoggle{lineNumbering}{\linenumbers}{}

\iftoggle{heath}{
    Let \drawPolygonAlt[ABCD]{ABD, DBC} be the given rectilineal figure and \drawAngleAlt{E} the given rectilineal angle; \par
    thus it is required to construct, in the given angle \drawAngleAlt{E}, a parallelogram equal to the rectilineal figure \drawPolygonAlt[ABCD]{ABD, DBC}. 
    \vskip \baselineskip
    Let \drawUnitLineAlt{DB} be joined, and let the parallelogram \drawPolygonAlt[FH]{FGHK} be constructed equal to the triangle \drawPolygonAlt{ABD}, in the angle \drawAngleAlt{HKF} which is equal to \drawAngleAlt{E}; \justification{I. 42}
    let the parallelogram \drawPolygonAlt[GM]{GLMH} equal to the triangle \drawPolygonAlt{DBC} be applied to the straight line \drawUnitLineAlt{GH}, in the angle \drawAngleAlt{GHM} which is equal to \drawAngleAlt{E}. \justification{I. 44}
    Then, since the angle \drawAngleAlt{E} is equal to each of the angles \drawAngleAlt{HKF}, \drawAngleAlt{GHM}, \par
    \centeredjust{the angle \drawAngleAlt{HKF} is also equal to the angle \drawAngleAlt{GHM}.}{C. N. 1}
    Let the angle \drawAngleAlt{KHG} be added to each; \par
    {\centering therefore the angles \drawAngleAlt{FKH}, \drawAngleAlt{KHG} are equal to the angles \drawAngleAlt{KHG}, \drawAngleAlt{GHM}. \par}
    But the angles \drawAngleAlt{FKH}, \drawAngleAlt{KHG} are equal to two right angles; \justification{I. 29}
    {\centering therefore the angles \drawAngleAlt{KHG}, \drawAngleAlt{GHM} are also equal to two right angles. \par}
    Thus, with a straight line \drawUnitLineAlt{GH}, and at the point \drawPointAlt{H} on it, two straight lines \drawUnitLineAlt{KH},\drawUnitLineAlt{HM}[0.65cm] not lying on the same side make the adjacent angles equal to two right angles; \par
    \centeredjust{therefore \drawUnitLineAlt{KH} is in a straight line with \drawUnitLineAlt{HM}[0.65cm].}{I. 14}
    And, since the straight line \drawUnitLineAlt{HG} falls upon the parallels \drawUnitLineAlt{KM}[0.65cm],\drawUnitLineAlt{FG}, \par
    \centeredjust{the alternate angles \drawAngleAlt[MHG]{GHM}, \drawAngleAlt{HGF} are equal to one another.}{I. 29}
    Let the angle \drawAngleAlt{HGL} be added to each; \par
    {\centering therefore the angles \drawAngleAlt[MHG]{GHM}, \drawAngleAlt{HGL} are equal to the angles\par}
    \centeredjust{\drawAngleAlt{HGF}, \drawAngleAlt{HGL}.}{C. N. 2}
    But the angles \drawAngleAlt[MHG]{GHM}, \drawAngleAlt{HGL} are equal to two right angles; \justification{I. 29}
    {\centering therefore the angles \drawAngleAlt{HGF}, \drawAngleAlt{HGL} are also equal to two right angles.}\par\justification{C. N. 1}
    \nolinenumbers
    \clearpage
    \linenumbers
    \centeredjust{Therefore \drawUnitLineAlt{FG} is in a straight line with \drawUnitLineAlt{GL}.}{I. 14}
    And, since \drawUnitLineAlt{FK} is equal and parallel to \drawUnitLineAlt{HG}, \justification{I. 34}
    and \drawUnitLineAlt{HG} to \drawUnitLineAlt{ML} also,\par
    \centeredjust{ \drawUnitLineAlt{KF} is also equal and parallel to \drawUnitLineAlt{ML};}{C. N. 1; I. 30}
    and the straight lines \drawUnitLineAlt{KM}[0.65cm],\drawUnitLineAlt{FL} join them [at their extremities]; \par
    \centeredjust{therefore \drawUnitLineAlt{KM}[0.65cm],\drawUnitLineAlt{FL} are also equal and parallel.}{I. 33}
    {\centering Therefore \drawPolygonAlt[KFLM]{FGHK, GLMH} is a parallelogram. \par}
    And, since the triangle \drawPolygonAlt{ABD} is equal to the parallelogram \drawPolygonAlt[FH]{FGHK}, and \drawPolygonAlt{DBC} to \drawPolygonAlt[GM]{GLMH}, the whole rectilineal figure \drawPolygonAlt[ABCD]{ABD, DBC} is equal to the whole parallelogram \drawPolygonAlt[KFLM]{FGHK, GLMH}. \par
    Therefore the parallelogram \drawPolygonAlt[KFLM]{FGHK, GLMH} has been constructed equal to the given rectilineal figure \drawPolygonAlt[ABCD]{ABD, DBC}, in the angle \drawAngleAlt{FKM} which is equal to the given angle \drawAngleAlt{E}. 
}{}

\iftoggle{raum}{
    Raum_Text
}{}
\nolinenumbers
\hfill Q. E. F.

\vskip \baselineskip


\resetlinenumber
}

\resetlinenumber

\startProposition{
    \markboth{Euclid in Color}{Proposition I.46}
    \defineNewPicture[1/3][5/2]{
	pair A, B, C, D, E;
	numeric d;
	d := 7/2u;
	A := (0, 0);
	B := A shifted (d, 0);
	D := A shifted (0, d);
	E := A shifted (d, d);
	C := 3/2[A,D];
	byAngleDefine(B, A, D, byblack, SOLID_SECTOR);
	byAngleDefine(E, B, A, byblue, SOLID_SECTOR);
	byAngleDefine(D, E, B, byred, SOLID_SECTOR);
	byAngleDefine(A, D, E, byyellow, SOLID_SECTOR);
	draw byNamedAngleResized();
	byLineDefine(A, B, byred, SOLID_LINE, REGULAR_WIDTH);
	byLineDefine(B, E, byyellow, SOLID_LINE, REGULAR_WIDTH);
	byLineDefine(E, D, byblack, SOLID_LINE, REGULAR_WIDTH);
	byLineDefine(C, A, byblue, SOLID_LINE, REGULAR_WIDTH);
	byLineDefine(D, A, byblue, SOLID_LINE, REGULAR_WIDTH);
	draw byNamedLineSeq(0)(AB,BE,ED,CA);
	draw byLabelsOnPolygon(B, A, D, C, E)(ALL_LABELS, 0);
    }
    ~\hfill\drawCurrentPicture\hfill~
}{
    \chapter*{Proposition I.46}
    \renewcommand{\currentpropheader}{Proposition I.46}
    \addcontentsline{toc}{chapter}{Proposition 46}
    \textit{On a given straight line to describe a square.}
    \vskip \baselineskip

    % Check line numbering toggle
\iftoggle{lineNumbering}{\linenumbers}{}

\iftoggle{heath}{
    Let  \drawUnitLineAlt{AB} be the given straight line;\par
    thus it is required to describe a square on the straight line  \drawUnitLineAlt{AB}.
    \vskip\baselineskip
    Let \drawUnitLineAlt{AC} be drawn at right angles to the straight line \drawUnitLineAlt{AB} from the point \drawPointAlt{A} on it, \par\justification{I. 11}
    and let \drawUnitLineAlt{AD} be made equal to \drawUnitLineAlt{AB}; \par
    through the point \drawPointAlt{D} let \drawUnitLineAlt{DE} be drawn parallel to \drawUnitLineAlt{AB}, \par
    and through the point \drawPointAlt{B} let \drawUnitLineAlt{BE} be drawn parallel to \drawUnitLineAlt{AD}. \justification{I. 31}
    Therefore \drawLineAlt{A,D,E,B} is a parallelogram; \par
    \centeredjust{therefore \drawUnitLineAlt{AB} is equal to \drawUnitLineAlt{DE}, and \drawUnitLineAlt{AD} to \drawUnitLineAlt{BE}.}{I. 34}
    But \drawUnitLineAlt{AB} is equal to \drawUnitLineAlt{AD}; \par
    {\centering therefore the four straight lines \drawUnitLineAlt{BA},\drawUnitLineAlt{AD},\drawUnitLineAlt{DE},\drawUnitLineAlt{EB} are equal to one another; \par}
    {\centering therefore the parallelogram \drawLineAlt{A,D,E,B} is equilateral. \par}
    I say next that it is also right-angled. \par
    For, since the straight line \drawUnitLineAlt{AD} falls upon the parallels \drawUnitLineAlt{AB},\drawUnitLineAlt{DE}, \par
    \centeredjust{the angles \drawAngleAlt{BAD}, \drawAngleAlt{ADE} are equal to two right angles.}{I. 29}
    But the angle \drawAngleAlt{BAD} is right; \par
    {\centering therefore the angle \drawAngleAlt{ADE} is also right. \par}
    And in parallelogrammic areas the opposite sides and angles are equal to one another; \justification{I. 34}
    {\centering therefore each of the opposite angles \drawAngleAlt{ABE}, \drawAngleAlt{BED} is also right. \par}
    Therefore \drawLineAlt{A,D,E,B} is right-angled. \par
    And it was also proved equilateral. \par
    Therefore it is a square; and it is described on the straight line \drawUnitLineAlt{AB}. 
 }{}

\iftoggle{raum}{
    Raum_Text
}{}
\nolinenumbers
\hfill Q. E. F.

\vskip \baselineskip


\resetlinenumber
}

\resetlinenumber

\startProposition{
    \markboth{Euclid in Color}{Proposition I.47}
    \defineNewPicture[1/3][9/4]{
	pair A, B, C, G, F, K, H, D, E, J, L, N, M, d[];
	A := (0, 0);
	B := A shifted (-9/10u, -9/7u);
	C = whatever[A, A shifted ((A-B) rotated 90)] = whatever[B, B shifted dir(0)];
	d1 := (B-A) rotated -90;
	G := A shifted d1;
	F := B shifted d1;
	d2 := (A-C) rotated -90;
	K := C shifted d2;
	H := A shifted d2;
	d3 := (C-B) rotated -90;
	D := B shifted d3;
	E := C shifted d3;
	J = whatever[A, A shifted dir(90)];
	J = whatever[B, C];
	L = whatever[A, A shifted dir(90)];
	L = whatever[D, E];
	N = whatever[B, K];
	N = whatever[A, C];
	M = whatever[A, E];
	M = whatever[B, C];
	draw byPolygon.GB(A,B,F,G)(byblack);
	draw byPolygon.HC(A,C,K,H)(byred);
	%draw byPolygon(C,N,K)(byred);
	draw byPolygon.CL(J,C,E,L)(byyellow);
	draw byPolygon.BL(B,J,L,D)(byblue);
	byAngleDefine(F, B, A, byblue, SOLID_SECTOR);
	byAngleDefine(D, B, C, byyellow, SOLID_SECTOR);
	byAngleDefine(A, B, C, bygreen, SOLID_SECTOR);
	byAngleDefine(B, A, C, byred, SOLID_SECTOR);
	byAngleDefine(B, A, G, byblack, SOLID_SECTOR);
	draw byNamedAngleResized();
	byLineDefine(A, B, byyellow, SOLID_LINE, REGULAR_WIDTH);
	byLineDefine(A, G, byred, SOLID_LINE, REGULAR_WIDTH);
	byLineDefine(G, F, byblue, SOLID_LINE, THIN_WIDTH);
	byLineDefine(F, B, bygreen, SOLID_LINE, REGULAR_WIDTH);
	draw byNamedLineSeq(0)(AB,AG,GF,FB);
	byLineDefine(C, A, byblue, SOLID_LINE, REGULAR_WIDTH);
	byLineDefine(A, H, byblack, SOLID_LINE, REGULAR_WIDTH);
	byLineDefine(H, K, bygreen, SOLID_LINE, THIN_WIDTH);
	byLineDefine(C, K, byyellow, SOLID_LINE, REGULAR_WIDTH);
	draw byNamedLineSeq(0)(CA,AH,HK,CK);
	byLineDefine(D, E, byblack, SOLID_LINE, THIN_WIDTH);
	byLineDefine(B, D, byblack, SOLID_LINE, REGULAR_WIDTH);
	byLineDefine(B, C, byred, SOLID_LINE, REGULAR_WIDTH);
	byLineDefine(C, E, bygreen, SOLID_LINE, REGULAR_WIDTH);
	draw byNamedLineSeq(0)(DE,BD,BC,CE);
	draw byLine(A, L, byblack, SOLID_LINE, THIN_WIDTH);
	draw byLine(B, K, byblack, SOLID_LINE, THIN_WIDTH);
	draw byLine(A, E, byblack, SOLID_LINE, THIN_WIDTH);
	draw byLine(A, D, byblack, SOLID_LINE, THIN_WIDTH);
	draw byLine(F, C, byblue, SOLID_LINE, THIN_WIDTH);
	byLineDefine.CAb(C, A, byblack, SOLID_LINE, REGULAR_WIDTH);
	byLineStylize (M, M, 1, 0, -1) (CAb);
	byLineDefine.AMb(A, M, byblack, SOLID_LINE, REGULAR_WIDTH);
	byLineStylize (C, C, 0, 1, -1) (AMb);
	byLineDefine.BCb(B, C, byblack, SOLID_LINE, REGULAR_WIDTH);
	byLineStylize (N, N, 0, 1, -1) (BCb);
	byLineDefine.BNb(N, B, byblack, SOLID_LINE, REGULAR_WIDTH);
	byLineStylize (C, C, 1, 0, -1) (BNb);
	draw byLabelsOnPolygon(B, F, G, A, H, K, C, E, L, D)(ALL_LABELS, -1);
    }
    ~\hfill\drawCurrentPicture\hfill~
}{
    \chapter*{Proposition I.47}
    \renewcommand{\currentpropheader}{Proposition I.47}
    \addcontentsline{toc}{chapter}{Proposition 47}
    \textit{In right-angled triangles the square on the side subtending the right angle is equal to the squares on the sides containing the right angle.}
    \vskip \baselineskip

    % Check line numbering toggle
\iftoggle{lineNumbering}{\linenumbers}{}

\iftoggle{heath}{
    Let \drawLineAlt{A,B,C} be a right-angled triangle having the angle \drawAngleAlt{BAC} right; \par
    I say that the square on \drawUnitLineAlt{BC} is equal to the squares on \drawUnitLineAlt{BA},\drawUnitLineAlt{AC}. 
    \vskip \baselineskip
    For let there be described on \drawUnitLineAlt{BC} the square \drawPolygonAlt[BDEC]{BL,CL}, and on \drawUnitLineAlt{BA},\drawUnitLineAlt{AC} the squares \drawPolygonAlt{GB}, \drawPolygonAlt{HC}; \justification{I. 46}
    through \drawPointAlt{A} let \drawUnitLineAlt{AL} be drawn parallel to either \drawUnitLineAlt{BD} or \drawUnitLineAlt{CE}, \par
    and let \drawUnitLineAlt{AD},\drawUnitLineAlt{FC} be joined. \par
    Then, since each of the angles \drawAngleAlt{BAC}, \drawAngleAlt{BAG} is right, \par
    it follows that with a straight line \drawUnitLineAlt{BA}, and at the point \drawPointAlt{A} on it, the two straight lines \drawUnitLineAlt{AC},\drawUnitLineAlt{AG} not lying on the same side make the adjacent angles equal to two right angles; \par
    \centeredjust{therefore \drawUnitLineAlt{CA} is in a straight line with \drawUnitLineAlt{AG}.}{I. 14}
    For the same reason \drawUnitLineAlt{BA} is also in a straight line with \drawUnitLineAlt{AH}. \par
    And, since the angle \drawAngleAlt{DBC} is equal to the angle \drawAngleAlt{FBA}: for each is right: let the angle \drawAngleAlt{ABC} be added to each; \par
    \centeredjust{therefore the whole angle \drawAngleAlt{DBA} is equal to the whole angle \drawAngleAlt{FBC}.}{C. N. 2}
    And, since \drawUnitLineAlt{DB} is equal to \drawUnitLineAlt{BC}, and \drawUnitLineAlt{FB} to \drawUnitLineAlt{BA}, \par
    {\centering the two sides \drawUnitLineAlt{AB},\drawUnitLineAlt{BD} are equal to the two sides \drawUnitLineAlt{FB},\drawUnitLineAlt{BC} respectively, \par}
    and the angle \drawAngleAlt{ABD} is equal to the angle \drawAngleAlt{FBC}; \par
    {\centering therefore the base \drawUnitLineAlt{AD} is equal to the base \drawUnitLineAlt{FC}, \par}
    \nolinenumbers
    \clearpage
    \linenumbers
    \centeredjust{and the triangle \drawLineAlt{A,B,D} is equal to the triangle \drawLineAlt{F,B,C}.}{I. 4}
    Now the parallelogram \drawPolygonAlt{BL} is double of the triangle \drawLineAlt{A,B,D}, for they have the same base \drawUnitLineAlt{BD} and are in the same parallels \drawUnitLineAlt{BD},\drawUnitLineAlt{AL}. \justification{I. 41}
    And the square \drawPolygonAlt{GB} is double of the triangle \drawLineAlt{F,B,C}, for they again have the same base \drawUnitLineAlt{FB} and are in the same parallels \drawUnitLineAlt{FB},\drawUnitLineAlt{GC}. \justification{I. 41}
    $\langle$But the doubles of equals are equal to one another.$\rangle$ \par
    {\centering Therefore the parallelogram \drawPolygonAlt{BL} is also equal to the square \drawPolygonAlt{GB}. \par}
    Similarly, if \drawUnitLineAlt{AE},\drawUnitLineAlt{BK} be joined, \par
    the parallelogram \drawPolygonAlt{CL} can also be proved equal to the square \drawPolygonAlt{HC};\par
    {\centering therefore the whole square \drawPolygonAlt[BDEC]{BL,CL} is equal to the two squares \drawPolygonAlt{GB}, \drawPolygonAlt{HC}.} \justification{C. N. 2}
    And the square \drawPolygonAlt[BDEC]{BL,CL} is described on \drawUnitLineAlt{BC}, and the squares \drawPolygonAlt{GB}, \drawPolygonAlt{HC} on \drawUnitLineAlt{BA},\drawUnitLineAlt{AC}. \par
    {\centering Therefore the square on the side \drawUnitLineAlt{BC} is equal to the squares on the sides \\ \drawUnitLineAlt{BA},\drawUnitLineAlt{AC}. \par}
    Therefore etc. 
}{}

\iftoggle{raum}{
    Raum_Text
}{}
\nolinenumbers
\hfill Q. E. D. 

\vskip \baselineskip


\resetlinenumber
}

\resetlinenumber

\startProposition{
    \markboth{Euclid in Color}{Proposition I.48}
    \defineNewPicture[1/4][3]{
	pair A, B, C, D;
	numeric d;
	d := 7/4u;
	A := (0, 0);
	C := A shifted (0, 4u);
	B := A shifted (d, 0);
	D := A shifted (-d, 0);
	byAngleDefine(C, A, B, byred, SOLID_SECTOR);
	byAngleDefine(D, A, C, byyellow, SOLID_SECTOR);
	draw byNamedAngleResized();
	draw byLine(A, C, byblue, SOLID_LINE, REGULAR_WIDTH);
	byLineDefine(A, B, byyellow, SOLID_LINE, REGULAR_WIDTH);
	byLineDefine(A, D, byred, SOLID_LINE, REGULAR_WIDTH);
	byLineDefine(C, B, byred, SOLID_LINE, REGULAR_WIDTH);
	byLineDefine(C, D, bygreen, SOLID_LINE, REGULAR_WIDTH);
	draw byNamedLineSeq(0)(AB,AD,CD,CB);
	draw byLabelsOnPolygon(D, C, B, A)(ALL_LABELS, 0);
    }
    ~\hfill\drawCurrentPicture\hfill~
}{
    \chapter*{Proposition I.48}
    \renewcommand{\currentpropheader}{Proposition I.48}
    \addcontentsline{toc}{chapter}{Proposition 48}
    \textit{If in a triangle the square on one of the sides be equal to the squares on the remaining two sides of the triangle, the angle contained by the remaining two sides of the triangle is right.}
    \vskip \baselineskip

    % Check line numbering toggle
\iftoggle{lineNumbering}{\linenumbers}{}

\iftoggle{heath}{
    For in the triangle \drawLineAlt{A,B,C} let the square on one side \drawUnitLineAlt{BC} be equal to the squares on the sides \drawUnitLineAlt{BA},\drawUnitLineAlt{AC}; \par
    I say that the angle \drawAngleAlt{BAC} is right. 
    \vskip \baselineskip
    For let \drawUnitLineAlt{AD} be drawn from the point \drawPointAlt{A} at right angles to the straight line \drawUnitLineAlt{AC}, \par
    let \drawUnitLineAlt{AD} be made equal to \drawUnitLineAlt{BA}, \par
    and let \drawUnitLineAlt{DC} be joined. \par
    Since \drawUnitLineAlt{DA} is equal to \drawUnitLineAlt{AB}, \par
    {\centering the square on \drawUnitLineAlt{DA} is also equal to the square on \drawUnitLineAlt{AB}. \par}
    Let the square on \drawUnitLineAlt{AC} be added to each; \par
    {\centering therefore the squares on \drawUnitLineAlt{DA},\drawUnitLineAlt{AC} are equal to the squares on \drawUnitLineAlt{BA},\drawUnitLineAlt{AC}. \par}
    But the square on \drawUnitLineAlt{DC} is equal to the squares on \drawUnitLineAlt{DA},\drawUnitLineAlt{AC}, for the angle \drawAngleAlt{DAC} is right; \par
    \justification{I. 47}
    and the square on \drawUnitLineAlt{BC} is equal to the squares on \drawUnitLineAlt{BA},\drawUnitLineAlt{AC}, for this is the hypothesis; \par
    {\centering therefore the square on \drawUnitLineAlt{DC} is equal to the square on \drawUnitLineAlt{BC}, \par}
    {\centering so that the side \drawUnitLineAlt{DC} is also equal to \drawUnitLineAlt{BC}. \par}
    And, since \drawUnitLineAlt{DA} is equal to \drawUnitLineAlt{AB}, and \drawUnitLineAlt{AC} is common, \par
    {\centering the two sides \drawUnitLineAlt{DA},\drawUnitLineAlt{AC} are equal to the two sides \drawUnitLineAlt{BA},\drawUnitLineAlt{AC}; \par}
    and the base \drawUnitLineAlt{DC} is equal to the base \drawUnitLineAlt{BC};\par
    \centeredjust{therefore the angle \drawAngleAlt{DAC} is equal to the angle \drawAngleAlt{BAC}.}{I. 8}
    But the angle \drawAngleAlt{DAC} is right;\par
    {\centering therefore the angle \drawAngleAlt{BAC} is also right. \par}
    Therefore etc. 
}{}

\iftoggle{raum}{
    Raum_Text
}{}
\nolinenumbers
\hfill Q. E. D.

\vskip \baselineskip
}

\end{document}
